% 本文件由 LaTeX 排版 Agent 根据有效 translation.json 自动生成。
% author=John Chrysostom; work_id=2304; title=得撒洛尼前书讲道集

\chapter{得撒洛尼前书讲道集}

% structure: p0001 source=h1 target=omitted reason=page-title-already-rendered-by-parent

讲道一\newline{}讲道二\newline{}讲道三\newline{}讲道四\newline{}讲道五\newline{}讲道六\newline{}讲道七\newline{}讲道八\newline{}讲道九\newline{}讲道十\newline{}讲道十一\par

\SourceRecord{Translated by John A. Broadus.；Nicene and Post-Nicene Fathers, First Series；Vol. 13.；Edited by Philip Schaff.；Buffalo, NY: Christian Literature Publishing Co.,；1889.；<http://www.newadvent.org/fathers/2304.htm>.}

% structure: p0001 source=h1 target=section reason=page-title

\section{第一篇讲道 论得撒洛尼前书}

\BibleRef{得前 1:1-3}\par

\BibleQuote{\Emph{保禄、息耳瓦诺和弟茂德，致得撒洛尼人的教会，就是属于天主父和主耶稣基督的：愿恩宠与平安归于你们。} \Emph{我们常为你们众人感谢天主，在祈祷中常记念你们；在天主和我们的父面前，不断地记念你们因信德所做的工作，因爱德所受的劳苦，因望德所怀的忍耐，都是赖我们的主耶稣基督。}}\par

\Emph{因此}，为何呢？当写信给厄弗所人时，弟茂德与他同在，他却没有将他与自己一同列入（问候语中），尽管弟茂德为他们所认识并受钦佩？因为他说：\BibleQuote{你知道他的明证，他怎样像儿子侍奉父亲一样，与我一同在福音中侍奉}\BibleRef{斐 2:22}；又说：\BibleQuote{我没有别人像他那样真心关心你们的景况}\BibleRef{斐 2:20}；但在这里他却将弟茂德与自己联合呢？在我看来，是因为他即将立刻派遣弟茂德，而他将赶在这封信之前到达，所以写信就是多余的了。因为他说：\BibleQuote{所以我希望立刻打发他}\BibleRef{斐 2:23}。但这里并非如此；弟茂德刚回到他那里，所以他自然参与到信函中。因为他说：\BibleQuote{但弟茂德从你们那里来到我们这里}\BibleRef{得前 3:6}。但为什么他把息耳瓦诺放在弟茂德前面，尽管他见证弟茂德无数的优点，且将他置于众人之上？或许弟茂德出于极大的谦卑，希望并请求他如此做；因为他看见他的老师如此谦卑，竟将自己的门徒与自己同列，他就更渴望如此，并热切追求。因为他说：\par

\BibleQuote{保禄、息耳瓦诺和弟茂德，致得撒洛尼人的教会。} 在这里他并未给自己任何称号——既非\InnerQuote{宗徒}，也非\InnerQuote{仆人}；我想，是因为这些人初受教导，尚未有任何关于他的经验，所以他未运用称号；而且那时正是他向他们宣讲的开始。\par

\BibleQuote{致得撒洛尼人的教会}，他说。这话说得好。因为很可能人数稀少，且尚未形成团体；为此他以教会的名称安慰他们。因为在那些已度过许多时日、教会会众已众多的地方，他并不用这个称呼。但教会这一名称多半是众多之人的名称，也是已结成一体的组织的名称，为此他如此称呼他们。\par

\BibleQuote{在天主父内}，他说，\BibleQuote{和主耶稣基督内}。\BibleQuote{致得撒洛尼人的教会}，他说，\BibleQuote{即在天主内的}。看，再次使用了\Quote{在}这个词汇，既用于父也用于子。因为有许多集会，既有犹太人的也有希腊人的；但他却说：\BibleQuote{致那在天主内的（教会）}。这是一个伟大的尊严，没有其他能与它相比，即它\Quote{在天主内}。因此，愿天主赐予这教会配得上这样的称呼！但我恐怕它离这称呼甚远。因为若有人是罪的奴隶，就不能被称为\Quote{在天主内}。若有人不按天主行事，就不能被称为\Quote{在天主内}。\par

\BibleQuote{愿恩宠与平安归于你们。} 你看出这封书信的开端竟是赞美吗？\BibleQuote{我们常为你们众人感谢天主，在祈祷中常记念你们。} 因为为他们感谢天主，是见证他们极大的进步，不仅他们自己受赞美，而且天主也因他们被感谢，好像这一切都是他自己所为。他也教导他们要节制，几乎是说，这一切都是天主的德能。他为他们感谢，是出于他们的善行；而他在祈祷中记念他们，则是出于他对他们的爱。然后，如他常做的，他说他不仅在他的祈祷中记念他们，而且超越他的祈祷。\Quote{不断地记念}，他说，\BibleQuote{你们因信德所做的工作，因爱德所受的劳苦，因望德所怀的忍耐，都是赖我们的主耶稣基督，在天主和我们的父面前。} 什么是\Quote{不断地记念}？或者是在天主和父面前记念，或者是记念你们在天主和父面前的因爱德所受的劳苦，或者简单地说：\Quote{不断地记念你们}。然后，免得你以为这\Quote{不断地记念你们}是随便说的，他加上了\Quote{在天主和我们的父面前}。因为没有人称赞他们的行动，没有人给他们任何报酬，他说这话：\Quote{你们在天主面前劳苦}。什么是\Quote{信德的工作}？就是没有什么能使你们的坚定动摇。因为这是信德的工作。你若相信，就忍受一切；你若不忍受，你就不相信。因为所应许的，岂不正是这样，使相信的人愿意忍受甚至千万种死亡吗？天国、不朽和永生都已摆在他面前。因此，相信的人将忍受一切。信德便通过他的工作显明出来。正可这样说：你们不仅相信了，而且通过你们的工作彰显了它，通过你们的坚定，通过你们的熱忱。\par

\BibleQuote{以及你们因爱德所受的劳苦。} 为什么？爱有什么劳苦？仅仅是爱并不算是劳苦。但真诚地爱却是极大的劳苦。因为请告诉我，当有千般事搅动我们远离爱，而我们却抗拒它们全部，这岂非劳苦？这些人难道没有忍受那些事，为不背弃他们的爱？那些攻打宣讲的人不是去了保禄的房东那里，没有找到他，就把雅松拉到城市首领面前吗？\BibleRef{宗 17:5} 请告诉我，当种子尚未扎根时，忍受如此大的风暴、如此多的考验，这难道是小劳苦吗？他们还向他索取保证金。雅松给出了保证金后，就送走了保禄。请告诉我，这是小事吗？雅松难道没有为他置自身于危险中吗？他称此为因爱德所受的劳苦，因为他们这样与他联结。\par

请注意：他首先提及他们的善行，然后才提及自己的，这样他既能避免显得自夸，也不会显得是在预先爱他们。他说：\Quote{还有坚忍。}因为那迫害并非仅限于一时，而是持续不断的；他们不但与身为导师的\Person{保禄}争战，也与他的门徒们争战。如果他们对那些行奇迹者、那些可敬的人尚且如此，那么你们想想，他们对那些住在他们中间、突然背离他们的同胞又会是怎样的感受？因此，他也为他们作证说：\BibleQuote{因为你们成了在\Place{犹太}的、属于天主的各教会的效法者。}\BibleRef{得前 2:14}\par

他说：\BibleQuote{以及在我们的主耶稣基督内的望德，在我们的天主和父面前。}\BibleRef{得前 1:3}因为这一切都源于信德和望德，所以发生在他们身上的事不仅显示了他们的刚毅，更表明他们完全确信那为他们预备的赏报。正因如此，天主立即容许迫害发生，好叫无人能说，那宣讲是靠轻率或谄媚建立的；也为了彰显他们热火般的热情，并表明那说服信徒灵魂的不是人的劝说，而是天主的德能——他们甚至准备面对万死。倘若那宣讲没有立即深深扎根并屹立不摇，这便不会发生。\par

第4、5节。\BibleQuote{弟兄们，天主所爱的，我们知道你们是如何被拣选的：因为我们的福音传到你们那里，不仅在于言语，也在于德能、在于圣神，并在充足的把握中；正如你们也知道，我们在你们中间为你们的缘故是怎样的人。}\BibleRef{得前 1:4-5}\par

知道什么？是\BibleQuote{我们在你们中间是怎样的人}\BibleRef{得前 1:5}？他在这里也隐约触及自己的善行，但却是含蓄的。因为他先想扩大对他们的称赞，他说的大意是这样的：我知道你们是伟大而高贵的人，你们是蒙拣选的。因此，我们也为你们的缘故忍受一切。因为这\BibleQuote{我们在你们中间是怎样的人}\BibleRef{得前 1:5} 这句话，表明我们以极大的热忱和极大的奋力，准备为你们的缘故舍弃自己的生命；为此，谢意不归给我们，而归于你们，因为你们是蒙拣选的。因此他在别处也说：\Quote{我为蒙拣选的人忍受这一切。}\BibleRef{弟后 2:10} 因为岂会有人不为天主所爱的人忍受一切呢？在述说自己的部分之后，他几乎要说：既然你们既是蒙爱的又是蒙拣选的，我们遭受一切便合情合理。因为对他的称赞不仅坚定了他们，也提醒他们自己也展现了与热忱相应的刚毅。他说：\par

第6节。\BibleQuote{你们就成了我们和主的效法者，在极大的苦难中，怀着圣神的喜乐接受了圣言。}\BibleRef{得前 1:6}\par

奇妙！这是何等的赞美！门徒们突然成了教师！他们不仅听了圣言，还迅速达到了与\Person{保禄}相同的高度。但这算不了什么；请看他是如何高举他们的，他说：\BibleQuote{你们成了主的效法者。}\BibleRef{得前 1:6}如何？\BibleQuote{在极大的苦难中，怀着圣神的喜乐接受了圣言。}\BibleRef{得前 1:6}不仅是苦难，而且是极大的苦难。这一点我们可以从\WorkTitle{宗徒大事录}中得知，他们如何兴起对他们的迫害。\BibleRef{宗 17:5} 他们扰乱了全城的统治者，煽动全城反对他们。还不仅如此，你们确实受了苦，并且相信了，但不是在悲伤中，而是在喜乐中。这也是宗徒们所做的：经上说：\BibleQuote{他们喜乐，因为他们算配为这名受辱。}\BibleRef{宗 5:41} 因为这才是可赞叹的。尽管那能忍受苦难本身也不算小事。但这现在乃是超越人性、似乎拥有不可伤害之身体的人的表现。\par

他们如何成了主的效法者？因为主也忍受了许多苦难，却喜乐了。因为他是自愿来此。为了我们，他空虚了自己。他将要受人唾骂、被鞭打、被钉十字架，他却如此喜乐地承受这些，以致他对父说：\BibleQuote{光荣我吧。}\BibleRef{若 17:1}\par

他说：\Quote{怀着圣神的喜乐。}免得有人问：你怎么说\Quote{苦难}？怎么又说\Quote{喜乐}？两者怎能调和呢？他补充说：\Quote{怀着圣神的喜乐。}苦难在于身体的事，而喜乐在于灵性的事。如何呢？发生在他们身上的事是痛苦的，但由此而生的事却不然，因为圣神不容许那样。因此，一个人既可以因自己的罪而受苦却不喜乐，也可以为基督的缘故受鞭打而愉悦。因为圣神的喜乐就是如此：它从看似痛苦的事中引出欢愉。他说：他们逼迫了你们，迫害了你们，但即使在那种境况中，圣神也没有离弃你们。正如那三个少年在火中得到露水清凉，你们也在苦难中得到了清凉。但正如那里，火的本质不会洒下露水，而是那\BibleQuote{呼啸的风}\BibleRef{达 3:50}；同样，这里，苦难的本质不会产生喜乐，而是为基督受苦的缘故，以及那滋润他们、在诱惑的熔炉中使他们安然的圣神。他说：不仅是喜乐，而且是\Quote{极大的喜乐}。因为这是出于圣神的事。\par

第7节。\BibleQuote{以致你们成了在\Place{马其顿}和\Place{阿哈雅}所有信徒的榜样。}\BibleRef{得前 1:7}\par

然而，他后来才到他们那里去。但他说：你们如此闪耀，以致成了那些在你们之前领受圣言者的教师。这正符合宗徒的作风。因为他没有说：你们成了信德方面的榜样；而是说：你们成了那些已经相信的人的榜样——你们教导他们应当如何相信天主，因为你们从一开始就投入了战斗。\par

他说：\Quote{以及在\Place{阿哈雅}}，即\Place{希腊}。\par

你们看见热忱是多么伟大的事吗？它不需要时间，不需要迟延，不需要拖延，只要敢于献出自己，一切就完成了。这样，虽然后来才加入宣讲，他们却成了比他们先来的人的导师。\par

训诲。因此，但愿没有人失望，即使他荒废了许多时间，一事无成。他也可以在短时间内做这么多事，甚至超过他以前所有的时间。因为连先前不信的人，在开头都如此光辉，何况那些已经相信的人呢！但愿也没有人因此松懈，因为他看到在短时间内可以挽回一切。因为未来是不确定的，主的日子像贼一样，在我们睡觉时突然来袭。但如果我们不睡觉，它就不会像贼那样袭击我们，也不会让我们措手不及。如果我们警醒和清醒，它不会像贼那样袭击我们，而是像君王的使者，召唤我们赴那为我们预备的美善。但如果我们睡觉，它就会像贼一样来到。因此，但愿没有人睡觉，也不要在德行上怠惰，因为那就是睡觉。你难道不知道吗？当我们睡觉时，我们的财产不安全，多么容易被算计？但当我们醒着时，就不需要那么多看守。当我们睡觉时，即使有很多看守，我们也常常灭亡。有门、闩、卫兵、外卫，贼却仍然到了我们身上。\par

我为什么这样说呢？因为如果我们醒着，就不需要别人的帮助；但如果我们睡觉，别人的帮助对我们毫无益处，即使有这帮助我们也要灭亡。享受圣人的祈祷是好事，但也要我们自己警醒。你说，如果我本人警醒，还需要别人的祈祷吗？实在，不要让自己陷入需要它的境地；我不希望你这样；但如果我们正确思考，我们总是需要它。保禄没有说：我需要祈祷做什么？那些祈祷的人并不配他，或者说不及他；而你却说：我需要祈祷做什么？伯多禄没有说：我需要祈祷做什么？因为经上说：\BibleQuote{教会为他向天主恳切祈祷。}\BibleRef{宗 12:5}而你却说：我需要祈祷做什么？正因如此，你需要它，因为你认为你不需要。即使你变得像保禄一样，你仍然需要祈祷。不要自高，免得你被贬抑。\par

但是，正如我所说的，如果我们自己也努力，为我们的祈祷也有效。请听保禄说：\BibleQuote{因为我知道，这事藉着你们的祈祷和耶稣基督之圣神的供应，终必使我得救。}\BibleRef{斐 1:19}又说：\BibleQuote{使那藉着许多人为我们所得的恩赐，因着许多人的感谢而越发显多。}\BibleRef{格后 1:11}而你却说：我需要祈祷做什么？但如果我们懒惰，没有人能帮我们。耶肋米亚能帮犹太人什么？他不是三次接近天主，第三次听到：\BibleQuote{你不要为这百姓祈祷，不要为他们呼号祈祷，因为我不会听你。}\BibleRef{耶 7:16}撒慕尔能帮撒乌耳什么？他不是哀悼他直到末日，而且不只是为他祈祷？他能帮以色列人什么？他不是说：\BibleQuote{我决不让停止为你们祈祷成为我的罪过。}\BibleRef{撒上 12:23}他们不是全都灭亡了吗？那么，你说，祈祷就毫无益处吗？它们非常有帮助：但也要我们自己有所作为。因为祈祷确实合作并协助，但一个人是与正在行动的人合作，是帮助正在工作的人。但如果你偷懒，你就不会得到多大的益处。\par

因为如果祈祷能让我们不用做任何事就进入天国，为什么所有希腊人不都成为基督徒？我们不是为普世祈祷吗？保禄不也这样做吗？我们不是恳求所有人都能悔改吗？为什么恶人不自己出一点力就变好呢？因此，当我们自己也尽本分时，祈祷就大有助益。\par

你愿意知道祈祷有多大益处吗？我想，想一想科尔乃略和塔彼达吧。\BibleRef{宗 10:3}也听雅各伯对拉班说：\BibleQuote{若不是我父亲敬畏的天主与我同在，你如今必定打发我空手而去。}\BibleRef{创 31:42}再听天主说：\BibleQuote{为我自己和我仆人达味的缘故，我必保护这城。}\BibleRef{列下 9:34}这是在什么时候？是在义人希则克雅的时代。既然祈祷甚至对最恶的人有效，为什么当拿步高到来时，天主没有说这话，却把城交出去了？因为邪恶更强大。同样，撒慕尔本人也为以色列人祈祷，并得了胜。但什么时候？当他们也使天主喜悦的时候，他们就打败了敌人。你说，如果我本人使天主喜悦，还需要别人的祈祷吗？人哪，绝不要说这话。需要，而且需要许多祈祷。因为请听天主论到约伯的朋友们说：\BibleQuote{他要为你们祈祷，你们的罪必得赦免。}\BibleRef{约 42:8}他固然犯了罪，但不是大罪。而这位以祈祷拯救了朋友的义人，在犹太人的时节却不能拯救正在灭亡的犹太人。为使你明白这一点，请听天主藉先知说：\BibleQuote{若诺厄、达尼尔和约伯在那里，他们也不能救自己的子女。}\BibleRef{则 14:14}因为邪恶占了上风。又说：\BibleQuote{即使梅瑟和撒慕尔站在我面前。}\BibleRef{耶 15:1}\par

并且请看，这话是对两位先知说的，因为他们二人都为他们恳求，却没有得到应允。因为厄则克耳说：\BibleQuote{啊，主，祢要灭绝以色列的遗民吗？}\BibleRef{厄 9:8} 然后，主显示祂这样做是公正的，便向他们展示他们的罪过；并且显示祂并非因轻视他才不肯接受他的代祷，祂说：这些事足以说服你，我并非轻视你，而是因他们的许多罪过，我不接受你的祈求。然而，祂又加上：\BibleQuote{即便诺厄、约伯和达尼尔站在。}\EditorialNote{出自厄则克耳14章} 祂更有理由对厄则克耳说这话，因为他曾忍受了那么多的事。他说：祢命令我吃粪便，我便吃了；祢命令我剃头，我便剃了；祢命令我侧卧，我便侧卧；祢命令我从墙洞出去，背负重担，我便出去了；祢夺去了我的妻子，并命令我不要哀悼，我便没有哀悼，反而刚强地承受了。\BibleRef{厄 24:18} 我曾为他们做了成千上万的事；我为他们恳求，祢岂不答应吗？祂说：我这样做并非轻视你，但即使诺厄、约伯和达尼尔在那里，为他们的儿女恳求，我也不会答应。因此，知道这些事，我们既不可轻视圣人的祈祷，也不可将一切全推给他们：免得我们一方面懒惰、漫不经心地生活，另一方面又剥夺自己的巨大益处。而是要既恳求他们为我们祈祷、举手，又要坚持修德；这样我们才能藉着我们的主耶稣基督的恩宠和慈爱，获得应许给爱祂之人的福乐，偕同祂，等等。\par

再说给耶肋米亚，他在天主的命令上所受的苦较少，但在他们的邪恶上所受的苦较多，天主对他说什么？\BibleQuote{你没看见他们做什么吗？}\BibleRef{耶 7:17} 他说：\Quote{是的，他们这样做——但求祢为我的缘故这样做。} 因此，主对他说：\Quote{即使梅瑟和撒慕尔站在。} 他们的第一位立法人，曾多次将他们从危险中拯救出来，他曾说：\BibleQuote{现在，如果祢宽恕他们的罪，就宽恕吧；如果不，也请把我从你的册子上抹去。}\BibleRef{出 32:32} 因此，如果他如今活着，这样说，他也不会成功——撒慕尔也一样，他曾拯救他们，并且自幼就受人钦佩。因为关于前者，我曾说，我与他交谈，如同朋友对朋友，并非通过隐晦的言语。关于后者，我曾说，在他幼年时我就向他启示，并且因他的缘故，我被说服，打开了曾被封闭的预言。因为经上说：\BibleQuote{在那时候，上主的话是稀有的，异象並不常见。}\BibleRef{撒上 3:1} 因此，如果这些人站在我面前，将无济于事。关于诺厄，祂说：\BibleQuote{诺厄是个义人，在他的世代中是个完美的人。}\BibleRef{创 6:9} 关于约伯，祂说他是\BibleQuote{无可指摘、公正、诚实、敬畏天主的人。}\BibleRef{约 1:1} 关于达尼尔，他们甚至把他当成神；但他说，他们不会救自己的儿女。因此，知道这些事，我们既不可轻视圣人的祈祷，也不可将一切全推给他们：免得我们一方面懒惰、漫不经心地生活，另一方面又剥夺自己的巨大益处。而是要既恳求他们为我们祈祷、举手，又要坚持修德；这样我们才能藉着我们的主耶稣基督的恩宠和慈爱，获得应许给爱祂之人的福乐，偕同祂，等等。\par

\SourceRecord{Translated by John A. Broadus.；Nicene and Post-Nicene Fathers, First Series；Vol. 13.；Edited by Philip Schaff.；Buffalo, NY: Christian Literature Publishing Co.,；1889.；<http://www.newadvent.org/fathers/230401.htm>.}

% structure: p0001 source=h1 target=section reason=page-title

\section{得撒洛尼前书第二篇讲道}

\BibleRef{得前 1:8-10}\par

\BibleQuote{因为主的言语从你们那里，不仅传到了马其顿和阿哈雅，而且你们对天主的信德也传到各处，以致不需要我们再说什么。因为他们自己传述我们怎样来到了你们那里，你们怎样离开了偶像归向了天主，为服事那生活而真实的天主，并期待他的圣子自天降下，就是他使之从死者中复活的耶稣，拯救我们脱免那将要来临的震怒。}\par

就像芬芳的香膏不把它的香气封闭在自己里面，而是远远散播，以它的香气充满空气，也传到邻近的感官；同样，卓越而令人钦佩的人也不把他们的德行封闭在自己里面，而是藉着美誉使许多人受益，使他们变得更好。这也在那时发生了。因此他说：\BibleQuote{如此你们就成了马其顿和阿哈雅所有信徒的榜样。}\BibleQuote{因为从你们那里，}他说，\BibleQuote{主的言语不仅传到了马其顿和阿哈雅，而且在各处你们的信德对天主的事已经传扬出去了。}因此，你们以教导充满了一切邻人，以奇迹充满了世界。这就是所谓\Quote{在各处}所表达的意思。他没有说你们的信德被传扬，而是说\Quote{发出响声}；如同一个地方近处都被响亮号筒的声音充满，同样，你们刚毅的报告是响亮的，且像那样发出声音，足以充满世界，并均匀地落在所有周围之人身上。因为伟大的事迹在那里更为响亮地受到庆祝，即它们发生的地方；远处固然也庆祝，但并非如此之多。\par

但是你们的情况并非如此，而是美誉的声音传到了地上的每一个角落。有人说，我们怎么知道这些话不是夸张呢？因为马其顿这个民族，在基督来临之前，就已经是名扬四海，比罗马人更受赞誉。罗马人也因将他们俘获而受钦佩。因为马其顿王的功绩超过了一切传闻，他从一个小城出发，却征服了世界。因此先知也看见他，一只带翅的豹子，显示他的迅捷、猛烈、如火的本性，仿佛突然之间用他胜利的战利品飞越了整个世界。他们说，他听到一个哲学家说有无穷的世界时，他痛苦地叹息，因为当世界无数时，他连一个都未征服。他是如此高傲、雄心勃勃，且处处闻名。随着国王的声誉，民族的荣耀也并驾齐驱。因为他被称为\Quote{亚历山大，马其顿人}。所以在那里发生的事情自然也广为谈论。因为与显赫之人有关的事无法被隐藏。因此，马其顿人不逊于罗马人。\par

这也源于他们的热忱。因为好像他在谈论某种有生命之物，他引入了\Quote{传扬出去}这个词；他们的信德是如此刚毅而有力。所以他说：\BibleQuote{我们无需说什么，因为他们自己传述我们在你们中间有怎样的进入。}他们不等我们听讲述，而是那些没有在场、没有看见的人，抢先于那些在场并看见你们善行的人。你们的善行处处藉着报告变得如此明显。因此，我们无需通过讲述你们的行动来激发他们同样的热诚。因为他们本该从我们这里听到的事，他们自己却在谈论，抢在我们前面。然而，在这种情况下，常常有嫉妒，但事情的极度伟大甚至战胜了嫉妒，他们成了你们争战的传报者。尽管被抛在后面，他们仍不沉默，反而抢先于我们。他们如此行事，不可能不信我们的报告。\par

\Quote{我们在你们中间有怎样的进入}是什么意思？那就是它充满了危险和无数死亡，但你们没有为这些事所困扰。如同没有发生任何事一样，你们依恋我们；好像你们没有遭受任何苦难，而是享受了无限的福乐，你们在这一切之后如此接待了我们。因为这是第二次进入。他们往\Place{贝洛雅}去，受迫害，之后他们来到时，人们仍然接待他们，好像他们也因这些人而受荣耀，甚至为他们舍命。\Quote{我们进入}这句话很复杂，包含对你们和我们的赞美。但他自己把这转化为他们的益处。他说：\BibleQuote{你们怎样离开了偶像归向了天主，为服事那生活而真实的天主}；也就是说，你们乐意地、热切地做了这事，不需要太多劳苦就能使你们归向。\Quote{为了服事，}他说，\Quote{那生活而真实的天主。}\par

这里他也引入了劝勉，这是为了使他的话语不那么刺耳的一部分。\BibleQuote{并期待，}他说，\BibleQuote{他的圣子自天降下，就是他使之从死者中复活的耶稣，拯救我们脱免那将要来临的震怒。}\Quote{并期待，}他说，\Quote{他的圣子自天降下}；被钉的那位，被埋葬的那位；从天上期待他。如何\Quote{从天上}呢？\Quote{就是他使之从死者中复活。}你同时看到了一切：复活、升天、第二次来临、审判、义人的赏报、恶人的惩罚。\Quote{耶稣，}他说，\Quote{拯救我们脱免那将要来临的震怒。}这既是安慰，也是劝勉和鼓励。因为如果他使他从死者中复活，他在天上，并将会从那里来，（而且你们相信了他；因为如果你们没有相信他，你们就不会遭受那么多），这本身就是足够的安慰。这些人将受惩罚，他在第二封书信中这样说，而你们将有不少的安慰。\par

\Quote{等待，}他说，\Quote{祂的圣子从天降下。}可怕之事已在眼前，但美好之事尚在未来，就是基督从天降下之时。请注意，何等大的望德是必需的：被钉者已复活，被提升天，祂将审判生者死者。\par

第二章第1-2节：\Quote{你们自己，弟兄们，知道我们进入你们中间，并不是徒然的；但我们在你们那里，曾在斐理伯受苦、受辱，就如你们所知道的，我们仍靠我们的天主放胆，在很大的争战中，把天主的福音传给你们。}\par

的确，你们的行为也很伟大，但我们也并未依靠人的话语。但他上面所说的，在这里又重复了，为要从两方面显出宣讲的性质：从奇迹、从宣讲者的决心、以及从接受者的热忱与激情。他说：\Quote{你们自己，知道我们进入你们中间，并不是徒然的}，也就是说，不是按照人的方式，也不是普通的。因为我们刚从大危险、死亡和鞭打中出来，就立即陷入危险。\Quote{但是，}他说，\Quote{我们曾在受苦、受辱之后，就如你们所知道的，在斐理伯，我们仍靠我们的天主放胆。}你看，他如何又把一切归于天主？\Quote{把天主的福音传给你们，在很大的争战中。}不能说：在那里我们确实有危险，但在这里却没有；你们自己也晓得，危险多么大，我们如何在你们中间带着许多争斗。他在给格林多人的书信中也说：\Quote{我在你们那里，又软弱、又恐惧、又甚战栗。}\BibleRef{格前2:3}\par

第3-4节：\Quote{因为我们的劝勉不是出于错误，也不是出于不洁，也不是出于诡诈：但既然我们蒙天主考验合格，被托付了福音，我们就如此讲论，不是要讨人喜欢，而是要讨那察验我们心的天主的喜欢。}\par

你看，正如我所说，从他们的坚忍中，他们证明宣讲是神圣的。因为若不是如此，若是欺骗，我们就不会忍受那么多危险，那些危险甚至不让我们喘息。你们在苦难中，我们也在苦难中。那么，这又是什么呢？若不是将来的事激动我们，若不是我们确信有美好的希望，我们不会因受苦而更加奋发。谁会为了眼前的事物而选择忍受如此多的痛苦，生活在焦虑和充满危险中呢？他们能说服谁呢？这些事本身不是就足以使门徒不安吗？当他们看见自己的老师身处危险时？但你们却不是这样。\par

\Quote{因为我们的劝勉，}也就是我们的教导，\Quote{不是出于错误。}他说，这事不是诡诈或欺骗，以致我们应放弃。它不是出于可憎之事，如同变戏法和巫术的伎俩。\Quote{也不是出于不洁，}他说，\Quote{也不是出于诡诈，}也不是为了像特乌达那样的叛乱。\Quote{但既然我们蒙天主考验合格，被托付了福音，我们就如此讲论，不是要讨人喜欢，而是要讨天主喜欢。}你看，这不是虚荣吗？\Quote{但那察验我们心的天主，}他说。我们行事不是为讨人喜欢。他说，我们为了谁而做这些事呢？然后称赞他们之后，他说，不是为讨人喜欢，也不是寻求从人来的荣耀，他加上：\Quote{但既然我们蒙天主考验合格，被托付了福音。}除非祂看见我们摆脱了一切世俗的考虑，祂不会拣选我们。因此，正如祂考验我们合格，我们仍然如此，如同\Quote{蒙天主考验合格}的人。祂在何处考验我们合格，并托付我们福音？我们在天主面前显为合格，所以我们仍然如此。我们被托付了福音，这是我们的德行证明；若我们有什么不好，天主就不会考验我们合格。但\Quote{考验我们合格}的说法，在此并不表示审查。而是，我们通过考验所作的，祂在没有考验的情况下作了。就是说，祂发现我们合格，并信任我们，我们就如此讲论；正如那些被考验合格并被托付、足以担当福音的人应该做的，我们就如此讲论，\Quote{不是要讨人喜欢}，也就是说，我们做这一切不是为你们。因为他先前称赞了他们，为使他的言论不引起怀疑，他说：\par

第5-6节：\Quote{因为我们从来没有用过谄媚的话，这是你们知道的，也没有藏过贪心的借口，天主可以作证；也没有寻求从人来的荣耀，无论是从你们，还是从别人，虽然我们作为基督的宗徒，本可加重你们的负担。}\par

因为他说：\Quote{无论何时，\InnerQuote{我们都没有用过谄媚的话}；}那就是说，我们没有谄媚，这是欺骗者的做法，他们想要获得控制权并统治。没有人能说我们谄媚是为了统治，也没有人能说我们为此寻求财富。对于这明显的事，他后来请他们作见证。他说：\Quote{\InnerQuote{我们是否谄媚，}你们知道。}但对于不确定的事，即是否出于贪财，他呼求天主作证。\Quote{\InnerQuote{也不求从人来的荣耀，无论是从你们，还是从别人，虽然我们作为基督的宗徒，本可成为负担；}}那就是说，我们也不寻求尊荣，不自夸，也不要求随从护卫。即使我们这样做了，也并非不合身份。因为如果君王所派遣的人尚且受到尊荣，何况我们呢？他没有说\Quote{\InnerQuote{我们被藐视}}，也没有说\Quote{\InnerQuote{我们没有享受尊荣}}，那会是指责他们，而是说\Quote{\InnerQuote{我们没有寻求它们}}。因此，我们本可寻求尊荣却没有寻求，甚至在传道需要时也不寻求，我们怎么会为了荣耀而行事呢？即使我们寻求了，那也无可指责。因为那些由天主派遣、如同从天上而来的使者，理应享受极大的尊荣。\par

但我们以极大的宽容，不做这些事中的任何一件，为要堵住对手的口。不可说，我们对你们如此，对别人却不这样。因为他在致格林多人的书信中也这样说：\Quote{\InnerQuote{因为如果有人奴役你们，吞吃你们，掳掠你们，高举自己，打你们的脸，你们尚且容忍他。}}\BibleRef{格后 11:20}又说：\Quote{\InnerQuote{他亲临的软弱，言语也微不足道。}}\BibleRef{格后 10:10}又说：\Quote{\InnerQuote{请你们原谅我这一次的不当。}}\BibleRef{格后 12:13}他在那里也表明，因遭受如此多的事，他是极其谦卑的。但在这里，他也论及钱财说：\Quote{\InnerQuote{虽然我们作为基督的宗徒，本可成为负担。}}\par

第7、8节。\Quote{\InnerQuote{但我们在你们中间是温柔的，如同乳母抚育自己的儿女：我们如此眷爱你们，不但乐意将天主的福音分给你们，连自己的性命也愿意给你们，因为你们已成为我们所疼爱的。}}\par

他说：\Quote{\InnerQuote{但我们是温柔的}}；我们没有表现出任何冒犯或麻烦的事，没有令人不悦或自夸的事。而\Quote{\InnerQuote{在你们中间}}这句话，就如同说，我们如同你们中的一员，不取更高的地位。\Quote{\InnerQuote{如同乳母抚育自己的儿女}}。教师应当如此。乳母会为了获得荣耀而谄媚吗？她会向小孩子要钱吗？她会对他们冒犯或成为负担吗？她们难道不是比母亲更宽容吗？这里他显明了他的情感。\Quote{\InnerQuote{我们如此眷爱你们}}，他说，我们对你们如此依恋，他说，我们不仅不从你们那里取什么，甚至如果需要把我们的灵魂分给你们，我们也不会拒绝。请告诉我，这难道出于人的看法吗？谁这样愚蠢，竟说这话？他说：\Quote{\InnerQuote{我们不但乐意将天主的福音分给你们，连自己的性命也愿意给你们}}。所以这比那更大。有什么益处呢？因为福音是益处，但给出我们的性命，在困难方面比那更大。因为仅仅宣讲与舍命并不相同。那确实更宝贵，但后者是更难的事。他说，我们愿意，如果可能，甚至为你们付出性命。我们也本应愿意这样做；因为若我们不愿意，就不会忍受这必要。既然他先前称赞，现在也称赞，所以他这样说：我们做这些事，不是为求钱财，不是谄媚你们，也不是贪图荣耀。请注意：他们曾多方争战，因而应当受褒扬，甚至超乎寻常地受赞美，好使他们更坚定；这赞美是可疑的。因此他说这一切，为消除这猜疑。他也提到危险。而且，为避免人以为他如此说危险是因劳苦，并要求他们尊重他，因此当他不得不提到危险时，他又加上\Quote{\InnerQuote{因为你们已成为我们所疼爱的}}；我们情愿为你们舍命，因为我们深爱你们。我们宣讲福音，是因天主命令；但我们如此爱你们，以至于如可能，我们甚至愿将性命给你们。\par

凡爱人的，应当如此爱，即使被要求舍命，而且可能的话，也不拒绝。我不说\Quote{\InnerQuote{若被要求}}，而是说他甚至主动跑去献上礼物。因为没有什么比这样的爱更甜美；没有什么会落在那里是痛苦的。的确，\Quote{\InnerQuote{忠实的朋友，是生命的良药。}}\BibleRef{德 6:16}的确，\Quote{\InnerQuote{忠实的朋友，是坚固的保障。}}\BibleRef{德 6:14}因为一个真正的朋友什么不能做到？他不能提供什么快乐？什么益处？什么安全？即使你举出无限的宝藏，也没有一样能与真正的朋友相比。首先，让我们谈谈友谊本身的大喜乐。朋友因见到朋友而欢喜，并洋溢着喜乐。他与朋友心灵合一，产生难以言喻的快乐。他只要想起他，心灵就振奋，欣喜若狂。\par

我所说的乃是真正的朋友，那些同心合意、甚至愿意为彼此而死、热切相爱的人。请不要因想到那些仅存几分爱意、只是同席的伙伴、有名无实的朋友，就以为我的话被驳倒了。若有人拥有如我所说的朋友，他必承认我言为真。这样的人虽日日见其朋友，仍不感餍足。他为朋友祈求的事，与为自己祈求的相同。我认识一人，他为朋友呼求圣人，恳请他们先为朋友祈祷，然后为自己。良友如此珍贵，以致时间和地点也因他而可爱。正如发光的身体将光芒洒向邻近之处，朋友也将他们自己的恩宠留在所到之处。常常在朋友缺席时，我们站在那些地方，不禁流泪，追忆共同度过的时日，叹息不已。言语无法表达交往所带来的快乐有多大。唯有亲历者方知。向朋友请求恩惠或接受恩惠，皆无猜疑。当他们有所嘱托时，我们便觉亏欠；但当他们迟迟不这样做时，我们便忧伤。我们所有的一切，无不属于他们。我们常常轻看此世的一切，却因他们而不愿离去；他们比光明更令我们渴慕。\par

诚然，朋友比光明更令人渴慕——我说的是真正的朋友。不必惊奇：因为宁可太阳熄灭，也不愿失去朋友；宁可在黑暗中生活，也不愿没有朋友。我将告诉你为何如此。因为许多看见太阳的人仍处于黑暗，而那些拥有朋友的人，即使身处患难，也永不陷于黑暗。我所说的乃是属灵的朋友，他们不将任何事物置于友谊之上。保禄就是如此，他甘愿献出自己的灵魂，即使无人要求，甚至愿意为他们下地狱。我们应当以如此炽热的心去爱。\par

我愿给你们一个友谊的例证。朋友——即按照基督的朋友——胜过父亲和儿子。请不要向我提及当今的朋友，因为这一善物也已随其他事物消逝。但请思想宗徒时代，我并非指那些首领，而是泛指信徒们；他说：\BibleQuote{一切，都同心合意，没有一人说他的财物中有什么是自己的……照各人所需用的分给每人。}\BibleRef{宗 4:32}那时没有\Quote{我的}和\Quote{你的}这类词语。这就是友谊：人不应视自己的财物为私有，而应视为邻人的；他的所有物属于他人；他应如同照顾自己的灵魂一样照顾朋友的灵魂，朋友亦然。\par

有人会说：在哪里能找到这样的人呢？因为我们没有意愿；但这是可能的。若不可能，基督就不会命令我们，也不会如此详细论述爱了。友谊是伟大的，其伟大无人能学尽，也无言语能表达，唯有经验本身才能显明。正是它导致了异端。它使希腊人成为希腊人。爱者不愿命令，也不愿统治，反倒因受统治和命令而感恩。他更愿施恩而非受恩，因为他爱，且感到尚未满足他的渴望。他因受惠而得的喜悦，不及施惠时的喜悦。他宁愿施惠，而非欠债；或者说，他既愿受惠于人，也愿人欠他。他既愿施恩，又看似未施恩，而自己反为负债者。我想或许你们中许多人未能理解我所说的话。他愿在施恩中居先，却不显得居先，反而像是在还报恩情。天主对待人类也是如此。他定意将自己的圣子为我们舍出，但为不显得施恩，而像是亏欠我们，便命令亚巴郎献上他的儿子，如此，在做极大善事时，他看似并未行什么大事。\par

因为当没有爱时，我们以恩惠责备人，并夸大微小之事；但当有爱时，我们便隐藏这些事，并愿使大事显得微小，以免让人觉得朋友欠我们，反倒我们自己因朋友欠我们而成为负债者。我知道大多数人不能理解我所说的，其原因是，我现在所说之事如今已居于天上。因此，若我谈论某种生长在印度的植物，无人曾经验过，即使我说上万言，言语也无法描述它。同样，现在我所说的，都是徒然，因为无人能理解我。这是一株栽种在天上的植物，它的枝条并非沉甸甸的珍珠，而是德行的生活，远比珍珠更悦人。你想说什么快乐，是污秽的还是可敬的？友谊的快乐超越一切，即使你谈论蜂蜜的甘甜。因为蜂蜜会使人饱腻，但朋友永远不会——只要他仍是朋友；对他的渴望反而增加，这样的享乐永远不会令人厌烦。朋友比现世的生命更甜美。因此，许多人在朋友去世后不愿再活下去。有朋友相伴，即使流放也能忍受；没有朋友，即使居住本国也不愿选择。有朋友，贫穷也可忍受；没有朋友，健康与财富都不可忍受。他是另一个自我：我受困，因为我无法举例证明。因为若能用例子，就会显出我所说的远不及当说的。\par

这些事在这里确实如此。但来自天主的友谊的赏报如此巨大，无法言传。祂赐予赏报，好使我们彼此相爱，而相爱正是我们理当为之付出赏报的事。\Quote{祈祷吧，}祂说，\Quote{并领受赏报，}因为那是我们理当为之付出赏报的事，因为我们求善事。\Quote{因为你所求的，}祂说，\Quote{领受赏报。守斋，并领受赏报。行善，并领受赏报，}尽管你反倒欠着赏报。但是，如同父亲们，当他们使儿女成为有德之人后，还进一步给予他们赏报；因为他们是欠债的，因为他们给了儿女快乐；天主也是如此行事。\Quote{领受赏报吧，}祂说，\Quote{如果你有德，因为你令你的父亲喜悦，为此我欠你一份赏报。但如果你邪恶，就不是这样：因为你激怒了生你的那位。}那么，我们不要激怒天主，而要令祂喜悦，好使我们能获得天国，在我们的主耶稣基督内，愿光荣与德能归于祂，直到永远。阿们。\par

\SourceRecord{Translated by John A. Broadus.；Nicene and Post-Nicene Fathers, First Series；Vol. 13.；Edited by Philip Schaff.；Buffalo, NY: Christian Literature Publishing Co.,；1889.；<http://www.newadvent.org/fathers/230402.htm>.}

% structure: p0001 source=h1 target=section reason=page-title

\section{得撒洛尼前书第三篇讲道}

得撒洛尼前书 2:9-12\BibleRef{得前 2:9-12}\par

\Quote{弟兄们，你们原记得我们的劳苦和勤劳：我们日夜操作，为不加重你们任何人的负担，而向你们宣讲了天主的福音。你们和天主都可作证，我们怎样以圣洁、正义和无可指摘的态度对待你们信友。你们也知道，我们怎样如同父亲对待自己的儿女一样，分别劝勉、鼓励和恳求你们每一位，使你们的行事为人，配得上那召你们进入 祂 的国和光荣的天主。}\par

老师为使学生得救，不该有任何负担感。如果蒙福的\Person{雅各伯}日夜护守羊群而劳累，那么受托照顾灵魂的人，更应忍受一切辛劳，即使工作艰苦卑贱，只须着眼一件事：学生的得救，以及由此归于天主的光荣。请看\Person{保禄}，这位宣讲者、世界的宗徒、被高举到如此大荣的人，亲手劳作，为不使门徒受累。\par

\Quote{你们原记得}，他说，\Quote{我的弟兄们，我们的劳苦和勤劳。}他先前说过：\Quote{我们作为基督的宗徒，本可加重你们的负担}，正如他在致格林多人书里也说：\Quote{你们不知道吗？那些在圣所供职的人，就吃圣所中的物；同样，主也命定传福音的人，靠福音生活。} \BibleRef{格前 9:13} 但我说，我不愿意，我却劳苦了；他不仅工作，而且十分勤勉。请看他怎么说：\Quote{你们原记得}，他没有说我从你们得到的益处，而是\Quote{我们的劳苦和勤劳：我们日夜操作，为不加重你们任何人的负担，而向你们宣讲了天主的福音。} 他致格林多人书时说了另一番话：\Quote{我掠夺了别的教会，取了他们的工资，为能服事你们。} \BibleRef{格后 11:8} 即使在那里他也工作，但他没有提及这一点，却强调了更惊人的事，好像说：我服事你们时，是受别人供养的。但这里并非如此。那么是什么呢？\Quote{日夜操作。} 而且在那里他说：\Quote{当我在你们那里时，虽有所缺乏，却没有拖累任何人}，以及\Quote{我取了工资，为能服事你们。} \BibleRef{格后 11:8} 这里他表明那些人贫穷，但那里并非如此。\par

为此他屡次请他们作证。因为\Quote{你们是见证人}，他说，\Quote{天主也是}；天主是可信的，但这另一件事最能叫他们确信。对于不知情的人，那件事不确定；但对所有人，这却是无疑的。不要问这是否保禄说的。他超出所需地给他们保证。因此他说：\Quote{你们是见证人，天主也是：我们怎样以圣洁、正义和无可指摘的态度对待你们信友。}\par

\Quote{正如你们知道，我们怎样劝勉、安慰你们每一位，如同父亲对待自己的儿女一样。} 上面说了他的行为，这里他说他的爱，这爱超过他管理他们的职权。所说的话也显明他不骄傲。\Quote{如同父亲对待自己的儿女，劝勉你们，鼓励你们，恳求你们，使你们的行事为人，配得上那召你们进入 祂 的国和光荣的天主。} 当他说\Quote{恳求你们}时，他就提到了\Quote{父亲}；我们虽恳求，但不是强求，而是如同父亲。\Quote{你们每一位。} 怪哉！在如此众多的会众中，不忽略任何人，无论大小贫富。\Quote{劝勉}你们，他说：要容忍。\Quote{安慰和恳求。} \Quote{劝勉}，所以他们不寻求光荣；\Quote{恳求}，所以他们不奉承。\Quote{使你们行事为人，配得上那召你们进入 祂 的国和光荣的天主。} 再看，他在叙述中如何既教导又安慰。因为 祂 既召他们进入 祂 的国，既召他们进入光荣，他们就该忍受一切。我们\Quote{恳求}你们，不是要你们施恩于我们，而是要你们获得天国。\par

第13节。\Quote{为此，我们也不断地感谢天主，因为你们由我们接受了所听的天主圣言，并没有拿它当人的言语，而实在当天主的圣言领受了，这言语在你们信友中发生效力。} \BibleRef{得前 2:13}\par

不能说，他说，我们固然行一切无可指摘，但你们却做了与我们生活不相称的事。因为你们听我们时，所留意的仿佛不是听人，而是听天主亲自劝勉。这从何显明呢？因为他从自己的试炼和他们的见证，以及他的行事方式，证明他没有用谄媚或虚荣宣讲；同样，从他们的试炼，他也证明他们正确领受了圣言。因为他说，除非你们听时好像天主在说话，你们怎能忍受那样的危险呢？请看他的尊严。\par

第14-16节。\Quote{因为你们，弟兄们，曾效法了在基督耶稣内、在犹太的各天主教会，因为你们也由本乡人受了苦，正如他们由犹太人受了苦。犹太人杀了主耶稣和他们的先知，又驱赶了我们；他们不悦乐天主，且与众人为敌，阻止我们向外邦人讲道，使他们得救，以致常填满自己的罪恶。但天主的义怒已降到他们身上到了极点。} \BibleRef{得前 2:14-16}\par

\BibleQuote{因为你们，}他说，\BibleQuote{成了在犹太的天主的教会的效法者。}这是一种极大的安慰。他说，他们向你们做这些事不足为奇，因为他们对本国人也做了同样的事。这也不小地证明了宣讲的真实性，即连犹太人也经受得住一切。\BibleQuote{因为你们也，}他说，\BibleQuote{从自己的同胞那里受了同样的苦，正如他们从犹太人那里所受的。}他说\BibleQuote{他们在犹太也同样做了}，其中还有更深的意义；它表明他们到处喜乐，就像高尚地争斗过一样。因此他说，\BibleQuote{你们也受了同样的苦。}再说，如果他们对主也敢于做这样的事，那对你们又有什么可奇怪的呢？\par

你看他如何提出这一点作为极大的安慰？他不断地提到它；仔细检查可以发现几乎在他所有的书信中，在各种情况下，在一切诱惑的场合，他都引述基督。因此注意，在这里也是如此，当他谴责犹太人时，他提醒他们想起主和主的苦难；他深知这是最大的安慰。\par

\BibleQuote{他们既杀了主，}——但也许他们不认识他——他们当然认识他。那么怎样？他们不是也杀了并用石头打了自己的先知吗？甚至他们的书卷还随身携带。他们这样做不是为了真理。因此，这不仅是诱惑中的安慰，而且他们被提醒不要认为（犹太人）是为了真理而这样做，并因此困扰。\BibleQuote{又驱逐了我们，}他说。我们也，他说，遭受了无数灾祸。\BibleQuote{不悦乐天主，且与众人为敌；禁止我们对外邦人说话，使他们得救。}\BibleQuote{与众人为敌，}他说。怎样？因为如果我们应当对世界说话，而他们禁止我们，他们就是世界共同的敌人。他们杀了基督和先知，侮辱天主，是世界的共同敌人，他们驱逐我们——我们是来为他们的得救。如果他们甚至在犹太也做了这样的事，那他们对你们做了这些又有什么奇怪呢？\BibleQuote{禁止我们对外邦人说话，使他们得救。}因此，阻碍众人的得救是嫉妒的表现。\BibleQuote{常常充满自己的罪过。但义怒已临到他们身上到极点。}\BibleQuote{到极点}是什么？这些事不再像从前。这里没有回转，没有界限。义怒近在眼前。这从何显现？从基督所预言的。因为不仅有同受苦难者的安慰，而且听到迫害者将受罚也是一种安慰。如果拖延是烦恼，那么他们将永不再抬头就是安慰；或者说，他缩短了拖延，通过说\BibleQuote{义怒}，表明这早已应得、预定的和预言的。\par

\BibleRef{得前 2:17} \BibleQuote{但我们，暂时离开你们——在面容上，不在心里——越发竭力地、怀着极大的渴望要见你们的面。}\par

他没有说\BibleQuote{分离}，而是说了更重的。他之前谈到谄媚，表明他没有谄媚，没有寻求荣耀。这里他谈到爱。因为他之前说过\BibleQuote{如同父亲对自己的儿女}，\BibleQuote{如同乳母}，这里他用了另一个说法：\BibleQuote{成了孤儿}，这是指失去父亲的孩子们。然而他们成了孤儿。\BibleQuote{不}——他说——\BibleQuote{而是我们}。因为如果有人审视我们的渴望，就像无保护的小孩子，遭受了过早的丧亲，渴望父母，不仅出于自然情感，而且由于他们的被遗弃状态，我们也是如此。由此他也表明自己对分离的沮丧。他说，我们不能说我们已经等了很长时间，而是\BibleQuote{短时间内}，并且是\BibleQuote{在面容上，不在心里}。因为我们心里总有你们。看他的爱何其大！虽然心里总有他们，他也寻求面对面相见。不要对我讲你高深的哲学！这是真正的炽热之爱：既看见，又听见，又说话；这可能大有裨益。\BibleQuote{我们越发竭力。}\BibleQuote{更加}是什么意思？他或者要说\BibleQuote{我们强烈地依恋你们}，或者\BibleQuote{既然暂时被剥夺，我们设法见你们的面}。看看蒙福的保禄。当他不能亲自满足他的渴望时，他通过别人来做，比如当他派弟茂德去斐理伯人那里，又派同一个人去格林多人那里，当他不能亲自时，就通过别人与他们交往。因为在爱他们上，他像某个疯狂的人，无法自制，无法控制自己的情感。\par

\BibleRef{得前 2:18} \BibleQuote{因此，我们曾愿意到你们那里去。}\par

这是爱的表现；但这里他只提到一个需要，即\BibleQuote{为了见你们。} \BibleQuote{我保禄，一次再次地，但撒殚阻止了我们。}\par

你怎么说？是撒殚阻碍吗？是的，的确，因为这不是天主的工作。因为在\BibleBook{罗马}{书}中，他说了这一点，即天主阻碍了他\BibleRef{罗 15:22}；而在别处，路加说，\Quote{圣神}阻止了他们进入亚细亚。\BibleRef{宗 16:7} 在\BibleBook{格林多}{前书}中，他说这是圣神的工作，但这里只是撒殚的工作。但他说的是撒殚的什么阻碍呢？一些意想不到的和猛烈的诱惑：因为经上说，犹太人正设谋陷害他，他就在希腊被拘留了三个月。但为了救恩计划而甘愿留下是另一回事。因为那里他说：\Quote{因此，我在这些地方再没有可做之处了}\BibleRef{罗 15:23}，以及\Quote{为了宽免你们，我忍住了未到格林多去。}\BibleRef{格后 1:23} 但这里却没有这类事。而是什么呢？即\Quote{撒殚阻碍}了他。\Quote{我保禄}，他说，\Quote{一次又一次地}。请注意，他是多么热切，以及他如何彰显自己，表明他最爱他们。他说\Quote{我保禄}，意思是不管别人如何。因为别人只是愿意，而我甚至尝试了。\par

第19节。\Quote{因为我们的希望、或喜乐、或夸耀的冠冕是什么？不就是你们在主耶稣基督来临之时，在他面前吗？}\BibleRef{得前 2:19}\par

告诉我，马其顿人，他们就是你的希望吗，有福的保禄？不只有他们，他说。因此他加了\Quote{你们不也是吗？}因为他说：\Quote{我们的希望、或喜乐、或夸耀的冠冕是什么？}请注意，这些话像是满怀柔情的妇女对她们幼小的孩子说话。他说\Quote{夸耀的冠冕}。因为\Quote{冠冕}这名称不足以表达辉煌，还加了\Quote{夸耀的}。这是何等的热忱！无论是母亲还是父亲，即使他们聚在一起，混合他们的爱，也绝不可能表现出与保禄相等的爱。他说：\Quote{我的喜乐和冠冕}，意思是，我在你们身上的喜乐胜过在冠冕中的喜乐。试想，一整个由保禄栽种和扎根的教会出现在眼前，这是何等的大事。谁会不为如此众多的子女以及这些子女的良善而喜乐呢？因此这也不是奉承。因为他没有说\Quote{你们}，而是说\Quote{你们也}，与其他人一同。\par

第20节。\Quote{因为你们就是我们的荣耀和我们的喜乐。}\BibleRef{得前 2:20}\par

第3章1至2节。\Quote{因此，我们再也忍耐不住了，就决定我们独自留在雅典。}意思是\Quote{我们选择了}。\Quote{打发我们的兄弟弟茂德，就是天主的仆役，并我们在基督福音上的同工。}\par

他说这话，不是为了夸耀弟茂德，而是为了尊敬他们，因为他把同工和福音的仆人差遣给了他们。好像他说：我们把他从劳苦中抽出来，派了天主的仆役并在基督福音上的同工到你们那里去。\par

\Quote{为坚固你们，并安慰你们的信德。}\par

第3节。\Quote{免得有人因这些困苦而动摇。}\BibleRef{得前 3:3}\par

那么他在这里说什么呢？因为教师的诱惑扰乱了他们的门徒，而他当时落入了许多诱惑之中，正如他自己所说：\Quote{撒殚阻碍了我们}，一直这样说；他说：\Quote{我一次再次想要到你们那里去}，却未能成行，这是极大暴力的证明。而这一点理所应当使他们困扰，因为他们受到自己诱惑的困扰，不如看到自己的教师受伤那样严重；正如士兵不被自己的试练所困扰，不如看到将军受伤那样严重。他说：\Quote{为坚固你们}；并不是说他们在信德上有什么欠缺，也不是说他们需要学习什么。\par

\Quote{并安慰你们的信德，免得有人因这些困苦而动摇；因为你们自己知道，我们是为此而定的。}\par

第4节。\Quote{因为我们在你们那里时，已预先告诉你们，我们将受迫害；正如所发生的，你们也知道。}\BibleRef{得前 3:4}\par

你们不该，他说，困扰，因为没有什么奇怪、没有什么出乎意料的事发生；这足以使他们兴起。因为你看，基督也因此预先告诉了他的门徒。请听他说：\Quote{现在事情发生以前，我告诉了你们，使你们在事情发生时，能够相信。}\BibleRef{若 14:29} 因为大大地，确实大大地有助于安慰别人，就是预先从他们的教师那里听到将要发生的事。因为如同一个病人，如果从医生那里听说这样或那样的情况正在发生，就不会太困扰；但如果有什么意料之外的事发生，就好像连医生也束手无策，病势超出了他的医术，他就痛苦并困扰；这里也是如此。保禄预知这一点，预先告诉了他们：\Quote{我们将受迫害}，\Quote{正如所发生的，你们也知道。}他不仅说这事发生了，而且说他预知了许多事，并且它们发生了。\Quote{我们是为此而定的。}所以你们不仅不该为过去的事困扰和不安，即使将来有类似的事发生，也不该，\Quote{因为我们是为此而定的。}\par

伦理教训。让我们侧耳倾听。基督徒正是为此被召选的。因为关于所有信友，都说了这话：\Quote{我们正是为此被召选的。}而我们，好像我们被召选是为了安逸，若遭受任何苦难便觉得奇怪；然而我们有什么理由觉得奇怪呢？因为没有一种患难或试探的时期临到我们，不是人所共有的。现在正是合适的时机对你们说：\BibleQuote{你们与罪恶争斗，还没有抵抗到流血的地步}（\BibleRef{希 12:4}）。或者更确切地说，这对我们说并不合适——但什么合适呢？你们还没有轻视财富。因为对那些确实已经失去自己一切所有的人来说，这些话说得有理，但这是对保留自己财产的人说的。谁曾为基督的缘故被抢去财富？谁曾被打？谁曾受辱？我是指言语上的侮辱。你有什么可夸耀的？你有什么勇气说什么呢？基督为我们这些仇敌受了那么多苦。我们能显明我们为祂受了什么吗？实在没有任何受苦，反而从祂那里领受了无数好处。那时我们将有何自信呢？你们不知道，士兵也是如此，当他能显示无数伤口和伤疤时，才能在君王面前荣耀。但如果他没有善行可显，即使他没有作恶，也将居于最末之列。\par

但你说，这不是战争的时期。但如果是的话，请告诉我，谁会应战？谁会攻击？谁会突破方阵？也许没有人。因为当我看到你们不肯为基督的缘故轻视财富，我怎能相信你们会轻视打击呢？告诉我，你们能勇敢地忍受那些侮辱你们的人，并祝福他们吗？你们不能——反而悖逆。没有危险的事你们不做；那么你们会忍受痛苦和苦难的打击吗？你们不知道，在和平时期也应该保持战争的操练吗？你们没有看见这些士兵，虽然没有任何战争扰乱他们，而是深刻的和平，他们擦亮武器，每天跟随教给他们战术的教师到广阔平坦的平原上，极其严格地保持战争的操练吗？我们属灵的士兵，谁这样做了？没有人。因此我们在战争中变得软弱、卑贱，容易被任何人俘虏。\par

但这是何等愚昧，当保禄呼喊说：\BibleQuote{凡愿意在基督耶稣内虔敬生活的，都要遭受迫害}（\BibleRef{弟后 3:12}），基督也说：\BibleQuote{在世界上你们要受苦难}（\BibleRef{若 16:33}），而蒙福的保禄又大声呼喊说：\BibleQuote{我们的战斗不是对抗血和肉}（\BibleRef{弗 6:12}），又说：\BibleQuote{所以要站稳，用真理束上你们的腰}（\BibleRef{弗 6:14}），你们却认为现在不是战争的时期？告诉我，为什么在没有战争的时候武装我们？为什么给我们无谓的麻烦？你在士兵可以休息和恢复的时候给他们穿上铠甲。但他会说：当然，即便不是战争，也应当关心战争的事。因为和平时期考虑战斗事务的人，在战争时期将是可怕的；但不懂战争事务的人，即使在和平时期也会更烦恼。为什么？因为他会为他所拥有的东西哭泣，由于不能为之战斗而痛苦。因为懦弱、无经验、不善战者的财产，属于所有勇敢善战的人。所以，首先我为此武装你们。但其次，我们一生的全部时间都是战争时期。怎样和在哪方面？魔鬼常常在附近。听经上所说：\BibleQuote{如同咆哮的狮子巡游，寻找可吞食的人}（\BibleRef{伯前 5:8}）。无数的身体情欲攻击我们，有必要列举出来，以免我们徒然欺骗自己。因为请告诉我，什么不与我们作战？财富、美貌、快乐、权力、权威、嫉妒、光荣、骄傲？因为不仅我们自己的光荣与我们作战，阻止我们降卑；而且他人的光荣也与我们作战，引我们嫉妒和恶意。但它们的相反者——贫穷、羞辱、被轻视、被拒绝、无权——这些也攻击我们。从人方面来的有邪恶、阴谋、欺骗、毁谤、无数的攻击。同样从魔鬼方面来的有\BibleQuote{率领者、掌权者、黑暗世界的霸主、天界里邪恶的鬼神}（\BibleRef{弗 6:12}）。我们中有些人在欢乐，另一些人在悲伤，两者都是偏离正路。但健康和疾病（与我们作战）。从哪一面人不会陷入罪呢？你们要我从头说起，甚至从亚当开始吗？是什么俘虏了最初受造的人？快乐、饮食、贪爱统治。下一个儿子呢？嫉妒和怨恨。诺厄时代的人呢？肉体的快乐以及由此产生的邪恶。他的儿子呢？傲慢和不敬。索多玛人呢？傲慢、放荡、粮食饱足。但贫穷也常有此效果。因此一位智者说：\BibleQuote{不要使我贫穷，也不要使我富裕}（\BibleRef{箴 30:8}）。然而，既不是贫穷也不是富裕，而是不能善用它们的意志。经上说：\BibleQuote{要知道你是在罗网中行走}（\BibleRef{德 9:13}）。\par

蒙福的保禄说得极好：\Quote{我们正是为此被召选的}。他没有仅仅说我们受诱惑，而是说\Quote{我们正是为此被召选的}，好像说：我们为此而生。这是我们的本分，这是我们的人生，而你们寻求安息吗？刽子手没有站在我们身旁，撕裂我们的肋旁，强迫我们献祭；但贪财和贪多的欲望正在撕毁我们的眼睛。没有士兵点燃柴堆，或把我们放在烤架上；但比这更甚的是，肉体的火焰点燃了我们的灵魂。没有君王在场应许无数赏赐，使我们羞愧。但有一种贪求光荣的狂热，比那更甚地搔痒我们。一场大战，真是，极其浩大，只要我们警醒。\par

当前这个季节也有其冠冕。且听\Person{保禄}说：\BibleQuote{从此以后，正义的冠冕已为我预备下了，就是公义的审判者在那一天要赏给我的；不但赏给我，而且也赏给所有爱慕祂显现的人。}\BibleRef{弟后 4:8}当你失去了一个你所爱且唯一的儿子，你曾用大量财富养育他，寄予厚望，他是你唯一的继承人；不要抱怨，而要感谢天主，荣耀那召去他的人，在这方面你将不逊于\Person{亚巴郎}。因为正如\Person{亚巴郎}在天主命令时将儿子献给天主，同样，当天主召去了他，你没有抱怨。你曾患重病，许多人前来劝你，有的用咒语，有的用护身符，有的用其他东西来医治这恶疾？而你因敬畏天主，坚定不屈地忍受，宁愿承受一切，也不愿行那些偶像崇拜的事？这给你带来殉道的冠冕。不要怀疑。我将告诉你如何以及通过什么方式。因为正如这样的人坚定忍受酷刑的痛苦，以免朝拜偶像，同样，你也忍受疾病的痛苦，不愿要他人提供的那些疗法，也不做他所吩咐的事。\Quote{那些痛苦更剧烈}——是的，但这些持续时间更长，所以最终是一样的；不，往往这些也更剧烈。因为告诉我，当高烧在体内肆虐燃烧，你拒绝了别人推荐的咒语，难道你没有戴上殉道的冠冕吗？\par

再者，有人丢了钱财吗？许多人劝你求助于占卜者；但你因敬畏天主，因为这是被禁止的，宁愿不要收回钱财也不愿违背天主——你就与那施舍给穷人的人获得同样的赏报，如果你丢失后感谢天主，并且当能够求助于占卜者时，你忍受着不去收回，而不愿这样收回。因为正如他因敬畏天主将一切给了穷人，同样，你也因敬畏天主，当别人抢掠了你时，你没有追回它。\par

我们是自己受伤害或不伤害的主宰。如果愿意，让我们以盗窃本身为例来说明全部情况。窃贼凿穿了墙壁，冲进了房间，拿走了贵重的金器、宝石，总之，他清空了你的全部宝藏，并且没有被抓获。事实令人痛心，看似是损失；但仍未成定局，而是取决于你使之成为损失或收益。如何能成为收益呢？我会尽力向你展示如何成为巨大收益，如果你愿意；但如果不愿意，损失将比已经发生的更为严重。因为正如工匠，当材料摆在他们面前时，技艺精湛的人善用之，而技艺不精的人则糟蹋它，使之成为他的损失，在这些事上也是如此。那么如何成为收益呢？如果你感谢天主，如果你不痛哭流涕，如果你说出\Person{约伯}的话：\BibleQuote{天主赐的，天主收回。我赤身出离母胎，也要赤身归去。}\BibleRef{约 1:21}\par

你说：\Quote{什么？}你说：\Quote{天主取走了？窃贼取走了，你怎么能说天主取走了？}不要奇怪，因为就连\Person{约伯}，对于魔鬼取走的东西，他说：这些是天主取走的？难道你不也该说，窃贼取走的东西，是天主取走的吗？告诉我，你钦佩谁？是将所有财物施舍给穷人的那位，还是说这些话的\Person{约伯}？那没有当时施舍的人，比那施舍了的人差吗？不要说：\Quote{我不觉得感恩。这件事并非出于我的同意、知晓或意愿。强盗取走了它。我有什么赏报呢？}这些事也不是出于\Person{约伯}的知晓或意愿。怎么会呢？然而，他奋斗了。\par

而你有能力获得同样大的赏报，如同你自愿丢弃它一样。或许我们更钦佩这个感恩地承受冤屈的人，超过那个自发施舍的人。为什么？因为后者确实被赞美滋养，得到良心的支持，并有美好的希望；他先前勇敢地忍受了财物的匮乏，然后才丢弃它们；但是前者，还依恋着它们时，被强行剥夺了。这与先被说服割舍财富然后那样施舍，是不同的，与还在渴望时被剥夺也是不同的。如果你说出这些话，你将得到比\Person{约伯}多数倍甚至更多的赏报。因为他在此得到双倍，但\Person{基督}许诺你百倍。因敬畏天主，你没有亵渎吗？你没有求助于占卜者吗？受冤屈时，你感恩了吗？你就像一个轻视财富的人，因为若非你先轻视了它，你做不到这一点。长久操练轻视财富，与突然承受已发生的损失，不是同一回事。这样，损失变成了收益，你将不受伤害，反而因魔鬼得利。\par

但损失又如何成为沉重的呢？当你失去你的灵魂时！告诉我，窃贼夺走了你的财产：你难道要夺走自己的救恩吗？因此，你因受他人加害而悲伤，为何要将自己投入更大的恶中呢？他或许使你陷入贫困；但你却悖逆地伤害自己于致命之事。他夺走了你身外之物，那些以后会违背你意愿而离开你的东西。但你自己却夺走了永久的财富。魔鬼因拿走你的财富而令你悲伤；你也要令它悲伤，不要取悦它。如果你求助于占卜者，你取悦了它。如果你感谢天主，你给了它致命的一击。\par

且看所发生的事。你若去找占卜者，仍然找不到，因为知道这事不在他们的能力之内；即使他们偶然告诉了你，你不仅丧失自己的灵魂，还会被弟兄们嘲笑，并且最终悲惨地失去它。因为魔鬼知道你承受不了损失，甚至为这些事否认你的天主，便再次赐给你财富，好让他有机会再次欺骗你，使你跌倒。若占卜者告诉你，不必惊奇。魔鬼是无形的，他四处游荡。正是他武装了强盗本身。因为这些事情不是没有魔鬼参与的。因此，如果他武装他们，他也知道财物存放在哪里。他并非不熟悉自己的仆役。这并不奇怪。若他看见你因损失而悲伤，就会再加添另一个损失给你。若他看见你嘲笑它、轻视它，他就会停止这样做。因为正如我们对待仇敌那样，用能令他们悲伤的事物攻击他们，但若我们看到他们不悲伤，我们便会停止，因为无法折磨他们；魔鬼也是如此。\par

你说什么？你难道没有看见那些在海上航行的人，当风暴来临时，他们不顾财富，甚至把货物扔到海里？\Quote{人啊，你说什么？你是在与风暴和船难合作吗？在波浪夺走你的财富之前，你就亲手把它扔掉了？为什么在船难之前，你先使自己沉船？}但这可能是一个没有经历过海上风浪的乡野之人说的话。而航海者，真正知道什么是风平浪静的原因、什么是风暴的原因的人，甚至会嘲笑说这种话的人。因为他说：\Quote{我把它扔到海里，这样大海就不会淹没船只。}同样，一个在人生事件和考验上有经验的人，当他看到风暴逼近，恶神想要制造船难时，他会把剩余的财富也扔掉。你被抢劫了吗？施舍，你就减轻了船的负担。强盗掠夺了你吗？把剩下的给基督。这样，你就能从先前的损失中安慰你的贫穷。减轻船的负担，不要紧紧抓住剩下的，以免船被水灌满。他们为了保全身体，把货物扔到海里，不等袭来的波浪颠覆船只。而你难道不愿停止船难，好拯救灵魂吗？\par

请试一下，我恳求你——你若不信，试一下，你就会看到天主的荣耀。当任何痛苦的事情发生时，立即施舍；感谢这事的发生，你就会看到多么大的喜乐临到你。因为属灵的收益，即使很小，也大到足以使一切身体的损失黯然失色。只要你还有东西可以给基督，你就是富有的。告诉我，如果你被抢劫时，国王来到你面前，伸出手来，求你给他一点东西，你难道不会认为自己比所有人都富有吗？因为国王甚至在如此大的贫穷之后也没有嫌弃你？不要被财富冲昏头脑，只要胜过你自己，你就会胜过魔鬼的攻击。获得巨大的收益是在你能力之内的。\par

让我们轻视财富，这样我们才不会轻视灵魂。但人怎能轻视财富呢？你难道没有看到，在美丽的身体和爱慕它们的人身上，只要他们在眼前，火就被点燃，火焰明亮升起；但当有人把他们移得远远的，一切就熄灭，一切就沉睡；同样，在财富方面，也不要让人预备金子、宝石、项链；这些东西一被看见，就缠住眼睛。但如果你想像古人那样富有，不要富在金子，而要富在必需品上，好将你所拥有的分施给别人。不要喜爱装饰。这样的财富既容易被强盗觊觎，又会给我们带来烦恼。不要有金银器皿，而要有储备的食粮、酒和油，这些不是要再卖出去换钱，而是要供应给需要的人。如果我们从这些多余之物中抽身，我们就能获得天上的美物；愿天主恩赐我们都获得这些美物，在我们的主耶稣基督内，偕同祂，等等。\par

\SourceRecord{Translated by John A. Broadus.；Nicene and Post-Nicene Fathers, First Series；Vol. 13.；Edited by Philip Schaff.；Buffalo, NY: Christian Literature Publishing Co.,；1889.；<http://www.newadvent.org/fathers/230403.htm>.}

% structure: p0001 source=h1 target=section reason=page-title

\section{得撒洛尼人前书第四篇讲道}

\BibleRef{得前 3:5-8}\par

\Quote{为此，我既不能再忍，就打发人去，要知道你们的信德，恐怕那诱惑者诱惑了你们，我们的劳苦就归于徒然。但当弟茂德刚从你们那里来到我们这里，把你们信德和爱德的好消息，以及你们常常记念我们，切慕见我们，如同我们切慕见你们一样的事传报给我们；为此，弟兄们，我们在一切困苦和患难中，因着你们的信德得了安慰：因为你们若在主内站得稳，我们就活了。}\par

一个备受争议的问题摆在我们面前，它从许多来源汇集而成。但这是什么问题呢？他说：\Quote{为此，我不能再忍，就打发\Person{弟茂德}去，要知道你们的信德。}你怎么说？那位知道许多事、听到不可言传的话、甚至上到第三层天的人，难道身在\Place{雅典}时就不知道吗？况且距离不远，他与他们分别也不久。因为他说：\Quote{暂时与你们分离。}他不知道\Place{得撒洛尼人}的情况，却不得不打发\Person{弟茂德}去知道他们的信德，他说：\Quote{恐怕那诱惑者诱惑了你们，我们的劳苦就归于徒然。}\par

那么该说什么呢？就是圣人们并不知道一切事。这可以从许多例子中得知，不论是早期的圣人，还是之后的圣人。例如，\Person{厄里叟}不知道有关那妇人的事（见：\BibleRef{列下 4:27}）；\Person{厄里亚}对天主说：\Quote{只剩下我一个人，他们还要寻索我的性命。}于是，天主对他说：\Quote{我为自己留下七千人。}（见：\BibleRef{列上 19:10}）又如\Person{撒慕尔}被派遣去给达味傅油时，上主对他说：\Quote{不要看他的容貌，也不要看他身材高大，因为我已弃绝了他；因为上主不像人看人：人看外貌，上主却看人心。}（见：\BibleRef{撒上 16:7}）\par

这出于天主极大的眷顾。如何并怎样呢？既是为了圣人自己，也是为了那些信靠他们的人。因为正如祂允许有迫害一样，祂也允许他们不知道许多事，好使他们保持谦卑。为此，\Person{保禄}也说：\Quote{有一根刺加在我身上，就是撒殚的使者来攻击我，免得我过于高举自己。}（见：\BibleRef{格后 12:7}）又为免得别人也对他们有过高的想法。因为如果他们因异事而以为他们是神，那么如果他们一直知道一切事，更是如此。他接着又说：\Quote{恐怕有人把我看得超过他所看见我的，或从我所听见的。}（见：\BibleRef{格后 12:6}）再听\Person{伯多禄}医治瘸子时说的话：\Quote{为什么你们定睛看我们，好像因我们自己的能力和虔诚使这人行走呢？}（见：\BibleRef{宗 3:12}）如果他们说了这些，行了这些，并借着这些微小有限的奇迹，尚且生出这样的恶念，那么若是大的奇迹，就更不用说了。\par

允许这些事还有另一个原因。即没有人能说他们是因为非同常人而行那些卓越的事，以致人人都变得懈怠，祂便显示他们的软弱，以他们的愚拙断开一切无耻的借口。为此他有所不知，为此他也常常决定了却不来，好让他们知道有很多事他不知道。由此产生了大益处。因为若已有人声称\Quote{这人就是那称为大能的天主的能力}（见：\BibleRef{宗 8:10}），有人说是这人或那人；倘若不是这样，他们会想什么呢？\par

但这里似乎是对他们的责难。但完全相反，这正显示出他们卓越的品行，并证明了他们试探的过度。如何呢？请留意。因为如果你先说\Quote{我们被指定为此}，且\Quote{不要有人动摇}，为什么你又打发\Person{弟茂德}，担心发生你不愿的事呢？这确实出于他极大的爱。因为那些爱的人，出于过分的热情，甚至对安全之事也疑惧。但这由他们巨大的试探所引起。因为我确实说了我们被指定为此，但试探的过度使我惊慌。所以他不是说，我打发他作为责难你们，而是\Quote{我不能再忍}，这更是爱的表达。\par

\Quote{恐怕那诱惑者诱惑了你们}是什么意思？你看见在患难中动摇是出于魔鬼，出于他的诱惑吗？因为当他不能动摇我们自己时，他便采取另一种方式，通过我们来动摇较软弱的人，这正是极其软弱且无可推诿的表现；正如他在\Person{约伯}的事上所行的，激动他的妻子说：\Quote{你说一句话反对天主，然后死吧！}（见：\BibleRef{约 2:9}）你看他如何诱惑了她。\par

但他为什么不说\Quote{动摇}，却说\Quote{诱惑}呢？因为，他说，我只怀疑到如此程度，即你们曾被诱惑。因为他称他的攻击为诱惑。谁接纳他的攻击，谁就动摇了。奇怪！\Person{保禄}的情感何其伟大！他不看自己的患难，也不看对他所设的计谋。因为我想他那时仍留在那里，正如\Person{路加}所说：\Quote{他在\Place{希腊}住了三个月，犹太人就设计害他。}（见：\BibleRef{宗 20:3}）\par

因此他所关心的不是自己的危险，而是门徒们。你看见他如何超越所有生身父母吗？因为我们在自己的患难和危险中，常失去对一切的记忆。但他为自己的孩子们如此恐惧战栗，竟打发\Person{弟茂德}到他们那里去——独有\Person{弟茂德}是他的安慰，是他的同伴和同劳者——而且是在危险当中打发他。\par

他说：\Quote{我们的劳苦就归于徒然了。}为什么呢？因为即使他们偏离了，也不是由于你的过失，不是由于你的疏忽。但尽管如此，我以为，出于我对弟兄们极大的爱，我的劳苦就归于徒然了。\par

\BibleQuote{免得诱惑者诱惑了你们}。但他诱惑，不知道是否能推翻。那么，他即使不知道，仍攻击我们，而我们明知将完全战胜他，却不警醒吗？但他攻击我们，虽然不知道，却在\Person{约伯}的事上显明了。因为那恶魔对天主说：\BibleQuote{祢岂不是四面围护他和他的家，并他一切所有的吗？且伸手毁他一切所有的，他必当面诅咒祢。}\BibleRef{约 1:10}他试探：若见什么软弱，就发动攻击；若强固，就罢手。\BibleQuote{我们的劳苦，}他说，\BibleQuote{归于徒然。}让我们都听一听，\Person{保禄}如何劳苦。他没有说工作，而是说\Quote{劳苦}；他没有说你们失落了，而是说\Quote{我们的劳苦}。因此，即使有什么事发生，那也是有一定缘由的。但竟没有发生，这是一个伟大的奇事。他说，这些事我们确实预料到了，但结果却相反。因为我们不仅没有从你们那里增添痛苦，反而得到了安慰。\par

\BibleQuote{但是弟茂德刚刚从你们那里来到我们这里，把你们的信德和爱德的好消息报给了我们} \BibleQuote{给我们带来了喜讯，}他说。你看见\Person{保禄}的极大喜乐吗？他没有说给我们带来了消息，而是说\Quote{给我们带来了喜讯}。他认为他们的坚定和爱德是何等的美事。因为一个人若坚定，另一方也必须坚定。他因他们的爱德而喜乐，因为那是他们信德的标记。\BibleQuote{并且你们对我们常怀美好的纪念，渴望见到我们，如同我们渴望见到你们一样，}他说。那就是，带着赞美。不是我们在你们中间时，也不是我们行奇迹时，而是现在，我们远离你们，受鞭打，遭受无数痛苦时，\Quote{你们对我们怀有美好的纪念}。请听，门徒因对老师怀有美好纪念而受到称赞，他们如何被称为有福的。让我们效法他们。因为我们是使自己受益，而不是那些被我们所爱的人。\BibleQuote{渴望见到我们，}他说，\BibleQuote{如同我们也渴望见到你们。}这同样安慰了他们；因为对于一个爱人的人来说，感受到所爱的人知道自己被爱，是一种极大的安慰和慰藉。\par

\BibleQuote{为此，弟兄们，我们在一切困苦和患难中，因你们的信德，在你们身上得了安慰。因为你们若在主内站稳，我们就活了。}有什么能与\Person{保禄}相比？他以为近人的救恩是自己的，如此对待众人，如同对待肢体一样。现在谁能说出这样的话来？或者，谁能有这样的想法？他没有要求他们因他为他们所受的试炼而感谢他，反而因他们没有因他的试炼而动摇而感谢他们。就好像他说，那些试炼伤害你们甚于我们；你们受诱惑甚于我们，你们一无所受，甚于我们这些受苦的。因为，他说，\Person{弟茂德}把这些好消息带给我们，我们就感觉不到任何痛苦，反而在一切患难中得了安慰；不单是在这患难中。因为一位好老师，只要他的学生们的事情合他的心意，就没有别的事能触动他。他说，因着你们，我们得了安慰；你们坚定了我们。然而情形却相反。因为他们在受苦时不屈服，而是勇敢地站立，这足以坚定门徒。但他却将整个事情反转，将赞美归给他们。他说，你们膏抹了我们，你们使我们重新呼吸；你们没有让我们感受到我们的试炼。他没有说我们重新呼吸，也没有说我们得了安慰，而是说什么？\BibleQuote{我们现在活着，}表明他认为除了他们的绊倒，没有什么是试炼或死亡，而他们的进步就是生命。还有谁能够如此表达对自己门徒软弱的悲伤或喜乐？他没有说我们喜乐，而是说\Quote{我们活着}，指未来的生命。\par

所以，没有这个，我们甚至不认为活着就是生命。老师应当如此感受，门徒也当如此；这样，任何时候都不会有差错。然后，他进一步缓和语气，请看他说什么：\par

第9、10节。\BibleQuote{我们为你们，在我们的天主面前，因你们所享有的一切喜乐，该怎样感谢天主呢？我们恳切地日夜祈求，为能见到你们的面，并能补足你们信德所欠缺的。}\par

他说，你们不只是我们生命的缘由，也是极大喜乐的缘由，甚至我们不能恰当地感谢天主。他说，你们的好行为，我们认为是天主的恩赐。你们向我们显示了如此多的恩惠，以致我们认为那是出于天主；是的，甚至它就是出于天主。因为这样的心态并非来自人的灵魂或谨慎。\par

\BibleQuote{日夜，}他说，\BibleQuote{恳切地祈求。}这也是一种喜乐的标记。因为正如农夫，听到自己耕种的土地上谷穗累累，便渴望亲眼见到如此美景，同样，\Person{保禄}渴望见到\Place{马其顿}。\BibleQuote{恳切地祈求。}请注意这种极度的祈求；\BibleQuote{为能见到你们的面，并能补足你们信德所欠缺的。}\par

这里有一个大问题。因为如果现在你们活着，是因为他们站得稳，\Person{弟茂德}带来了你们\BibleQuote{信德和爱德的喜讯}，你们充满了如此大的喜乐，以致不能相称地感谢天主，那么你们怎么说他们的信德有缺欠呢？那些是谄媚的话吗？绝不，远非如此。因为先前他作证说，他们忍受了许多斗争，并且没有比\Place{犹太}各教会更受影响。那么是什么呢？他们没有完全享受他的教导的益处，也没有学到一切他们应当学的东西。他在结尾处表明了这一点。也许在他们中间有关于复活的争议，有许多人扰乱他们，不是通过诱惑，也不是通过危险，而是扮演教师的角色。这就是他所说的他们信德中缺少的东西，因此他这样解释了自己，并没有说\Quote{你们应当被坚固}，那里他确实担心信德本身：\Quote{我派了}他说，\Quote{\Person{弟茂德}\InnerQuote{坚固你们}}，但这里说，\Quote{补充那缺少的}，这更是教导的事，而不是坚固的事。正如他另处所说：\Quote{使你们成全于一切善行。}\BibleRef{格前 1:10} 成全的事就是其中有一些小缺欠的事：因为那是被带向成全的事。\par

\BibleQuote{愿我们的天主和父自己，以及我们的主耶稣基督，亲自给我们开路，使我们到你们那里去。愿主使你们彼此之间的爱德，并对众人的爱德，增加且丰盈起来，就像我们对你们所有的一样。}\par

这是过分爱德的证明，因为他不仅亲自为他们祈祷，甚至在他的书信中插入自己的祈祷。这证明了一个炽热的灵魂，一个真正无法抑制的灵魂。这也是对那里也做的祈祷的证明，同时也是借口，表明他们去不了不是自愿的，也不是由于懒惰。好像他曾说：愿天主亲自剪除到处分散我们的诱惑，使我们能直接到你们那里去。\Quote{愿主使你们增加且丰盈。}你看到他话语中表现出的无法抑制的爱德疯狂吗？使你们增加且丰盈，代替使你们成长。好像有人说，他渴望以某种超丰富程度被他们所爱。\Quote{就像我们对你们所有的一样，}他说。我们的部分已经完成，我们祈祷你们的得以完成。你看到他希望爱德不仅对彼此，而且对所有人延伸吗？因为这才是虔诚爱德的真正本质，它拥抱所有人。如果你确实爱这个人，但不爱那个人，那是人的爱。但我们的不是这样。\Quote{就像我们对你们所有的一样。}\par

\BibleQuote{为使你们的心在圣德上无可指摘，在我们的天主和父面前，当我们的主耶稣基督与众圣人降临时。} \BibleRef{得前 3:13}\par

他指出爱德使他们自己得益，而不是被爱的人得益。他说，我希望这爱德满溢，使你们无可挑剔。他没有说\Quote{坚固你们}，而是\Quote{你们的心}。\BibleQuote{因为从心里发出来的是恶念。}\BibleRef{玛 15:19} 因为可能不做坏事，人却是恶的；例如，有嫉妒、不信、诡诈、幸灾乐祸、缺乏爱心、持守扭曲的教义，这些都是出于心；对这些纯洁就是圣洁。的确，贞洁惯常被特别称为圣洁，因为淫乱和奸淫也是不洁。但普遍地，所有罪都是不洁，所有德行都是纯洁。因为，经上说：\BibleQuote{心里洁净的人是有福的。}\BibleRef{玛 5:8} 祂所说的\Quote{洁净的人}指的是各方面洁净的人。\par

因为其他事物也知道玷污灵魂，并不逊色。因为邪恶玷污灵魂，请听先知的话：\BibleQuote{耶路撒冷啊，当从心里洗去罪恶。}\BibleRef{耶 4:14} 又说：\BibleQuote{你们要洗濯，自洁，从你们心里除掉邪恶。}\BibleRef{依 1:16} 他没有说\Quote{淫乱}，所以不仅淫乱，其他事物也玷污灵魂。\par

\Quote{使你们的心得以坚固，}他说，\Quote{无可指摘，在圣德上，在我们的天主和父面前，在我们的主耶稣基督与众圣人降临时。}因此，基督那时将是审判者，但不仅在他面前（单独），而且我们也将在父面前受审。或者他意味着这个：在天主面前无可指摘，正如他常说的\Quote{在天主眼中}，因为这是真正的美德——不是在人的眼中？\par

是爱德使他们无可指摘。因为它确实使人真正无可指摘。有一次当我与某人谈论此事，说爱德使人无可指摘，并且爱邻人的爱德不容许罪恶有任何进路，在我谈话中又全面追溯其余一切时——我的一位熟人插进来说：\Quote{那么淫乱呢？难道不可能既爱又犯淫乱吗？而且这确实是从爱德产生的。吝啬、奸淫、嫉妒、敌对阴谋以及诸如此类的一切，都可以通过爱邻人的爱德而停止；但淫乱如何呢？}他说。于是我告诉他：\Quote{连这个爱德也能停止。因为如果一个人爱一个犯淫乱的女人，他将会竭力既把她从其他男人那里拉过来，也不再增加她的罪。所以与一个女人犯淫乱是极度恨与那女人犯淫乱者的行为，但一个真正爱她的人会将她从那可憎的习俗中拉出来。}而且没有一种罪，爱德的力量像火一样不能烧毁的。因为让一把坏柴抵抗一堆大火，比让罪的本性抵抗爱德的力量更容易。\par

因此，让我们把这道理种植在我们自己的灵魂里，好使我们能与诸圣一同站立。因为诸圣都是因着对近人的爱而中悦了天主。亚伯尔为何被杀害，而他自己却没有杀人？是因他对兄弟的挚爱，他甚至不容这样的念头产生。加音为何接受了这毁灭性的嫉妒瘟疫？因为我不再称他为亚伯尔的兄弟了！因为爱的根基在他内没有稳固确立。诺厄的儿子们为何得到赞誉？岂不是因他们挚爱父亲，不忍见他的赤身？而另一个为何受诅咒？岂不是因他不爱父亲？亚巴郎为何得到赞誉？岂不是因他在为侄子所行的事上，以及为索多玛人恳求时所表现的爱？因为诸圣深受爱与同情的影响，强烈，强烈地。\par

请想一想，我恳求你们。保禄，那敢于面对烈火的人，坚硬如钻石，坚定而不可动摇，四面坚固，被对天主的敬畏所钉牢，且不可弯曲；因为\Quote{谁（他说）能使我们与基督的爱隔绝呢？是困苦、窘迫、迫害、饥饿、赤身露体、危险、刀剑吗？}\BibleRef{罗 8:35}这位敢于面对这一切、天地的一切、嘲笑死亡之钻石般大门、无物可阻挡的人——当他看见他所爱之人的眼泪时，竟如此破碎和压伤——那钻石般的人——甚至不隐藏他的情感，而是立刻说：\Quote{你们做什么，哭泣而使我的心破碎？}\BibleRef{宗 21:13}你说什么，告诉我？一滴眼泪竟有力量压伤那钻石般的灵魂吗？是的，他说，因为我抵挡一切，除了爱。这爱胜过我，征服我。这是天主的性情。深渊的水没有压伤他，而几滴眼泪却压伤了他。\Quote{你们做什么，哭泣而压碎我的心？}因为爱的力量是伟大的。你们不是又看见他哭泣吗？你为什么哭泣？告诉我。\Quote{三年之久，我昼夜不停地流泪劝诫每一个人。}\BibleRef{宗 20:31}出于他伟大的爱，他害怕，恐怕某种瘟疫会混入他们中间。又如：\Quote{我因许多的忧苦和心中的痛苦，流了许多眼泪，写信给你们。}\BibleRef{格后 2:4}\par

再说若瑟，告诉我，那坚定的人，他抵挡了如此巨大的暴虐，在面对如此巨大的情欲火焰时显出如此高贵，如此战胜并克服了他女主人的狂妄。当时有什么不能诱惑他呢？美丽的外表，尊贵的骄傲，昂贵的衣裳，香水的芬芳（因为这一切都懂得如何软化灵魂），比一切更柔和的言语！因为你知道，那爱得如此炽烈的女人，没有什么卑贱的话她不愿说出口，摆出恳求者的姿态。这个妇人，尽管身穿金饰，具备王后的尊严，竟如此破碎，以致跪在可能是个俘虏少年面前，甚至也许哭泣着抱住他的膝盖，不止一次，而是屡次如此。然后他可能看见她的眼睛闪耀着最明亮的光辉。因为很可能她不是简单地，而是极其精心地装饰她的美貌，希望用许多网罗捕捉基督的羔羊。这里还加上许多魔法咒语。然而，这个坚定、坚固、如岩石般坚硬的人，当他看见那些曾出卖他、把他扔进坑里、卖掉他、甚至想杀害他、是他监牢和荣耀之因的兄弟们，听见他们如何欺骗父亲（因为我们说，经上记载说，有一个被野兽撕裂了\EditorialNote{创 37:20}，而\EditorialNote{64:28}），他破碎了，软化了，压伤了。\Quote{他就哭了，}经上说，而且不能忍受自己的情感，便进去镇定自己\BibleRef{创 43:30}，即擦去眼泪。\par

这是怎么回事？你哭泣吗，若瑟？然而目前的情况不值得哭泣，而值得愤怒、忿怒、愤慨及巨大的报复和惩罚。你的仇敌——那些杀兄弟者——在你手中；你可以满足你的愤怒。而且这也不是不义。因为你并不是首先行不义，而是自卫反击那些行了不义的人。不要看你的尊严。这并非出于他们的谋划，而是出于天主，祂将祂的恩宠倾注于你。你为什么哭泣？但他会说：愿我藉这冤屈的记忆而倾覆我在一切事上所获得的赞誉，这种事决不会发生。这确实是哭泣的时候。我不比野兽更愚钝。它们向本性奠祭，无论自己遭受什么伤害。我哭泣，他说，是因为他们曾这样待我。\par

让我们也效法这人。让我们为那些伤害我们的人哀悼哭泣。不要向他们发怒。因为他们确实值得眼泪，为那使他们自己招致的惩罚和定罪。我知道，你们现在如何哭泣，如何欢喜，既钦佩保禄，又惊讶于若瑟，并称他们为有福的。但如果谁有仇敌，现在让他回想起来，让他记在心里，趁着心因对诸圣的回忆而仍火热，他才能消解愤怒的顽固，并软化那刚硬和麻木的。我知道，你们离开这里之后，在我停止说话之后，若有任何温暖和热忱存留，也不会像现在你们听我说话时那样大。所以，如果有谁变得冷淡了，让他融化那冰霜。因为对伤害的记忆实在是霜和冰。但让我们呼求正义的太阳，让我们恳求祂将祂的光辉照在我们身上，就再不会有厚冰，而是可饮的水。\par

如果义德的太阳之火触及了我们的灵魂，它就不会留下任何冰冻、任何坚硬、任何灼烧、任何不结果实的东西。它会使一切成熟、一切甘甜、一切充满极大的喜乐。如果我们彼此相爱，那光芒也必会到来。请允许我，我恳求你们，恳切地说这些事。让我听到，藉着这些话我们产生了一些效果：有人已去用双臂抱住他的仇人，拥抱他，缠绕着他，深情地亲吻他，并且哭了。即使对方是野兽、石头或无论什么，也会被这样深情的仁慈变成温柔。他凭什么是你的敌人？他侮辱了你吗？但他根本没有伤害你。难道你为了钱财，竟让你的兄弟与你为敌吗？不要这样做，我恳求你们。让我们消除一切。这是我们的时节。让我们善用它。让我们砍断我们罪恶的绳索。在我们前去受审之前，让我们不要彼此论断。\BibleQuote{不要让太阳}（经上说）\BibleQuote{在你们的忿怒上落下} \BibleRef{弗 4:26} 不要让任何人推迟。这些推迟产生拖延。如果你今天推迟了，你更加脸红；如果再加上明天，羞耻更大；如果第三天，更糟。那么，让我们不要使自己蒙羞，而是让我们宽恕，好使我们被宽恕。如果我们被宽恕，我们将获得一切祝福，借着我们的主\Person{耶稣基督}，偕同祂，等等。\par

\SourceRecord{Translated by John A. Broadus.；Nicene and Post-Nicene Fathers, First Series；Vol. 13.；Edited by Philip Schaff.；Buffalo, NY: Christian Literature Publishing Co.,；1889.；<http://www.newadvent.org/fathers/230404.htm>.}

% structure: p0001 source=h1 target=section reason=page-title

\section{论\BibleBook{得撒洛尼}{前书}第五篇讲道}

\BibleRef{得前 4:1}\par

\BibleQuote{最后，弟兄们，我们在主耶稣内恳求和劝勉你们，既然你们从我领受了应当如何行走及悦乐天主，就该更加充盈不已。因为你们知道我们曾藉着主耶稣基督给了你们什么命令。原来天主的旨意就是你们的圣化。}\par

当他处理了紧迫的、手头的事，即将着手进入那些永久的、他们应持续聆听的事时，他加上了这句话：\BibleQuote{最后}，即是常远的、永远的。\BibleQuote{我们在主内恳求和劝勉你们。}奇怪！他甚至不把自己当作足够可信的劝勉者。然而谁又能如此可信呢？但他把基督与自己联合。他说：我们劝勉你们，藉着天主。这也是他对格林多人说的：\BibleQuote{天主藉着我们恳求（劝勉）你们。}\BibleRef{格后 5:20}\BibleQuote{就如你们从我们领受了。}这\Quote{领受}不仅关乎言语，也关乎行动，即是\BibleQuote{你们应当怎样行走}，他藉此指整个生活方式。\BibleQuote{并悦乐天主，使你们更加充盈不已。}这就是说，藉着更加充盈，你们不要止于诫命的界限，甚至要超越它们。因为这就是\BibleQuote{你们更加充盈不已}之所在。在前文中，他接纳了他们坚定信德的奇妙，但在此处他规范他们的生活。因为这就是长进，甚至超越诫命和规条。因为这一切不再出于老师的强迫，而是出于他们自己的自愿选择。正如大地不应只承载播撒之物，同样，灵魂也不应止于被灌输的事物，而应超越它们。你们看见他恰当地说\Quote{超越}了吗？因为美德分为这两件事：远离邪恶和行善。因为远离邪恶并不足以到达美德，它是一条路，是一个引向它的开端；我们仍需要极大的勤勉。因此，他按照诫命的次序告诉他们应避免的事。这是合理的。因为做这些事的确招致惩罚，但不做却不会带来赞美。然而，美德的行为，例如施舍财物等，不属于诫命的范畴，他说。但什么是值得的呢？\BibleQuote{能够领受的，就让他领受。}\BibleRef{玛 19:12}因此，有益的是，正如他怀着极大的恐惧和战栗将这些诫命给了他们，他也藉着这些书信提醒他们他的关怀。因此他不是重复它们，而是提醒他们这些。\par

\BibleQuote{因为}他说，\BibleQuote{这是天主的旨意，就是你们的圣化，使你们远离一切邪淫。}此外，注意，他没有任何其他地方如此强烈地提及任何其他事，如同这一\OriginalTerm{圣化}{sanctification}。在其他地方他也这样写道：\BibleQuote{你们要追求与众人和平，并追求圣化，没有圣化，谁也见不到主。}\BibleRef{希 12:14}你们为什么惊奇，如果他处处给门徒写这个主题，甚至在他的\BibleBook{弟茂德}{前书}中也说：\BibleQuote{你当保持纯洁。}\BibleRef{弟前 5:22}此外，在他的\BibleBook{格林多}{后书}中他说：\BibleQuote{以极大的忍耐、斋戒、藉着纯洁。}\BibleRef{格后 6:5}人们可以在许多地方找到这一点，既在这封\BibleBook{罗马}{书}中，也在各处，在他所有的书信中。因为实际上，这是一种对所有有害的邪恶。正如一头满身泥污的猪，无论它进入哪里，都以其臭气充满各处，并以其粪秽窒息感官，\OriginalTerm{邪淫}{fornication}也是如此；它是一种不易洗去的邪恶。但当有些人甚至已有妻子却还这样做时，那是多么过分的不义！\BibleQuote{因为}他说，\BibleQuote{这是天主的旨意，就是你们的圣化，使你们远离一切邪淫。}因为有许多种放荡行为。放荡的快乐有各种花样，是不堪提及的。但他说了\BibleQuote{远离一切邪淫}后，留给那些知道它们的人。\par

第四节至第五节：\BibleQuote{你们各人应知道怎样在圣化和尊贵中持守自己的器皿，不在情欲的激情中，正如那不认识天主的外邦人。}\par

他说：\BibleQuote{你们各人应知道怎样在圣化和尊贵中持守自己的器皿。}因此，这是一件必须学习的事，且要殷勤，不可放纵。但我们的\OriginalTerm{器皿}{vessel}，当它纯洁时，我们拥有它；当它不洁时，罪恶拥有它。这是合理的。因为它不做我们所愿望的事，而是做罪恶所命令的事。\BibleQuote{不在情欲的激情中，}他说。\BibleQuote{正如外邦人，}他说，\BibleQuote{那不认识天主的。}因为这种人就是不期待自己将受惩罚的人。\par

\BibleRef{得前 4:6} \BibleQuote{没有人在这事上越轨，亏负他的弟兄。}\par

他说得很好：\Quote{谁也不可越分}。天主为每人指定了妻子，为天性划定了界限，即只与一人结合；因此与另一人结合便是越分，是占有不属于自己的东西，是抢夺；甚至比任何抢夺更残忍；因为我们的财富被夺走时，我们固然悲伤，但婚姻被侵犯时，悲伤更甚。你称他为弟兄，却亏负他，且是在不法之事上亏负他？此处他论及奸淫，但上文也论及\Quote{一切邪淫}。因为当他即将说\Quote{谁也不可越分、亏负他的弟兄}时，他说：不要以为我只在弟兄中间这样说；你们根本不可占有别人的妻子，甚至那些没有丈夫、属于公共的女子也不可。你们必须戒绝\Quote{一切邪淫}；\Quote{因为}，他说，\Quote{主在这一切事上必要报应}。他先劝诫他们，又羞辱他们，说：\Quote{如同外邦人一样}。然后从推理中表明亏负弟兄的不当。此后他加上主要的话：\Quote{因为}，他说，\Quote{主在这一切事上必要报应，正如我们预先警告并见证过的}。因为我们做这些事必不不受惩罚，我们享受的快乐也不及所受的惩罚。\par

第7节。\Quote{因为天主召叫我们，不是为污秽，而是为圣化。} \BibleRef{得前 4:7}\par

因为他曾提到\Quote{他的弟兄}，又加上了天主是报应者，表明即使受亏负的是不信者，做这事的人也要受惩罚；他说：祂惩罚你并非因为替那人报应，而是因为你侮辱了祂自己。祂亲自召叫了你，你却侮辱了召叫你的那一位。为此，他加上了：\par

第8节。\Quote{因此，那拒绝的，不是拒绝人，而是拒绝天主，因为祂将祂的圣神赐给你们。} \BibleRef{得前 4:8}\par

所以，即使你玷污了皇后，他说，或甚至你本人的婢女，而她有丈夫，所犯的罪相同。为什么？因为祂所报应的不是受伤害的人，而是祂自己。因为你的玷污同样，你同样侮辱了天主；因为两者都是奸淫，两者都是婚姻。即使你没有犯奸淫，而是行了邪淫，尽管娼妓没有丈夫，但天主仍然报应，因为祂报应的是祂自己。因为你做这事，与其说是藐视人，不如说是藐视天主。这从以下事实显而易见：你行这事时对人隐瞒，却装作天主看不见你。因为请告诉我，如果有人堪当紫袍和国王（皇帝）的无限尊荣，被命令按照这尊荣生活，却去与任何女人玷污自己；他侮辱了谁？是她，还是赐给他一切的国王？她的确也受了侮辱，但不相等。\par

因此，我恳求你们，让我们谨防这罪。因为正如我们惩罚女人，当她们已嫁给我们却把自己给了别人时，我们也同样受惩罚，虽然不是按罗马的法律，而是按天主的法律。因为这也是奸淫。因为奸淫不仅发生在与有夫之妇行这事的人身上，也发生在与妻子相连的男人身上。请留心听我说。因为所说的话虽然触犯许多人，却有必要说出来，以在将来纠正此事。不仅当我们玷污一个已嫁与男人的女人时是奸淫；而且当我们自己已娶了妻子，却玷污一个自由无夫的女子时，这事也是奸淫。因为如果那与之行淫的女子没有约束，那算什么？你却是受约束的。你犯了法。你伤害了自己的肉身。请告诉我：如果你的妻子与一个没有妻子的自由人行奸淫，你为什么惩罚她？因为那是奸淫。为什么？然而玷污她的那个人没有妻子，但她却是与丈夫结合的。那么，你也是与妻子结合的，所以同样，你的罪行也是奸淫。因为经上说：\Quote{凡休妻的，除非为了奸淫的缘故，就是叫她犯奸淫；凡娶那被休的妇人的，也是犯奸淫。} \BibleRef{玛 5:32} 如果娶离婚妇人是犯奸淫，那么有自己的妻子，又玷污另一个女人——这对每个人都是显而易见的。但或许对于你们男人，关于这题目已经说得够了。因为关于这样的人，基督说：\Quote{他们的虫子不死，火也不灭。} \BibleRef{谷 9:44} 然而为了青年人的缘故，有必要对你们说话，与其说对青年人，不如说对你们。因为这些话不仅适合他们，也适合你们。怎么讲？我现在就告诉你们。那没有学会行奸淫的人，也不会知道如何犯奸淫。但那在娼妓中打滚的人，很快就会陷入另一件事，并且会玷污自己，即使不与有夫之妇，也会与无夫的女子。\par

那么我劝什么，好铲除这些根呢？你们中间有年幼的儿子，正将他们抚育成世俗之人的，要赶快使他们戴上婚姻的轭。因为当他还年轻时，情欲搅扰他，在婚前这段时间，要用训诫、威胁、恐惧、应许和无数其他方法约束他们。但在结婚的时候，不要有人推迟。看，我说的是媒人的话，要你们让儿子娶妻。但我这样说并不羞愧，因为就连保禄也不羞愧地说：\Quote{你们不要彼此亏负} \BibleRef{格前 7:5}，这似乎比我说的话更可耻，然而他并不羞愧。因为他不关注言辞，而关注被言辞纠正的行为。当你的儿子长大成人，在他从军或从事任何其他生活道路之前，要考虑他的婚事。如果他看见你很快就要为他娶妻，而且中间的时间短暂，他就能耐心忍受情欲的火。但如果他觉察到你懈怠迟缓，要等到他获得大笔收入才为他成婚，他因时间漫长而绝望，就会轻易陷入淫乱。但哀哉！这里万恶的根源也是贪财。因为没有人关心他的儿子应如何清醒和谦逊，反而所有人都为黄金疯狂，因此无人将此视为要事。因此我劝你们首先要好好管束他们的灵魂。如果他找到贞洁的新娘，并且只认识那一具身体，那么他的欲望会强烈，他对天主的敬畏也更大，婚姻真正可敬，接受纯洁无瑕的身体；后代将充满祝福，新娘和新郎会彼此顺服，因为双方都没有沾染别人的习气，他们会互相谦让。但有人在年轻时就开始放荡，经验娼妓之道，在第一、二个晚上会称赞自己的妻子；但之后他会很快回到那种放荡中，寻求那种放荡不羁的欢笑、充满卑劣意味的话语、放荡的行为，以及其他所有不能容忍提及的猥亵。但一个自由女子不会忍受做出这样的表露，也不会玷污自己。因为她被许配给丈夫，是要作他生活的伴侣，为生育子女，不是为了猥亵和欢笑；她应当管理家务，也教导他庄重，而不是给他提供淫乱的燃料。\par

但娼妓的姿态在你看来似乎令人愉悦。我知道。因为圣经说：\Quote{淫妇的嘴唇滴下蜂蜜} \BibleRef{箴 5:3}。正因如此，我费尽周折，使你们不致尝到那蜂蜜，因为它立刻化为苦胆。圣经又说：\Quote{她在片刻间使你的喉咙舒畅，但日后你必发现她比苦胆更苦，比双刃的剑更锋利} \BibleRef{箴 5:3}。你说什么？请容忍我说些粗俗的话——我若可以这样说——并像不知羞耻、毫不脸红的人那样表达。因为我不是自愿这样，而是因为那些行为无耻的人，我被迫说这类话。我们在圣经中甚至看到许多这样的事。因为连厄则克耳在责备耶路撒冷时，也说了许多这样的话，并不以为耻。这是公义的。因为他不是出于自己的倾向说这些话，而是出于关切。因为尽管这些话似乎不雅，但他的目的并非不雅，反而非常适合一个希望从灵魂中清除不洁的人。因为无耻的灵魂若不亲耳听到这些话，就不会受触动。就像一个医生想除去腐烂的疮口，先要把手指伸进伤口，若不先玷污他那治疗的双手，就不能治愈它。我也是这样。除非我先玷污我那医治你们情欲的口，否则我不能医治你们。但恰恰相反，我的口并未被玷污，如同他的手也未玷污一样。为什么？因为那不洁不是本性的，也不是来自我们自己的身体，如同那情况也不是来自他的手，而是来自别人的身体。但如果在一个别人的身体上，他不拒绝浸入自己的双手，告诉我，在我们的身体上，我们岂能拒绝呢？因为你们就是我们的身体，虽有病态和不洁，但终究是我们的。\par

那么，我说的这是什么，又为何这样长篇大论？你仆人穿的一件衣服，因它的污秽你绝不情愿穿，甚至宁愿赤身也不愿使用它。但一个不洁污秽的身体，不仅被你的仆人使用，还被无数其他人使用，你却要滥用它，而不厌恶吗？你听了这话感到羞耻吗？但应当为行为羞耻，而不是为言语。我略过其他所有事：粗鲁、品行的败坏、以及其余生活中的奴性和卑鄙。告诉我，你与你的仆人是否去找同一个女人？我但愿只是你的仆人，也许甚至是刽子手！然而你不能忍受握刽子手的手，却与那与他成为一体的人亲吻拥抱，并不战栗，也不惧怕！你不羞愧吗？你不赧颜吗？你不心如刀割吗？\par

我确实曾向你们的祖先说过，他们应当早早引导你们进入婚姻；然而你们也并非没有受罚的责任。因为若没有其他比你们更多的青年，无论是从前还是现在，度着贞洁的生活，那么你们或许还有借口。但若有人能如此，你们又怎能说我们无法抑制情欲的火焰呢？那些能够做到的人，因分享了同一本性而成为你们的控告者。请听保禄所说：\Quote{追求和平……并追求圣德，若无圣德，无人能见主。}\BibleRef{希 12:14}这警告难道还不足以吓倒你们吗？你们看见别人终身完全贞洁，庄重地度过一生，而你们竟不能在这青年时期约束自己吗？你们看见别人千万次克制了快乐，而你们连一次也不能自制吗？请允许我告诉你们原因：因为青年并非根源，否则所有青年都会放荡。但我们自己投身于火中。当你们前往剧场，凝视女人的裸体，以此满足眼目时，当下固然欢喜，但事后却因此燃起巨大的热病。当你们看见女人仿佛以身体姿势被展示，还有那些只包含放荡爱情的戏剧和歌曲——据说某女子爱某男子，未得成就而自缢；以及以母亲为对象的非法爱情——当你们也通过听觉、通过女人、通过形象，甚至通过老男人（许多人戴上假面具扮演女人）接受这些时，请告诉我，之后你们怎能保持贞洁？这些故事、这些景象、这些歌曲占据你们的灵魂，随后便产生这类梦境。因为灵魂的本性大多会在日间幻想它所希望和渴望的事物。因此，当你们在那里既看到卑劣的行为，又听到更卑劣的言语，受了伤却不施药，那伤口怎能不自然加重？疾病怎能不愈演愈烈，且比身体上的疾病更甚？因为如果我们愿意，我们的意志比身体更容易矫正。身体需要药物、医生和时间，而这里只需有意愿，便能成为善或恶。所以，你们是主动让紊乱进入的。因此，当我们为自己聚集有害之物，却对有益之物毫不关心时，又怎能恢复健康呢？\par

为此，保禄说：\Quote{正如不认识天主的外邦人。}让我们羞愧，让我们恐惧，如果不认识天主的外邦人常能贞洁，而我们却比他们更糟，让我们转念蒙羞。如果我们愿意，如果我们远离那些有害之事，就容易成就贞洁；因为如果我们不愿意，连避免奸淫也不容易。有什么比在市场上行走更容易呢？但由于过度的懒惰，它变得困难了——不仅对女人，有时甚至对男人也是如此。有什么比睡觉更容易呢？但我们甚至使这事也变难了。许多富人整夜辗转反侧，因为他们不等待睡意就来就寝。简而言之，人若愿意，便没有难事；若不愿意，便没有易事；因为我们是这一切的主人。为此，圣经说：\Quote{如果你们愿意，并听从我，}\BibleRef{依 1:19}又说：\Quote{如果你们不愿意，也不听从，}\BibleRef{依 1:20}因此一切都在于愿意或不愿意。为此，我们既受惩罚，也受赞扬。但愿我们成为那些受赞扬者，借恩宠和慈爱获得所应许的福分，等等。\par

\SourceRecord{Translated by John A. Broadus.；Nicene and Post-Nicene Fathers, First Series；Vol. 13.；Edited by Philip Schaff.；Buffalo, NY: Christian Literature Publishing Co.,；1889.；<http://www.newadvent.org/fathers/230405.htm>.}

% structure: p0001 source=h1 target=section reason=page-title

\section{论《得撒洛尼前书》第六篇讲道}

\BibleRef{得前 4:9-10}\par

\BibleQuote{论到弟兄之爱，我们无需写信给你们；因为你们自己由天主所教训，要彼此相爱；其实你们对所有在全部马其顿的弟兄们也正是这样行。}\par

那么，为何在恳切地论及洁德之后，即将论及工作的责任以及不为逝者忧伤之事时，他却插入一切善事之首——爱德，仿佛略过不谈似的，说：\BibleQuote{我们无需写信给你们}？这也是出于他的大智慧，属于灵性的训导。因为在此他显示两件事：第一，这事如此必要，以至于不需要教导。因为非常重要之事对众人是显而易见的。第二，这样说比直接劝戒更使他们羞愧。因为若他认为他们行为端正，因此不加劝戒，即使他们行为不正，也会更快地引导他们走向正途。请注意，他没有说对众人的爱，而是对弟兄的爱。\BibleQuote{我们无需写信给你们}。他本应沉默不语，若没有必要。但如今他藉着说没有必要，成就了一件比直接说出更大的事。\par

\BibleQuote{因为你们自己由天主所教训}\OriginalTerm{由天主所教训}{\GreekText{θεοδίδακτοι}}。且看，他用了何等高的赞颂，使天主成为在这件事上的导师。他说，你们无需从人学习。这也正如先知所说：\BibleQuote{他们都要受天主的教导。}\BibleRef{依 54:13} \BibleQuote{因为你们自己，}他说，\BibleQuote{由天主所教训，要彼此相爱。其实你们对众弟兄和全马其顿的弟兄们也正是这样行}；他也指其他所有的人。这些话大大鼓励他们这样做。我不但说你们由天主所教训，而且从你们所做的事上我知道这一点。在这方面他给了他们许多见证。\par

\BibleQuote{但我们劝勉你们，弟兄们，要更加充盈，并勉力；}意思是，增加并努力。\par

第11至12节。\BibleQuote{要安静，做自己的事，亲手劳作，正如我们所吩咐你们的：好使你们在外人面前行为端正，并且一无所缺。}\par

他显示懒惰是许多罪恶的根源，勤勉是许多益处的根源。他藉着发生在人间的事表明这一点，正如他常做的，且做得明智。因为这些事比超性的事更能引领多数人。再者，爱近人的标记不是从他们领受，而是施予他们。请注意。他正要劝戒和提醒，却将他们的善行放在中间，好使他们既能从先前的劝戒和警告中振作，那时他说：\BibleQuote{所以那弃绝的，不是弃绝人，乃是弃绝天主，}也不至于对此反感。这劳作的果效，就是人不从他人领受，也不游手好闲，而是藉劳动施予他人。因为经上说：\BibleQuote{施予比领受更为有福。}\BibleRef{宗 20:35} \BibleQuote{并要劳作，}他说，\BibleQuote{亲手。}那些追求超性劳作的人在哪里？你看，他如何除去每一个借口，说\BibleQuote{亲手}？难道有人用手禁食？或彻夜祈祷？或睡卧地上？这话无人能说。但他说的是超性的劳作。因为藉劳作施予他人，实为超性之事，没有什么可与之相比。\BibleQuote{好使你们行事为人，}他说，\BibleQuote{端正。}你看见他从何处触动他们了吗？他没有说免得你们因乞讨而羞愧。但他暗示了同样的事，却说得更为温和，既打击又不严厉。因为若我们中间的人为此而反感，更何况外人呢？他们找到无数的指责和把柄，当他们看见一个健康且能自立的人，却向他人乞讨求助。因此他们也称我们为\OriginalTerm{贩卖基督者}{Christ-mongers}。为此，他说，\BibleQuote{天主的名被亵渎了。}\BibleRef{罗 2:24} 但他没有陈述这些，而是说了那最能触动他们的事——事情的羞耻性。\par

第13节。\BibleQuote{弟兄们，我们不愿你们对长眠者无知，免得你们忧伤，像那些没有望德的人一样。}\BibleRef{得前 4:13}\par

这两件事，贫穷和沮丧，最困扰他们，也牵涉到所有的人。且看他是如何补救的。他们的贫穷来自他们的财产被夺走。但若他命令那些为基督的缘故被夺去财产的人，藉劳动维持生活，更何况其他人呢？那些财产被夺走，从他的话中明显可见：你们成了天主教会的一份子。如何成为一份子？\BibleQuote{你们欢喜地承受了财产的剥夺。}\BibleRef{希 10:34}\par

现在他继续开始有关复活的论述。为什么呢？难道他先前没有向他们论述过这点吗？是的，但他在这里提及了更深的奥秘。那么这是什么奥秘呢？他说：\BibleQuote{我们这些活着还存留到主来临的人，决不会在已安眠的人以先。}关于复活的论述足以安慰那忧伤的人。但现在所说的也足以使复活极其可信。但首先让我们谈谈前面的话：\BibleQuote{弟兄们，关于安眠的人，我们不愿意你们不知道，免得你们悲伤，像那些没有望德的人一样。}请看他也如何温和地对待他们。他没有说：\Quote{你们这样无知吗？}如同他对格林多人说的\BibleQuote{愚拙}？既然知道有复活，你们还这样悲伤，如同不信的人一样；但他说话非常温和，显示出对他们在其他德行上的尊重。他没有说\Quote{关于死了的人}，而是说\Quote{安眠的人}，一开始就暗示了安慰。他说：\BibleQuote{免得你们悲伤，像那些没有望德的人一样。}因此，为自己去世的人忧伤，就是像那些没有望德的人的行为。这很合理。因为一个对复活一无所知的灵魂，以为这个死亡就是终结，自然会忧伤、哀哭、痛哭，如同失去的人一样无法忍受。但你，既期待复活，为何还要哀叹呢？哀叹乃是没有望德之人的份。\par

请听这话，你们这些爱哭嚎的女人，你们这些在哀悼时忍受不了悲伤的女人，你们的行为如同外邦人。但如果为去世的人悲伤是外邦人的行为，那么告诉我，捶胸顿足、撕裂脸颊是谁的行为？如果你相信他会复活，他没有消亡，这只是打盹和睡眠，那你为什么哀叹呢？你说，是因为他的陪伴、他的保护、他对我们事务的关怀，以及他所有其他的服务。那么当你失去一个尚未能做什么的幼小孩子时，你为何哀叹？为何要试图召回他？你说，他显示了好希望，我期待他会成为我的支持者。为此我思念我的丈夫，为此我思念我的儿子，为此我哭嚎哀叹，并非不信复活，而是失去了支持，失去了我的保护者、我的伴侣、与我分享一切的人——我的安慰者。为此我哀悼。我知道他会复活，但我无法忍受中间的分离。许多烦恼涌向我。我暴露在所有愿意伤害我的人面前。那些从前怕我的仆人现在藐视我，践踏我。如果有人曾受过恩惠，他已经忘记从他那里得到的恩惠；如果有人曾被去世者亏待，为了报复他，他就对我大发怒气。这些事使我无法忍受我的寡居。正是为了这些事我折磨自己，正是为了这些事我哀哭。\par

那么我们将如何安慰这样的人？我们该说什么？如何驱散他们的悲伤？首先我要努力说服他们，他们的哭嚎并非由于他们所说的这些事，而是出于不合理的激情。因为如果你因这些事而哀悼，你就应该一直为去世的人哀悼。但如果一年过去后，你忘记了他，好像他从未存在过，你就不是在哀悼去世的人也不是他的保护。但你无法忍受分离，也无法忍受关系的断绝？那些甚至再次结婚的人又能说什么呢？确实！他们思念的是前夫。但让我们不要对她们说话，而是对那些对去世者保持慈爱之情的人说话。你为什么哀叹你的孩子？为什么哀叹你的丈夫？前者，你说，是因为我还没有享受过他；后者，是因为我期待能享受他更久。而正是这件事，显出何等缺乏信德，竟然认为你的丈夫或你的儿子是你的保障，而不是天主！你难道不想这会激怒祂吗？因为常常正因此祂将他们取去，使你不至于如此依恋他们，从而将你的希望从他们身上撤回。因为天主是忌邪的，祂愿意被我们爱超过一切；而且，因为祂极其爱我们。因为你们知道，这是那些爱得痴迷之人的习惯。他们过分嫉妒，宁愿舍弃自己的生命，也不愿在任何一位竞争者面前被轻视。也正因此，天主因着这些话语而将他取去。\par

请你告诉我，古时为何没有寡居和过早的孤苦？为何祂允许亚巴郎和依撒格长寿？因为亚巴郎即使活着也把天主置于他之前。祂确实对他说：杀；他就杀了他。为何祂让撒辣活到如此高龄？因为她活着时，他听从天主而不是她。为此，天主对他说：\Quote{你要听从你妻子撒辣的话。} \BibleRef{创 21:12}那时，没有人因爱丈夫或妻子，或因保护孩子而惹天主发怒。但如今，因为我们堕落，大不如前，我们男人爱妻子胜过爱天主，我们女人敬重丈夫胜过敬重天主。因此，祂甚至违背我们意愿地吸引我们归向祂。不要爱你的丈夫胜过爱天主，你就永远不会经历寡居。或者，即使寡居发生，你也不会感受到它。为什么？因为你有一位不朽的护佑者，祂更爱你。如果你爱天主更深，就不要哀伤：因为更被爱的那位是不朽的，祂不让你感到失去较少被爱的那位的痛苦。我将用一个例子向你显明。告诉我，如果你有一个丈夫，事事顺从你，受人尊敬，使你在各处都受尊重，不被轻视，众人都敬重他，聪明智慧，又爱你，你因他而被认为幸福，并且你与他一同生了一个孩子，孩子尚未成年就离世，那时你会感到痛苦吗？绝不会。因为更被爱的那位使痛苦消失。如今，如果你爱天主胜过你的丈夫，祂必定不会很快把他带走。即使祂带走他，你也不会感受到痛苦。为此，有福的约伯听到孩子们同时死亡的消息时，没有感到严重的痛苦，因为他爱天主胜过他们。既然他所爱的祂活着，那些事就不能使他痛苦。\par

你说什么，女人？你的丈夫或你的儿子是你的保护者？难道你的天主不保护你吗？谁给了你丈夫？不是祂吗？谁造了你？不是祂吗？祂确实从无中把你带出，向你嘘入灵魂，赋予你理智，恩赐你认识祂，并为了你没有怜惜自己的独生子，祂不保护你吗？而你的同伴保护你吗？这些话应得何等义怒！你从丈夫那里得到过这类东西吗？你什么也说不出。因为即使他善待过你，那也是在他先受了善待之后，而你已先开始了。但在天主方面，无人能这样说。因为天主施恩于我们，并非因先接受了我们的任何恩惠，而是祂无所欠缺，仅凭自己的良善就施恩于人。祂应许你天国，赐予不朽的生命、光荣、兄弟情谊、义子地位。祂使你与祂的独生子同为承继人。在如此大的恩惠之后，你还记挂你的丈夫吗？他赐予过这类东西吗？祂使太阳照耀，赐下雨水，用年岁的出产供养你。我们有祸了，因为我们如此忘恩负义！\par

为此，祂带走你的丈夫，使你不去寻求他。但即使他离世，你仍依附他，而抛弃天主，尽管你本应感恩，将一切交托给祂？因为你从丈夫那里得到了什么？分娩的痛苦、劳苦、常有的侮辱和责骂、斥责和怒气。这些不正是从丈夫而来的吗？但你却说，也有其他好处。那么是哪些呢？他用昂贵的衣服装扮你的美貌吗？他把金饰放在你脸上吗？他使你在众人中受尊重吗？但如果你愿意，你可用比离世者更好的装饰来打扮自己。因为端庄使拥有者比金饰更令人钦佩。这位君王也有衣裳，不是这种，而是好得多的。如果你愿意，穿上那衣裳。那么是怎样的衣裳？有一件有金边的衣裳；如果你愿意，就装饰灵魂。但他使你不在人前被轻视吗？那有什么了不起？你的寡居使你在魔鬼前不被轻视。你曾管辖你的仆人，至少你若曾管辖过他们。但如今，你不再管辖仆人，却管辖无形的权势、掌权者、世界的统治者。你却不提你与他同受的烦恼：有时惧怕官长，有时邻舍得宠。如今你从这一切中解脱了，从恐惧和害怕中。但你担心谁将抚养留下的孩子？\Quote{孤儿的父亲。}请告诉我，谁生了他们？你没有听基督在福音中说：\Quote{生命不是贵于食物，身体不是贵于衣服吗？} \BibleRef{玛 6:25}\par

你看见吗？你的哀哭并非因失去他的陪伴，而是由于缺乏信德。但一个已故父亲的孩子并不同样出色。为什么？因为他们有天主作父亲，难道他们不出色吗？我能指给你看多少由寡妇抚养长大而成为著名的人，又有多少在父亲抚养下而毁掉的人！因为如果你从幼年就按他们应当被抚养的方式抚养他们，他们将享受到比父亲保护大得多的益处。因为抚养子女是寡妇的本分——听保禄说：\Quote{若她抚养了儿女} \BibleRef{弟前 5:10}；又说：\Quote{她必因生育而得救}（他没有说因丈夫），\Quote{倘若她们持守信德、爱德、圣德，并以端庄自持。} \BibleRef{弟前 2:15}从幼年就给他们灌输对天主的敬畏，祂会保护他们胜过任何父亲；这将是一道不可摧毁的壁垒。因为当有护卫坐在里面时，我们无需外面的谋略；但若护卫不在，我们一切外在的谋略都徒劳。\par

这为他们将成为财富、光荣和装饰。这将使他们光彩夺目，不在地上，甚至在天上。不要看那些束金带的人，也不要看那些骑马的人，或那些因父辈而在王宫中闪耀的人，更不要看那些有仆人和教师的人。因为这些事或许使寡妇为她们的孤儿哀哭，心想：我的儿子若其父尚在，也该享有如此福气；但现在他处于沮丧和羞辱之中，不值一提。妇人，不要想这些，要在思想上为你敞开天堂之门，默观那里的宫殿，注视那坐在那里的君王。试想，在地上的人能比你的儿子在那里更显赫吗？——然后你才叹息。但若有些人在世上名声好，这也不值一提。只要你愿意，他可以成为天国的士兵，被编入那支军队的行列。因为那里入伍的人不骑马，而是乘云。他们不行走在地上，而是被提到天上。他们没有奴隶在前开路，而是有天使本身。他们不站在会死的君王面前，而是站在那不死者、万王之王、万主之主面前。他们腰间没有皮带，却有那说不出的荣耀，他们比君王或任何最显赫的人更为辉煌。因为在那些王室中，不需要财富，不需要高贵出身，只需要德行；哪里有德行，要获得首席就什么也不缺。\par

如果我们愿意培养智慧，没有什么对我们来说是痛苦的。抬头看看天，它比宫殿的屋顶华丽得多。如果上面的宫殿的地面比下面的宏伟得多，以至相比之下，下面不过是尘土；若有人被认为配得完全看见那些宫殿，他将拥有何等的福乐！\par

他说：\BibleQuote{但那真正是寡妇，且无依靠的，只寄望于天主。}\BibleRef{弟前 5:5} 这是对谁说的？对那些没有子女的人说的，因为他们更受赞许，有更多机会悦乐天主，因为对他们来说，所有的锁链都已松开。没有人拘禁他们，没有人强迫他们拖着锁链跟随。你与丈夫分离了，却与天主结合了。你没有同仆为伴，却有你的主。当你祈祷时，告诉我，你不是在与天主交谈吗？当你阅读时，听祂在与你交谈。祂对你说什么？比你丈夫更温柔的话。因为即使你的丈夫恭维你，这荣誉也不大，因为他是你的同仆。但当主恭维奴隶时，那才是伟大的求爱。祂如何求爱我们？听祂用什么方式。祂说：\BibleQuote{凡劳苦和负重担的，你们都到我这里来，我要使你们安息。}\BibleRef{玛 11:28} 又藉着先知呼唤说：\BibleQuote{妇女岂能忘记她怀中孕育的胎儿？但即使妇女忘了，我也不会忘记你，上主说。}\BibleRef{依 49:15} 这些话是何等的大爱！又说：\BibleQuote{归向我吧}\BibleRef{依 45:22}；另处又说：\BibleQuote{归向我吧，你们必得救。}\BibleRef{依 44:22} 若有人愿意从\BibleBook{雅}{歌}中选取，以更奥秘的方式理解，他将听到祂与那适合祂的灵魂交谈并说：\BibleQuote{我的佳偶，我的鸽子。}\BibleRef{歌 2:10} 有什么比这些话更甘甜？你看到天主与人的对话了吗？但是什么？告诉我，你没有看到那些有福的妇人的许多儿女已去世，躺在坟墓里吗？又有多少遭受更严酷的苦难，与丈夫一起也失去了儿女？让我们注意这些事，让我们为这些事操心，那就没有什么对我们来说是沉重的，我们将持续在属灵的喜乐中度过一生；我们将享受永恒的福乐，愿我们都藉着恩宠和慈爱等分享这福乐。\par

\SourceRecord{Translated by John A. Broadus.；Nicene and Post-Nicene Fathers, First Series；Vol. 13.；Edited by Philip Schaff.；Buffalo, NY: Christian Literature Publishing Co.,；1889.；<http://www.newadvent.org/fathers/230406.htm>.}

% structure: p0001 source=h1 target=section reason=page-title

\section{得撒洛尼前书讲道集第七篇}

\BibleRef{得前 4:13}\par

\BibleQuote{弟兄们，关于安眠的人，我们不愿你们不知道，免得你们悲伤，像其他没有望德的人一样。}\par

许多事情单因无知就使我们悲伤，如果我们好好了解，就能消除忧愁。保禄也这样指出，说：\BibleQuote{我不愿你们不知道，免得你们悲伤，像其他没有望德的人一样。}难道为此你们就不愿他们知道吗？但你为何不提所定下的惩罚？他说，无知，是指复活的道理。但为什么？这一点从另一面看是明显的，而且也已被承认。但与此同时，还有这并非微不足道的收益。因为他们并不怀疑复活，然而仍然哀恸，因此他这样讲。他对那些不信复活的人用一种方式讲论，但对这些人却用另一种方式。因为显然，那些询问\BibleQuote{时间和时期}的人是知道的。\BibleRef{得前 5:1}\par

第14节。\BibleQuote{因为我们相信}，他说，\BibleQuote{耶稣死了，也复活了}，并且活着，\BibleQuote{同样，那些在耶稣内安眠的人，天主也要领他们与祂一同前来。}\BibleRef{得前 4:14}\par

否认肉躯的人在哪里？因为如果祂没有取肉躯，祂也没有死。如果祂没有死，祂也没有复活。那么他怎么从这些事上勉励我们有信德？照他们的说法，他岂不是一个轻浮的人和欺骗者？因为如果死亡出自罪恶，而基督没有犯罪，他现在怎么鼓励我们？如今，关于谁他说话呢，诸位？你们为谁哀悼？为谁悲伤？为罪人，还是单纯为那些死去的人？他为什么说：\BibleQuote{像其他没有望德的人一样}？\BibleQuote{从死者中首生者}，\BibleRef{哥 1:18}他说，初果。因此也必须有其余的人。并且你看，他在这里没有引入任何推理，因为他们易于受教。因为他在写给格林多人的书信中，也从推理中提出许多东西，然后他补充说：\BibleQuote{糊涂人！你所种的若不先死，决不会活过来。}\BibleRef{格前 15:36}因为这更有权威，但这是当他与信徒交谈时。但对于教外的人，这有什么权威呢？\BibleQuote{同样，}他说，\BibleQuote{那些在耶稣内安眠的人，天主也要领他们与祂一同前来。}再次，\Quote{安眠}：他绝不说\Quote{死者}。但关于基督，他的话说：\BibleQuote{祂死了}，因为后面提到复活，但这里说\Quote{安眠的人}。如何\BibleQuote{借着耶稣}？或者他们借着耶稣安眠，或者借着耶稣祂要带领他们。短语\BibleQuote{借着耶稣安眠}指信徒。这里异端者说，他指的是受洗的人。那么\Quote{同样}还有什么余地？因为耶稣并非借着洗礼安眠。但他为何说\Quote{安眠的人}？这样他论及的不是普遍的复活，而是部分的复活。他说，借着耶稣安眠的人，并且他在许多地方这样讲。\par

第15节。\BibleQuote{我们照主的话告诉你们这件事：我们还活着存留到主来临的人，决不会在那些安眠的人以先。}\BibleRef{得前 4:15}\par

谈及信徒和\BibleQuote{在基督内安眠的人}\BibleRef{格前 15:18}；又说：\BibleQuote{死人在基督内复活。}因为他的言论不仅关于复活，也关于复活和荣耀中的尊荣；那么所有人都要参与复活，他说，但并非所有人都在荣耀中，只有在基督内的。既然他愿意安慰他们，他不仅以此安慰他们，而且以丰富的尊荣和它的迅速来临，因为他们知道这一点。为证明他以尊荣安慰他们，他在下文说：\BibleQuote{我们将永远与主同在}；并且\BibleQuote{我们将被提到云中。}\par

但信徒如何在耶稣内安眠？意思是他们有基督在自己内。但\BibleQuote{祂要领他们一同前来}这一表达显示他们是从许多地方被领来。\Quote{这件事。}他要告诉他们一件奇怪的事。为此他也加上了使之可信的条件：\BibleQuote{从主的话}，他说，即我们不是出于自己说话，而是从基督学来的，\BibleQuote{我们还活着存留到主来临的人，决不会在那些安眠的人以先。}这也是他在写给格林多人的书信中所说的：\BibleQuote{在一刹那，一眨眼之间。}\BibleRef{格前 15:52}这里他借着方式也给了复活可信性。\par

因为这事似乎困难，他说，如同活着的人被提去是容易的，那么离世的人也是如此。但他说\Quote{我们}，并不是指他自己，因为他不会存留到复活，而是指信徒。为此他补充说：\BibleQuote{我们还活着存留到主来临的人，决不会在那些安眠的人以先。}好像他说：不要以为有任何困难。是天主成就这事。那时活着的人不会领先于那些已消散、已腐烂、已死了万年的人。但如同带来完整的人是容易的，同样，带来已消散的人也是容易的。\par

但有些人因不认识天主，而不信这事实。请告诉我：从无中造出一个人来，与使已分解的人复活，哪一样更容易？可是他们说些什么？某人遭遇海难，淹死在海中，许多鱼吞吃了他，每一条鱼吃掉了他的一些肢体。然后这些鱼中，一条在这个海湾被捕获，另一条在那个海湾被捕获；这鱼被人吃了，那鱼也被人吃了，而它们腹中藏有被吞吃的人肉。再者，那些吃了鱼的人——鱼曾吃掉了那人——在不同地方死去，他们自己或许也被野兽吞吃了。当有了如此巨大的混乱和分散时，那人怎能复活呢？谁能把灰尘收集起来？可是，人啊，你为什么说这话，编织琐碎的线索，使之成为难题？请告诉我：如果那人没有掉进海里，如果鱼没有吃他，鱼也没有被无数的人吃掉——而是他小心地被保存在棺材里，没有虫子或任何东西打扰他——那么那已分解的如何复活？灰尘和灰烬如何再次黏合？身体如何能再恢复其光彩？这难道不是一个困难吗？\par

如果提出这些疑问的是希腊人，我们会有无数的话对他们说。那么怎样？因为他们当中有一些人把灵魂转移到植物、灌木和狗身上。请告诉我：恢复自己的身体，还是恢复别人的身体，哪一个更容易？另一些人又说：他们被火焚烧，但衣服和鞋会复活，他们却不受嘲笑。还有人说原子。对这些人，我们根本不予争论；但对于信徒——如果我们仍可称那些提出疑问的人为信徒的话——我们还要说宗徒所说的话：一切生命都来自腐朽，所有植物，所有种子。你没有看见无花果树吗？它有怎样的树干，怎样的枝条，多少叶子，多少枝杈、茎和根，占据了那么多土地并深植其中？这棵如此巨大、如此繁茂的树，是从一粒种籽长出来的，这粒种籽被扔进地里，先腐败了。若它不腐烂、不分解，就不会有这些东西。请告诉我：这是怎样发生的？葡萄树也是如此，它看起来和吃起来都那么美好，却是从外表丑陋的东西长出来的。请告诉我：从天上降下的水不是同一种东西吗？它怎样变成了如此多的东西？这比复活更奇妙。因为在复活中，同一粒种籽和同一棵植物是主体，有很大的相似性。但在这里：请告诉我，水具有同一性质、同一本性，却变成如此多的东西？在葡萄树中，它变成酒，不仅是酒，还有叶子和汁液。因为不仅葡萄串，连葡萄树的其余部分也由它滋养。又在橄榄树中，它变成油，以及其他无数东西。奇妙的是：在这里它是湿的，在那里是干的；在这里是甜的，在那里是酸的；在这里是收敛的，在那里是苦的。请告诉我：它怎样变成如此多的东西？给我理由！但你办不到。\par

再就你自身而言——因为更接近——那被播下的种籽，怎样被塑造成如此多的东西？怎样变成眼睛？怎样变成耳朵？怎样变成手？怎样变成心脏？在身体中，形貌、大小、性质、位置、能力、比例有成千上万种差异，不是吗？神经、血管、肌肉、骨骼、膜、动脉、关节、软骨，以及除此之外许多由医生之子精确指定的东西，构成了我们的本性——这些都出自那一粒种籽！这岂不比那些事更困难得多吗？潮湿柔软的东西怎样凝结成干燥冰冷的东西——即骨头？怎样变成温暖潮湿的东西——即血中混合的？怎样变成寒冷柔软的东西——即神经？怎样变成寒冷潮湿的东西——即动脉？请告诉我：这些事从何而来？你对这些事不感到茫然吗？难道你没有看见每天在我们的生命时期中发生一种复活和死亡吗？我们的青春去哪里了？我们的老年从哪里来？为什么老去的人不能使自己变年轻，却能生出另一个非常年幼的孩子，而且他自己不能给予自己的，却给予别人？\par

我们在树木和动物身上也能看到这一点。然而，那给予别人的，本应首先馈赠自身。但这属人的理性要求。但当天主创造时，让一切都让步。如果这些事是那么困难，甚至极其困难，我就想起那些疯狂的人，他们好奇于圣子的\OriginalTerm{无形的生}{incorporeal Generation}。每天发生的事、在我们手中的事、被追问了千万次的事，还没有人能发现；那么请告诉我：他们怎能好奇于那隐秘的和不可言说的\OriginalTerm{生}{Generation}？这种人的心智在虚空中跋涉，不疲惫吗？不被卷入千万种眩晕中吗？不目瞪口呆吗？然而即使这样他们也不受教。当他们不能就葡萄和无花果说任何话时，他们却好奇于天主！请告诉我：那葡萄核怎样分解为叶子和茎？此前它们如何不在其中、不被看见？但你说：不是来自葡萄核，而是全部来自土地。那么，为什么没有这种子，土地本身什么也不产生？但让我们不要愚昧无知。所发生的既不是来自土地，也不是来自葡萄核，而是来自那既是土地又是种籽的主。为此，祂使同一事物既没有它们也能产生，有它们也能产生。首先，祂显示自己的权能，当祂说：\Quote{地要生出青草。}\newline{}\BibleRef{创 1:11}其次，除了显示权能，也教导我们勤劳努力。\par

我们爲何说了这些事呢？不是徒然，而是爲叫我们也能相信复活，并且当我们再次想凭理性理解某事却找不到时，不致恼怒而跌倒，而要谨慎退后，抑制我们的推理，投奔天主的德能与智慧。因此，知道了这些，让我们给推理套上辔头。让我们不越界，不超越爲我们知识所定的尺度。因爲他说：\BibleQuote{若有人}，\BibleQuote{以为自己知道什么，按他所当知道的，他甚么也不知道。}\BibleRef{格前 8:2}\par

我不是论及天主，他说，而是论及一切。何爲？你愿意学习关于大地的事吗？你知道什么？告诉我。它的尺度多大？它的体积多大？它的位置方式如何？它的本质如何？它的地位何在？它立于何处，靠什么支撑？这些事你一件也不能说吗？但你说它寒冷、干燥、色黑——仅此而已。再者，关于海洋呢？你同样会陷于不确定，不知它起于何处，止于何处，依托于什么，海底由什么支撑，有何种空间承载它，它之后是大陆还是止于水与空气。你又知道其中的什么？但什么？让我跳过元素。你愿意我们选取最小的植物吗？那无果的草，我们全都知道，告诉我，它是如何生长的？它的原料岂非水、土和粪？是什么使它如此美丽，具有如此令人赞叹的颜色？那美丽为何如此消逝？这不是水或土的作为。你看见处处需要信德吗？大地如何生发，如何产痛？告诉我。但你什么都说不出来。\par

人啊，学习这地下的事，不要好奇或多管天上之事。但愿只是天上，而非天上的主！你不认识你由之而生、在其中滋养、你居住、行走其上、甚至不能呼吸的大地，却好奇极远的事？实在\BibleQuote{世人虚幻}。\BibleRef{咏 39:5}若有人命你潜入深渊，探查海底的事，你必不堪承受。但无人强迫时，你竟自愿探測不可测的深渊？我恳求你，不要如此。让我们向上航行，不是漂浮，因我们很快就会疲乏而沉没；而是以\WorkTitle{圣经}爲船，扬起信德的帆。若我们航行在其中，天主的圣言必临在我们中间，作我们的领航者。若我们漂浮于人的推理，则不然。因为漂浮者中，谁有领航者同在？所以危险是双重的，既无船，又无领航者。因为连小船没有领航者也不安全，何况两者皆缺，还有何得救的希望？因此，不要让自己陷入明显的危险，而要登上安全的船只，以神圣的锚将自己系牢。这样，我们必满载货物，同时十分安全地航入平静的港湾，并获得爲爱祂的人所预备的福乐，在基督耶稣我们的主内。愿荣耀、权能、尊崇归于祂，偕同圣父及圣神，自今而始，迄无穷世。阿们。\par

\SourceRecord{Translated by John A. Broadus.；Nicene and Post-Nicene Fathers, First Series；Vol. 13.；Edited by Philip Schaff.；Buffalo, NY: Christian Literature Publishing Co.,；1889.；<http://www.newadvent.org/fathers/230407.htm>.}

% structure: p0001 source=h1 target=section reason=page-title

\section{《得撒洛尼前书》第八篇讲道}

\BibleRef{得前 4:15}\par

\Quote{我们藉主的话告诉你们这件事：我们这些活着存留到主来临的人，决不会在已经长眠的人以前。因为主必亲自从天降来，有呼喊声、有总领天使的声音，还有天主的号声；那些在基督内死了的人要先复活，然后我们这些活着还存留的人，同时与他们一起被提到云中，到空中迎接主：这样，我们就常与主同在。}\par

先知们确实为了表明他们所说之事的可信性，在所有事之前先说：\BibleQuote{依撒意亚所见的神视}（\BibleRef{依 1:1}），又说：\BibleQuote{那传给耶肋米亚的上主的话}（\BibleRef{耶 1:1}），又说：\Quote{上主这样说}，以及许多类似的表达。他们中许多人甚至看见了天主坐着，按他们所能看见的程度。可是保禄没有看见祂坐着，却有基督在他内说话，所以不说\Quote{上主这样说}，而说：\BibleQuote{你们既然寻求基督在我内说话的凭证}（\BibleRef{格后 13:3}），又说：\Quote{保禄，耶稣基督的宗徒。}因为\Quote{宗徒}所说的是派遣他者的言语，表明没有什么是出于自己的。又说：\BibleQuote{我认为我也有天主的圣神}（\BibleRef{格前 7:40}）。所以一切他都是藉圣神说的，但现在所说的，是直接从天主听见的。如同他在对厄弗所长老的讲话中所说的：\BibleQuote{施予比领受更为有福}（\BibleRef{宗 20:35}），那是在未记录的话语中听见的。\par

那么，我们来看他如今所说的话：\Quote{我们藉主的话告诉你们这件事：我们这些活着存留到主来临的人，决不会在已经长眠的人以前。因为主必亲自从天降来，有呼喊声、有总领天使的声音，还有最后的号声。}因为那时，他说：\BibleQuote{天上的万军也要震动}（\BibleRef{玛 24:29}）；可是为什么有号声？因为我们在西乃山也看见这事，那里也有天使。但\Quote{总领天使的声音}是什么意思？正如他在童女的比喻中所说的：起来！\BibleQuote{新郎来了}（\BibleRef{玛 25:6}）。这要么指此事，要么如国王的情形，那时也将如此，天使们在复活时服务。因为祂说：让死人复活，这事就完成了，天使们没有能力做这事，而是祂的话。如同国王命令并说话时，被关的人就出来，仆人们就领他们出来，但他们这样做并非出于自己的能力，而是出于那声音。基督在另一处也这样说：\BibleQuote{祂将派遣祂的天使，用大号角，从四方，从天这边到天那边，聚集祂的选民}（\BibleRef{玛 24:31}）。处处可见天使们奔波。所以我认为总领天使是管理那些被派遣者的，他这样呼喊：\Quote{让众人预备好，因为审判官临近了}。而\Quote{最后的号声}是什么意思？这里暗示有许多号声，在最后审判官降临。他说：\Quote{那些在基督内死了的人要先复活。然后我们这些活着还存留的人，同时与他们一起被提到云中，到空中迎接主：这样，我们就常与主同在。}\par

\Quote{为此，你们要以这些话彼此安慰。}（\BibleRef{得前 4:18}）\par

如果祂要降临，我们为什么被提上去？为尊荣的缘故。因为当一位国王驾车入城时，那些有尊荣的人出去迎接他，而被定罪的人则在城内等待审判官。当一位慈父来临，他的孩子和配作他孩子的人被带到车上，以便看见并亲吻他；而家中那些冒犯的人则留在里面。我们被载在父亲的车上。因为祂被接升天时在云中，\BibleQuote{我们将被提到云中}（\BibleRef{宗 1:9}）。你看这是多么大的尊荣！当祂下降时，我们出去迎接祂，而且比一切更幸福的是，我们将与祂同在。\par

\BibleQuote{谁能述说上主的大能，宣扬祂的一切荣耀？}（\BibleRef{咏 105:2}）祂赐给了人类多少恩惠！已死的人先复活，然后相遇一起发生。在所有之前死去的亚伯尔将与那些活着的人一起迎接祂。因此在这一点上他们并无优势，而是那已经朽坏、在泥土中多年的人将与活着的人一起迎接祂，所有其他人也一样。因为如果他们等待我们，使我们得以加冕，正如他在另一封书信中所说：\BibleQuote{因为天主为我们预备了更好的事，以致他们若不与我们同得，就不能成为完美}（\BibleRef{希 11:40}），那么我们更要等待他们；或者更好说，他们等待我们，我们却一点也不等待。因为复活是在\Quote{顷刻，眨眼之间}发生的。\par

关于他们被聚集的说法：他们各处复活，但被天使们聚集。前者是天主命令大地交出所托之能力的作为，其中没有服务员，正如祂那时呼唤拉匝禄：\BibleQuote{拉匝禄，出来！}（\BibleRef{若 11:43}）；而聚集是服务者的工作。但如果天使聚集他们并四处奔走，这里他们如何被提上去？他们是在下降之后、被聚集之后被提上去的。\par

因为这也在无人知晓的情况下发生。因为他们看见大地震动，尘土混合，尸体也许从各处升起，没有人服务这事，而\Quote{呼喊声}就已足够，整个大地充满（试想这事何等伟大：从亚当直到祂来临的所有人，那时将带着妻儿站立）——当他们看见地上如此大的骚动时，那时他们就会知道。因此，正如在肉身的救恩工程中，他们事先一无所知，那时也将如此。\par

当这些事发生时，也将有总领天使呼喊并命令天使的声音，以及号角声，或更确切地说，号筒的响声。那时，留在地上的人将是何等的战栗和恐惧！因为一个女人被提去，另一个被留下；一个男人被取去，另一个被撇下。\BibleRef{玛 24:40} 当他们看见一些人被提去，而自己却被留下时，他们的灵魂将处于何种状态？这些事岂不比任何地狱更可怕地摇撼他们的灵魂吗？让我们现在就以言语描绘这情景，仿佛它已临到。因为若突然的死亡、城中的地震和威胁如此使我们的灵魂战栗，那么当我们看见大地裂开、充满这一切，听见号角声和比任何号角更响亮的总领天使的声音，察觉天卷起、万有的天主亲自降临时——我们的灵魂又将如何？我恳求你们，让我们战栗恐惧，仿佛这些事正在发生。不要因延迟而自慰；因为当它必然发生时，延迟对我们毫无益处。\par

那时的恐惧和战栗将何其大！你们见过被押赴刑场的人吗？你们想，当他们走向城门时，灵魂处于何种状态？岂不是比许多死亡更痛苦？他们岂不愿做任何事、受任何苦，以便从那种黑暗的阴云中得救吗？我曾听许多人说，那些在被押解途中因君王的仁慈而被召回的人，甚至不视人为人，因为他们的灵魂如此惊慌、恐惧、魂不守舍。如果身体的死亡尚且如此惊吓我们，那么当永恒死亡来临时，我们又将如何？我为何提及那些被押解的人呢？那时周围站着一群人，大部分人甚至不认识他们。若有人窥探他们的灵魂，没有人会如此残忍、如此铁石心肠、如此刚硬，以至于不使自己的灵魂颓丧，并因恐惧和绝望而松弛。若他人被这种与睡眠无异的死亡带走时，不相干的尚且如此感受，那么当我们自己陷入更大的灾祸时，我们会是怎样的情形？相信我，言语无法形容那种痛苦。\par

但你会说：天主是爱人者，这些事都不会发生！那么圣经是白写了？不，你说，这只是威胁，好使我们变得明智！若我们不变得明智，反而继续作恶，祂难道不会，告诉我，施加惩罚吗？祂难道不会奖赏善人吗？是的，你说，因为行善甚至超越应得，是合宜的。所以，那些事确实是真的，必定会发生，而惩罚却根本不会发生，只是用于威胁和恐吓罢了！我不知道用什么方法才能说服你们。若我说：\BibleQuote{虫子不死，火也不灭}\BibleRef{谷 9:44}；若我说：\BibleQuote{他们要去到永火里}\BibleRef{玛 25:41}；若我把那已受惩罚的富贵者摆在你们面前，你们会说这都是威胁。那么我该如何说服你们呢？因为这是撒殚的推论，它用无益的宽容纵容你们，使你们懒惰。\par

我们如何才能消除它呢？凡我们所说的圣经之言，你们都会说是威胁。关于未来的事，这或许可以说，但对于已经发生并终结的事却不能这样说。你们听说过洪水。那些事也是威胁之言吗？它们不是实际发生了吗？那些人也说了许多这样的话，在方舟建造的一百年间，木头在被加工，义人大声呼喊，却无人相信。但因他们不信言语的威胁，便确实受了惩罚。我们若不信仰，也将如此。因此，祂将祂的来临与诺厄的日子相比，因为正如有些人不信那洪水，同样他们也会不信地狱的洪水。这些事是威胁吗？它们不是事实吗？那么，那位当时突然降罚于他们的，岂不更要如今也降罚吗？因为现在所犯的罪恶不小于那时代的过犯。如何呢？——因为那时，经上说：\BibleQuote{天主的儿子们进了人的女儿们那里}\BibleRef{创 6:4}，那些混杂是极大的罪过。但现在每一种恶行都有人尝试。那么你们相信洪水发生过吗？或者它看来像个寓言？然而就连方舟所停的山脉也作证，我说的是在\Place{亚美尼亚}的那些山。\par

不仅如此，甚至过分地，我将我的论述转向另一件比那更明显的事。因为我将不再提传闻，而是事实，然而那些比事实更清楚。因为圣经所说的一切，比我们所看见的事物更值得信赖。你们当中有谁曾到过巴勒斯坦旅行？我想是的。那么怎样呢？你们这些看见了那些地方的人，请为我在那些未曾到过的人面前作见证。因为从阿什凯隆和加沙直到约但河的尽头，有一片广阔而肥沃的土地——或者更准确地说，曾经有过——因为如今已不复存在。这曾经是如同园子一样的地方。因为经上记着说：\BibleQuote{罗特看见了约但河整个平原——且处处有水灌溉，如同上主的园子。}\BibleRef{创 13:10} 因此，这片如此繁茂、胜过所有地区、繁荣程度甚至超过天主的乐园的地方，如今比任何荒野都更荒凉。那里确实立着树木，并且结出果实。但果子是天主义怒的纪念碑。因为那里有石榴树，我说的是树和果子，外观非常好看，给无知的人以巨大的希望。但若拿在手中，剥开来，里面却没有果子，而是内中积满了尘埃和灰烬。整个土地也是如此。你若找到一块石头，会发现它满是灰烬。我何必提石头、木头和土呢？难道空气和水也参与了这个灾难吗？因为正如身体被焚烧耗尽，形状保留，出现火的轮廓，体积和比例还在，但能力不复存在；同样，在那里你看到土，却没有土的实质，尽是灰烬；树木和果子，却没有树木和果子的实质；空气和水，却没有水和空气的实质，因为连这些也变成了灰烬。然而空气怎能被烧？水怎能被烧，当它还是水时？木头和石头确实可烧，但空气和水完全不可能。对我们不可能，但对成就这些的那一位是可能的。所以，空气不过是火炉，水也是火炉。万物都不结果，全无产出，全是为了报复；是先前的激怒的图像，以及将来之事的证据。\par

这些也不过是威胁的话吗？这些不过是言语的声音吗？因为于我而言，前者并非不可信，未见之事与所见之事同样可信。但即便是对不信者，这些也足以产生信德。若有人不信地狱，就让他思考索多玛，让他默想哈摩辣，那已经施行且仍存留的惩罚。这是永恒惩罚的明证。这些事严酷吗？当你说没有地狱，天主不过是恐吓而已，这难道不严酷吗？当你松懈百姓的手时？是你这不信者迫使我言说这些；是你将我引出了这些话。如果你相信基督的话，我就不必被迫提出事实来诱导信德。但既然你回避了它们，你从此将被说服，无论愿意与否。关于索多玛，你还有什么可说？你是否还想知道当初行这些事的原因？那是一个罪，自然是严重而可咒骂的罪，但仅此一个。那时的人贪恋男童，因此遭受了这惩罚。但现在有成千上万相等甚至更严重的罪被犯下。那位曾为一个罪倾泄如此多的愤怒，且不顾亚巴郎的恳求，也不顾住在他们中间的罗特——那人出于对祂的仆人的尊敬，献出自己的女儿受凌辱——那么当有这么多罪时，他会宽恕吗？这些事真是可笑、琐碎、幻觉，以及魔鬼的欺骗！\par

你是否希望我再提出另一个例子？你当然听说过埃及王法郎；因此你也知道他受的惩罚，即他和整个军队、战车、马匹等一切，如何淹没在红海中。你还愿听其他的例子吗？他或许是一个不虔敬的人，或者并非\InnerQuote{或许}，而肯定是一个不虔敬的人。你愿看见那些属于信徒之列、紧握天主不放但没有正直生活的人，也受了惩罚吗？请听保禄说：\BibleQuote{我们也不要行邪淫，如同他们中有些人行了邪淫，一天内就倒毙了二万三千人。我们也不要抱怨，如同他们中有些人抱怨了，就被毁灭者所灭。我们也不要试探基督，如同他们中有些人试探了，就被蛇所灭。}\BibleRef{格前 10:8} 如果邪淫和抱怨有如此大的力量，那么我们的罪会有怎样的后果呢？\par

如果你现在没有受惩罚，不要惊讶。因为他们不知道有地狱，因此惩罚紧随着他们的脚后跟而来。但你，无论你犯什么罪，即使你逃过现世的惩罚，你将在那里全部偿还。祂如此惩罚那些几乎处于孩童状态、且没有犯下如此大罪的人——祂还会饶恕我们吗？这不合理。因为如果我们与他们犯同样的罪，我们当受比他们更大的惩罚。为什么？因为我们领受了更多的恩宠。但当我们的过犯众多，且比他们的更可憎时，我们将会遭受怎样的报复呢？他们——请没有人以为我这样说是在赞美他们或为他们开脱；绝无此事：因为当天主惩罚时，谁若作出相反的判决，是出于魔鬼的暗示。我这样说，不是赞美他们，也不是开脱他们，而是为显示我们的邪恶——他们虽然抱怨了，但他们是在旷野路上行走；而我们抱怨，尽管我们有家乡，住在自己的房屋里。他们虽然行了邪淫，但那是在他们刚脱离埃及的苦难之后，几乎还未听说过这样的法律。但我们却是在从祖先那里预先领受了救恩的道理之后去行的，因此我们应受更大的惩罚。\par

你愿听听别的事吗？他们在巴勒斯坦的苦难，在巴比伦人、亚述人手下所遭遇的饥荒、瘟疫、战争、被掳，以及来自马其顿人、哈德良和韦斯帕芗时代的苦难，又是如何？亲爱的，我有一件事想告诉你们；不要跑开！或者我先前再告诉你们另一件事。经上记载，有一次发生饥荒，王正在城墙上行走；一个妇人来向他说了这些话：\BibleQuote{这个妇人对我说：把你的儿子今天煮了，我们吃吧——明天吃我的。我们就把我的儿子煮了吃了，现在她却不肯把她儿子给我。}\BibleRef{列下 6:28}还有什么比这灾祸更可怕的呢？又在另一处，先知说：\BibleQuote{慈怜妇人的手，煮了自己的儿女。}\BibleRef{哀 4:10}犹太人那时受了这样的惩罚，我们岂不是更该受罚吗？\par

你们还愿听他们别的灾祸吗？请读一读约瑟夫，你们就会知道那整个悲剧，或许我们能由此说服你们，确实有地狱。请想，如果他们受了罚，我们为什么不受罚？或者，我们比他们犯更重的罪，现在却不受罚，这合理吗？岂不明显是因为惩罚为我们存留了吗？如果你们愿意，我就逐一个人告诉你们他们如何受了罚。加音杀了他的兄弟。这真是可怕的罪，谁能否认呢？但他受了罚，而且是很重的罚，抵得上万次死亡，因为他宁可死一万次。请听他说的：\BibleQuote{祢把我从这地上赶走，使我不能见祢的面，那么遇见我的人必会杀我。}\BibleRef{创 4:14}那么请问，如今不是有许多人做同样的事吗？当你杀的不是你的肉身兄弟，而是你的属灵弟兄时，你不也是做同样的事吗？虽然不用刀剑，但用别的手段：当你能够解除他的饥饿时，却不顾他。那么，如今没有人嫉妒他的弟兄吗？没有人把他陷入危险吗？但在这里他们没有受罚，然而将来他们要受罚。他从未听过成文的法律，也未听过先知，没见过大奇迹，却受了这么重的报应；那么，以别的方式做了同样的事，而且没有被这么多榜样所警醒的人，难道能不受罚吗？那么，天主的公义在哪里？祂的慈爱又在哪里呢？\par

又有一人因在安息日捡柴而被石头砸死，这不过是个小诫命，比割损更轻。那在安息日捡柴的人被石头砸死了；而那些常常违反法律做千万件事的人，却未受罚而去！如果没有地狱，祂的公义在哪里？祂不偏待人，那不偏待人的在哪里？然而祂还责备他们许多这样的事，说他们没有守安息日。\par

又有一人，加尔米，因偷窃了奉献之物，被石头砸死，连同他全家。那么，从那时起，没有人犯过亵圣罪吗？撒乌耳又因违背天主的命令而饶命，受了那么大的惩罚。从那时起，没有人饶过命吗？但愿没有！我们岂不更是如同野兽一样，违背天主的命令彼此吞噬，却没有人阵亡？厄里的儿子们，因在焚香之前吃了祭肉，连同他们的父亲受了最严厉的惩罚。那么，没有父亲对子女疏于管教吗？没有邪恶的儿子吗？但没有人受罚。那么，如果没有地狱，他们将在哪里受罚呢？\par

此外还可列举无数例子。阿纳尼雅和撒斐辣因私藏了所献的一部分而被立即惩罚。那么，从那时起没有人犯过这罪吗？为什么他们没有受同样的惩罚呢？\par

那么，我们说服你们有地狱了吗？还是你们需要更多例子？因此，我们也要从不成立的事例，就是从生活中现今发生的事来谈。因为我们必须从各方面收集这个观念，免得我们徒然自娱而害了自己。你们没有看见许多人遭受灾祸，身体残疾，受无穷的苦难，而另一些人却名声很好吗？为什么有人因杀人受罚，有人却没有？请听保禄说：\BibleQuote{有些人的罪是明显的，……有些人的罪是随后跟来的。}\BibleRef{弟前 5:24}多少杀人犯逃脱了！多少盗墓者！但让这些事过去吧。你们没有看见多少人受到最严厉的惩罚吗？有人被交付于长期的疾病，有人受持续的折磨，还有人受无数其他的灾祸。因此，当你看见一个人犯了与他们相同的罪，甚至更重，却未受罚时，你岂不连自己都不情愿地怀疑有地狱吗？想想那些在你之前受了重罚的人，想想天主是不偏待人的，尽管你做了无数恶事，却没有受这样的苦，你就会有地狱的观念了。因为天主已将这个观念深植在我们心中，以致没有人能不知道它。因为诗人、哲学家和寓言家，总之全人类，都讨论过那里的报应，说多数人在阴府受罚。如果那些是寓言，但我们所领受的却不是。\par

我说这些，并非想吓唬你们，也不是给你们的灵魂增添重担，而是要使它们变得明智，并使它们更轻省。我自己也宁愿没有惩罚，是的，我比任何人都更希望如此。为什么呢？因为你们每个人为自己灵魂担忧之时，我还得为管理你们这项职务负责。所以，我最无法逃脱。但惩罚和地狱不可能不存在。我有什么办法呢？那么他们问，天主对人的仁慈在哪里？在许多地方。但我宁愿在别的时候再探讨这个问题，免得我们混淆了关于地狱的讲论。暂且不要让我们已经获得的益处溜走。因为关于地狱的确信并非小益；回忆这样的讲论，如同苦药，若能常存于心中，便能清除一切恶习。因此，让我们善用这药，好使我们因纯洁的心，而得以配看见那些\BibleQuote{眼所未见，耳所未闻，人心所未曾想到的}。愿我们都能藉着我们的主\Person{耶稣基督}的恩宠和仁慈获得这一切。\par

\SourceRecord{Translated by John A. Broadus.；Nicene and Post-Nicene Fathers, First Series；Vol. 13.；Edited by Philip Schaff.；Buffalo, NY: Christian Literature Publishing Co.,；1889.；<http://www.newadvent.org/fathers/230408.htm>.}

% structure: p0001 source=h1 target=section reason=page-title

\section{关于《得撒洛尼前书》的第九篇讲道}

得撒洛尼前书 5:1-2\par

\BibleQuote{但弟兄们，论到时候和日期，不需要给你们写什么。因为你们自己深知主的日子来临，如同夜间的盗贼一样。}\BibleRef{得前 5:1-2}\par

似乎没有什么比人的本性更为好奇，更倾向于窥探隐晦和隐秘的事物。当心灵不安定、处于不完美状态时，这种情况常会发生。因为较为简单的孩童会不断地用他们频繁的问题来纠缠他们的保姆、导师和父母，其中无非是\Quote{这将在何时？}和\Quote{那将在何时？}这也是由于生活放纵和无所事事所致。因此，我们的心灵急于已经学习和领会许多事物，尤其是关于终结的时期；如果我们被如此影响，有什么奇怪呢？因为那些圣人自己也是最被这样影响的。在受难之前，宗徒们来到基督面前说：\BibleQuote{请告诉我们，什么时候要有这些事？你来临和世界终结的预兆是什么？}\BibleRef{玛 24:23}在受难和从死者中复活之后，他们对他说：\BibleQuote{主，就在这时候你要复兴以色列国吗？}\BibleRef{宗 1:6}他们问他的事，没有比这更快的了。\par

但后来，当他们被赐予圣神时，就不再是这样了。他们不仅自己不询问，也不抱怨这种无知，反而压制那些受不合时宜的好奇心之苦的人。例如，请听圣保禄现在所说的话：\BibleQuote{但论到时候和日期，弟兄们，不需要给你们写什么。}\BibleRef{得前 5:1}为什么他没有说无人知道？为什么他没有说还没有启示，而是说\BibleQuote{不需要给你们写什么}？因为那样他会更加使他们忧伤，而这样说话却安慰了他们。因为通过\Quote{不需要}这个说法，似乎既是多余又是不适当，他不让他们询问。\par

请告诉我，那有什么益处？假设终结将在二十、三十或一百年之后，这对我们有什么关系？每个人自己生命的终结难道不就是个人的终结吗？你为什么好奇，为普遍的终结而忧虑？在这件事上，我们和其他事上一样。因为在其他事上，我们离开自己私人的事务，却为普遍的事忧虑，说：某人是奸淫者，某人是通奸者，那人抢劫了，另一人伤害了；但没有人考虑自己的事，每个人都想着别人的私事；同样，在这里，每个人都忽略思考自己的终结，却急于听到关于普遍解散的事。那与你有什么相干？因为如果你使自己的终结成为好的，你就不会因另一个而受害；无论是遥远还是临近。这对我们无关紧要。\par

因此，基督没有告诉他们，因为那并不适宜。你可能会说，怎么不适宜？他知道为什么不适适宜。因为听他对他宗徒们说：\BibleQuote{父凭自己的权柄所定的时候、时期，不是你们可以知道的。}\BibleRef{宗 1:7}你为什么好奇？伯多禄，宗徒之长，和他的同伴，听到这样说，好像他们寻求的是过于他们所能知道的事。不错，你说；但这样就能堵住希腊人的口。如何？请告诉我。因为他们说，这个世界是一位神；如果我们知道它解散的时期，我们就能堵住他们的口。为什么？知道它何时被毁灭，还是知道它将被毁灭，会堵住他们的口？告诉他们，它会有终结。如果连这个他们都不信，他们也不会相信另一个。\par

请听圣保禄说：\BibleQuote{你们自己深知主的日子来临，如同夜间的盗贼一样。}\BibleRef{得前 5:2}不仅是指普遍的那一天，也是指每个人的那一天。因为一个与另一个相似，也与之相近。因为普遍的那一天从亚当起就开始了，然后是终结的完成；因为即使现在，称它为完成也不为过。因为每天有千万人死去，所有人都等候那一天，没有人在那一天之前复活，这难道不是那一天的作为吗？如果你想知道为什么它是被隐藏的，为什么它像夜间的盗贼一样来临，我会告诉你我认为可以很好解释的理由。没有人会在整个生命中修德；因为知道自己最后的一天，在犯了无数罪之后，然后来到\OriginalTerm{洗池}{Laver}，就会这样离去。因为如果现在，当因不确定性而产生的恐惧震撼所有灵魂时，所有人仍然在罪恶中度过了大半生，在临终时才领受圣洗——如果他们对此事完全确信，谁还会修德呢？如果许多人没有领受\OriginalTerm{光照}{Illumination}就离世，甚至这种恐惧也没有教导他们在活着时修德行事悦乐天主；如果连这种恐惧也被除去，谁还会清醒，谁还会温和？没有一个人！还有另一件事。对死亡的恐惧和对生命的爱约束了许多人。但如果每个人都知道他明天肯定会死，在那一天之前，他什么都不会拒绝尝试，他会随意杀人，通过向敌人复仇来逞快，并犯下无数罪行。\par

对于恶人，他们对今生的生命绝望，甚至不尊重身披紫袍的人。因此，凡确信自己无论如何必死的人，既向仇敌报复，又先满足了自己的灵魂，然后才死去。我再提第三点。贪生怕死、执著于世物的人，会因绝望和忧愁而毁灭。因为若有青年人知道自己未及老年就将死，如同最懒散的野兽被捕获后因等待死亡而变得更加懒散，他也将受到同样的影响。此外，即使是勇敢的人也不会得到赏报。因为他们若知道三年后必死，且在此之前不可能死，那么他们冒险犯难又怎能得到赏报呢？因为任何人都可能对他们说：\Quote{你们因为确信还有三年寿命，所以才投身危险之中，因为知道在这之前不会死去。}因为那期待每次危险都可能丧命的人，知道如果他不置身危险就能活下去，但若他尝试某些行动就会死，他就证明了他最大的热忱和对现世生命的轻视。我要用一个例子向你们说明。请问，如果族长亚巴郎预知他不必献祭自己的儿子，就把儿子带到那里，他还能得到赏报吗？又如保禄预知他不会死，就轻视危险，他又哪里值得钦佩呢？因为即使是极懒散的人，如果能找到一个可信赖的人确保安全，也会冲进火里去。但三位青年却不是这样。请听他们说：\BibleQuote{大王！有一位天主在天上，祂必从您手中和这火窑中拯救我们；即或不然，也请您知道，我们不侍奉您的神明，也不朝拜您所立的金像。}\BibleRef{达 3:17}\par

你看，有多少益处，还有比这些更多的益处来自不知道我们结局的时刻。同时，学习这些就足够了。为此，祂像夜间的贼一样来到；好叫我们不陷入邪恶，也不陷入懈怠；好叫祂不夺去我们的赏报。他说：\Quote{你们自己完全知道。} 那么，你们既然相信，为什么还要好奇呢？但未来不确定，从基督所说的可以得知。因为为此祂说了这话，请听祂说什么：\BibleQuote{所以，你们要警醒，因为你们不知道}那贼\BibleQuote{何时来到。}\BibleRef{玛 24:42} 为此保禄也说，\par

第三节：\BibleQuote{若人们正说\InnerQuote{平安无事}时，毁灭就突然临到他们，如同产痛临到孕妇；他们决不能逃脱。}\BibleRef{得前 5:3}\par

这里，他提到了他在第二封书信中也说过的事情。因为那些人确实在患难中，而攻打他们的人却安逸奢侈，当他藉着提及复活来安慰他们当前的苦难时，另有些人却用来自他们祖先的论点侮辱他们，说：\Quote{这事何时发生？}先知们也说过：\BibleQuote{祸哉，那些说：愿天主赶快，加快祂的工作，好使我们看见；愿以色列圣者的计划来临，好使我们知道！}\BibleRef{依 5:19}；以及\BibleQuote{祸哉，那些渴望主的日子的人！}\BibleRef{亚 5:18}他是指这日子；因为他不是简单地说那些渴望它的人，而是说那些因不信而渴望它的人。并且，他说：\BibleQuote{上主的日子是黑暗，而不是光明}——请看保禄如何安慰他们，好像他说：不要让他们认为处于昌盛状态就是审判不来临的证明。因为审判确实会这样来临。\par

但值得一问：如果假基督来了，厄里亚来了，那么怎么在人们正说\InnerQuote{平安无事}时，毁灭就突然临到他们呢？因为这些事不会让那日子悄悄来临，它们是它来临的征兆。但他并非指这个时期是假基督的时期和整个日子，因为那将是基督来临的征兆，但祂自己不会有一个征兆，而是突然且出乎意料地来临。至于产痛，你确实说，不会出乎意料地临到孕妇；因为她知道九个月后生产会发生。然而，它非常不确定。因为有些人在第七个月生产，有些在第九个月。无论如何，日子和时辰是不确定的。因此，保禄这样说。并且这个比喻很精确。因为没有很多肯定的产痛征兆；许多人确实在大路上，或出门在外时生产，没有预见到。他不仅在这里暗示了不确定性，还暗示了痛苦的剧烈。因为当她在嬉戏、欢笑、毫无防备时，突然被无法言说的疼痛抓住，被产痛折磨——同样，当那日子临到那些灵魂时，也会如此。\par

\BibleQuote{他们决不能逃脱。} 正如他刚才所说的。\par

第四节：\BibleQuote{但是你们，弟兄们，不在黑暗中，以致那日子像贼一样抓住你们。}\BibleRef{得前 5:4}\par

此处他论及一种黑暗不洁的生活。因为正如腐败邪恶之人行事尽在夜间，逃避众人眼目，将自己隐匿于黑暗之中。试问，奸淫者岂不守望夜晚？盗贼岂不等待黑夜？掘墓者岂不在夜间从事其勾当？那么，这日子岂不也如贼般突然临到他们？它岂不是也毫无预兆地降临？为何他说\Quote{你们无需人写信给你们}？此处他不是指不确定性，而是指灾祸，即不会如同灾祸般降临于他们。因为那日子虽会不确定地临到他们，却不会使他们陷入困境。他说：\Quote{叫那日子如同贼一样，不能突然抓住你们}。因为对于警醒并处于光明中的人，纵有盗贼闯入，也不能加害他们。同样，对于生活端正之人也是如此。但那些沉睡之人，贼会剥去他们一切而离去。这沉睡之人即指那些信赖今生事物之人。\par

第5节。他说：\Quote{你们众人都是光明之子，白昼之子}。\BibleRef{得前 5:5}\par

如何可能成为\Quote{白昼之子}？正如经上说\Quote{毁灭之子}和\Quote{地狱之子}。因此基督也向法利塞人说：\Quote{祸哉，你们——因为你们走遍海洋陆地，要使一个人入教，及至他成了，你们却使他成为一个地狱之子}。\BibleRef{玛 23:15} 保禄又说：\Quote{为了这些事，天主的忿怒才降在悖逆之子身上}。\BibleRef{哥 3:6} 即指那些行地狱之事和悖逆之事的人。同样，天主的子女是指那些行事悦乐天主的人；同样，白昼之子和光明之子是指那些行光明之事的人。\Quote{我们不属于黑夜，也不属于黑暗}。\par

第6、7、8节。\Quote{所以我们不要像其余的人一样沉睡，而要警醒谨慎。因为人睡觉是在夜里睡，醉酒的人是在夜里醉。但我们既然是属于白昼的，就应当谨慎}。\BibleRef{得前 5:6}\par

此处他表明，处于白昼取决于我们自己。因为在此世的白昼与黑夜，不取决于我们。黑夜纵然我们不愿也会来临，睡眠在我们不愿时也会袭来。但在那黑夜与睡眠方面，却非如此；我们总能拥有白昼，总能警醒，这全在于我们的能力。因为闭上灵魂的眼目，带来邪恶的睡眠，并非出于本性，而是出于我们的选择。他说：\Quote{我们当警醒谨慎}。因为人虽醒着却可能沉睡，若不行善。所以他加上\Quote{谨慎}。因为即使在白昼，若有人警醒却不谨慎，必陷入无数危险；可见谨慎是警醒的强度。他说：\Quote{人睡觉是在夜里睡，醉酒的人是在夜里醉}。此处所说的醉酒不仅指醉酒，也指一切恶习所导致的醉态。因为财富与贪财是灵魂的醉态，肉欲也是如此；你能指出的每一项罪都是灵魂的醉态。那么为何称恶习为睡眠？首先，因为恶习之人对德性无所作为；其次，因为他所见一切如同幻象，不按真相看事物，而是充满梦境，常有不合理的行为；纵使看见善事，也无坚定稳固。现世生命正是如此：充满梦幻与幻想。财富是梦，荣耀以及诸如此类皆是梦。睡眠之人看不见真实存有之物，却把无有之物幻想为实有。恶习以及恶习中的生命正是如此：它看不见真实存有之物，即属灵的、天上的、持久之物，却看见转瞬即逝、飞逝而去、不久便离开我们之物。\par

但仅仅警醒谨慎还不够，还须武装。因为若人警醒谨慎，却没有武器，盗贼很快就能制服他。因此，当我们既应警醒、谨慎，又应武装，而我们却赤手空拳、赤身露体、沉睡时，谁能阻止敌人利剑直刺要害？为此，保禄也指明我们需要武器，再接着说：\par

第8节。\Quote{穿上信德和爱德的胸甲，并戴上救恩的望德作头盔}。\BibleRef{得前 5:8}\par

他说\Quote{信德和爱德}。此处他涉及生活和教义。他表明警醒谨慎为何，即拥有\Quote{信德和爱德的胸甲}。他说：不是平常的信德，而是如同胸甲一般，难以被刺穿，是胸前的坚固屏障——你也当如此，以信德和爱德环绕你的灵魂，如此，仇敌一切的火箭便不能射入。因为若灵魂的力量预先被爱德的铠甲占据，一切阴谋者的计谋都枉然而无效。因邪恶、仇恨、嫉妒、谄媚、虚伪等皆不能穿透这样的灵魂。他不仅简单地说\Quote{爱德}，而且命令他们以爱德作为坚固的胸甲穿上。\Quote{并戴上救恩的望德作头盔}。因为头盔保护我们体内的要害部位，环绕头部并四面覆盖，同样，这望德不容理性动摇，而是使之挺立如头，不容外物击打。且外物不击打它，它自身也不滑脱。因为凡以这等武装自卫的人，绝不可能跌倒。因为\Quote{现今存留的有信德、望德、爱德}。\BibleRef{格前 13:13} 然后他说：穿上，武装自己；而他自己提供铠甲，使信德、望德、爱德得以生发并坚固。\par

第9节。\Quote{因为天主预定我们不是为受怒，而是为藉我们的主耶稣基督获得救恩，祂为我们死了}。\BibleRef{得前 5:9}\par

如此，天主并未倾向于此，即祂要毁灭我们，而是为拯救我们。从何显明这是祂的旨意？祂为我们赐下了自己的圣子。祂如此渴望我们得救，以至于赐下圣子，且不仅是赐下，更赐祂至死。由此生出望德。人啊，当你走向天主时，不要对自己绝望，祂连自己的圣子都为你没有保留。不要因眼前的苦难而气馁。祂赐下独生子，为救你脱离地狱，今后为你的救恩岂会吝惜什么？因此你当盼望一切美善。因为我们若要去一位审判者面前受审，而祂曾为我们显示了如此大爱，甚至牺牲了自己的儿子，我们便不应惧怕。因此，让我们盼望良善与伟大的事物，因我们已领受了最要紧的；让我们相信，因我们已看见了榜样；让我们爱，因为不爱那如此对待我们的人，是极端的疯狂。\par

第10、11节。他说：\Quote{无论我们醒着或睡着，都要与祂一同生活。为此，你们应彼此劝慰，互相建立，正如你们一向所做的。}\par

再次，\Quote{无论我们醒着或睡着}；这里的\Quote{睡着}与那里的含义不同。因为在这里，\Quote{无论我们睡着}是指身体的死亡；也就是说，不要惧怕危险；即使我们死了，我们仍将活着。不要因身处危险而绝望。你有坚强的保障。祂若没有被对我们炽热的爱所燃烧，就不会赐下自己的圣子。因此，即使你死了，你也将活着；因为祂自己也死了。所以，无论我们死或生，我们都要与祂一同生活。这是无关紧要的：我活着或死了，都无关紧要；因为我们必要与祂一同生活。因此，让我们为那生命行一切事；注目于它，行一切善工。亲爱的，恶是黑暗，是死亡，是黑夜；我们看不见当看的，行不出当行的。正如死者是丑陋且恶臭的，同样，堕落者的灵魂也充满诸多不洁。他们的眼睛闭着，口被堵住，躺在恶的床上不动；或者比自然死亡的人更可怜。因为那些人确实对两者都死了，而这些人对德性无感，却对恶活着。若有人击打死人，他感觉不到，也不报复，如同枯木。同样，他的灵魂也真是枯干了，失去了生命；每日遭受无数创伤，却毫无感觉，对一切麻木不仁。\par

若将这样的人比作疯狂、醉酒或谵妄者，并不为过。这一切都属于恶，而且恶比这一切更甚。疯狂者尚且被看见的人多加谅解，因为他的病并非出于选择，而仅出于本性；但活在恶中的人，谁会宽恕他？恶从何而来？多数人为何变坏？告诉我，疾病从何而来其恶性？癫狂从何而来？昏睡症从何而来？岂非从疏忽而来？如果身体的疾病源于选择，那么自愿的疾病更是如此。酗酒从何而来？岂非从灵魂的放纵而来？癫狂岂非因高烧过甚？高烧岂非因身体内元素过多？这元素过多岂非来自我们的疏忽？因为当我们因缺乏或过度而将内在的任何东西带出节制的界限时，我们就点燃了那火。再者，若火已点燃，我们仍继续忽略它，我们就为自己制造了一场无法扑灭的大火。恶也是如此。当我们不在其起始时加以抑制、将其斩断，后来便无法到达其尽头，反而变得超出我们的能力。因此，我恳求你们，让我们尽一切努力，永不陷入困倦。你们不见吗？当哨兵稍一打盹，他们便从长久的守望中得不到任何益处，因为那一点打盹毁了一切，使预备偷窃的人有了十足的安全。正如我们不像贼那样看见他们，同样魔鬼尤其如此，它总是伺机而动，埋伏等候，咬牙切齿。因此，让我们不要打盹。不要说，这边没有危险，那边也没有；我们常从意想不到之处被抢掠。恶也是如此；我们从意想不到之处灭亡。让我们谨慎环视一切，不要醉酒，我们就不会睡着。不要奢侈，我们就不会打盹。不要为外在事物疯狂，我们就会保持清醒。让我们从各方面操练自己。如同走绳索的人，即使稍微不留神，那一点就会造成大害：人一旦失去平衡，即刻坠落而亡；同样，我们也不能稍微松懈。我们行走在一条狭窄的路上，两旁是悬崖，不容两脚同时站立。你们不见需要何等小心？不见走这样路的人不仅保护双脚，也保护双眼？因为若他选择向一侧观看，即使脚站稳，眼睛因深度而眩晕，也会使全身坠落。但他必须谨慎自己和他的脚步；因此经上说：\Quote{不可偏右，也不可偏左。}（\BibleRef{箴 4:27}）恶的深度巨大，悬崖高耸，下面黑暗沉沉。让我们当心这窄路，以敬畏战兢行走。走这样路的人，不会沉迷于欢笑，也不会被醉酒压垮，而是以清醒和斋戒行走这样路。走这样路的人，不会携带任何多余之物；因为他即使轻装，也满足于能逃脱。没有人会缠住自己的脚，而是让脚自由活动。\par

但是我们，用无数的挂虑把自己锁住，带着今世无数的重担，四处张望，随意漫游，如何指望能走在那条窄路上呢？祂不单是说\BibleQuote{那路是窄的}\BibleRef{玛 7:14}，而且带着惊叹说：\BibleQuote{那路是多么窄！}即极其狭窄。我们在令人惊叹的事上也这样做。祂又说：\BibleQuote{那导入生命的门是多么狭。}祂说得很好。因为我们必须为我们的思念、言语、行为及一切事物交账，那路实在是窄的。但我们自己使它更窄，舒展、放宽自己，把脚拖出来。因为窄路对每个人都是艰难的，尤其对那为脂肪所累的人；那使自己消瘦的人就不会察觉它的狭窄。所以，凡习惯受挤压的人，就不会因它的压力而沮丧。\par

所以，不要任何人期望他能轻松地看见天堂。因为这是不可能的。不要任何人希望以奢侈走窄路，因为这是不可能的。不要任何人走在宽路上却希望得到生命。因此，当你看见某人沉迷于浴池、丰盛的筵席，或其它事上有成群仆从时，不要因不享这些事而自以为不幸，反要为他哀哭，因为他正走向灭亡的路。这条路有何益处呢？它终归是困苦。那狭窄有何害处呢？它引领进入安息。告诉我，如果有人被邀请赴王宫，要走过痛苦而险峻的窄路，另一个人被拖向死刑，穿过市场中央，我们称谁为有福？我们可怜谁？不是那走宽路的人吗？所以现在，让我们认为有福的不是那些奢侈的人，而是那些不奢侈的人。这些人是奔向天堂，那些人是奔向地狱。\par

或许他们中许多人甚至会嘲笑我们所说的话。但我为此最哀叹悲恸，因为他们甚至不知道应当嘲笑什么，应当为何特别哀恸，反而混淆、扰乱、颠倒一切。因此我为他们哀恸。你说什么呢，人啊？当你要复活，要为你的行为交账，要受最后的判决，你竟不顾这些，只想满足肚腹，醉酒吗？你还嘲笑这些事？但我为你哀恸，因为我知道等待你的灾祸，即将临到你身上的惩罚。我尤其为此哀恸：你竟在笑！与我一同哀哭，与我一同为你的罪恶哀恸吧。告诉我，如果你的一个朋友死了，你不是离开那些嘲笑他结局的人，视他们为仇敌，却爱那些为你哭泣同情你的人吗？那么，若你妻子的尸体停放，你必离开那发笑的人；但当你自己的灵魂被置于死地时，你竟离开那哭泣的人，自己发笑吗？你看见魔鬼如何使我们成为自己的仇敌和对手吗？让我们清醒一次，睁开眼，警醒，抓住永生，摆脱这长眠。有审判，有惩罚，有复活，有对我们行为的究察！主在云中来临时，经上说：\BibleQuote{在祂面前，有火燃起，四周有猛烈风暴。}\BibleRef{咏 49:3}有火河在祂面前翻滚，不死的虫，不灭的火，外面黑暗，咬牙切齿。即使你对我提起这些事发怒一万次，我也不停止提及它们。因为如果先知们虽被石头砸死却不缄默，我们更应当忍受仇视，不为了讨你们喜欢而讲论，免得因欺骗你们而自身被分割。有惩罚，不死，不减轻，无人为我们辩护。经上说：\BibleQuote{被蛇咬的，谁会可怜那行邪术的呢？}\BibleRef{德 12:13}当我们不怜悯自己时，告诉我，谁会怜悯我们？如果你看见一个人用剑刺自己，你能救他的命吗？绝不能。更何况，当我们有能力行善却不行善时，谁会宽恕我们？没有人！让我们怜悯自己。当我们向天主祈祷说：\Quote{主，求你怜悯我}时，让我们对自己这样说，并怜悯自己。我们自己是天主怜悯我们的仲裁者。祂赐给了我们这恩宠。如果我们做配得怜悯、配得祂慈爱的事，天主就会怜悯我们。但如果我们不怜悯自己，谁会宽恕我们？怜悯你的近人，你必会从天主自己获得怜悯。每天有多少人来到你面前说：\Quote{可怜我吧}，你却不转向他们；有多少赤身露体者，有多少残废者，我们并不俯就他们，反拒绝了他们的哀求。那么，当你自己不做任何配得怜悯的事时，你怎能要求获得怜悯呢？让我们成为有同情心的人，成为有怜悯心的人，这样我们才能中悦天主，获得祂为爱祂的人所预许的美善，因我们的主耶稣基督的恩宠和慈爱，同祂……等。\par

\SourceRecord{Translated by John A. Broadus.；Nicene and Post-Nicene Fathers, First Series；Vol. 13.；Edited by Philip Schaff.；Buffalo, NY: Christian Literature Publishing Co.,；1889.；<http://www.newadvent.org/fathers/230409.htm>.}

% structure: p0001 source=h1 target=section reason=page-title

\section{得撒洛尼前书讲道十}

得撒洛尼前书 5:12-13\par

\BibleQuote{但是我们恳求你们，弟兄们，要认识那些在你们中间劳苦，并在主内管理你们、劝戒你们的人；并因他们的工作，以爱格外敬重他们。你们之间也要和平相处。}\par

统治者必定会有许多招致敌意的机会。如同医生被迫给病人带来许多麻烦，为他们准备既令人不快却有益的药膳和药物；也如同父亲常常惹恼自己的孩子：同样，教师也是如此，而且更甚。因为医生即使令病人憎恶，但与病人的亲属和朋友关系良好，甚至病人本人也往往如此。父亲也因天性力量和外在法律，轻易地对儿子行使统治权；若他违背儿子意愿进行管教和责斥，无人能阻止，儿子本人也不敢抬头看他。但就\OriginalTerm{司铎}{Priest}的情况而言，困难极大。首先，他应当管理愿意服从并感激他统治的人；但这不太可能立即实现。因为凡被定罪和斥责的人，不论他是谁，必定不再感恩，反而成为仇敌。同样，被劝告、被劝戒、被劝勉的人也会如此。因此，若我说：将财富施舍给穷人，我说的是令人反感且沉重的话。若我说：克制你的愤怒，熄灭你的怒火，抑制你过度的欲望，削减你一点点奢侈，这一切都是沉重且令人反感的。若我惩罚一个懒惰的人，或将他逐出\OriginalTerm{教会}{Church}，或将他排除在公共祈祷之外，他悲伤，并非因被剥夺这些事物，而是因当众蒙羞。因为这是恶上加恶：被禁止参与属灵之事，我们却不因丧失如此伟大的恩惠而悲伤，只因在他人面前蒙羞而难过。我们毫不战栗，毫不畏惧事情本身。\par

为此，保禄从一端到另一端大量谈论这些人的事。基督确实以如此严厉的必要性使他们服从，祂说：\BibleQuote{经师和法利塞人坐在梅瑟的座位上。凡他们对你们所说的，你们都要实行和遵守；但不要效法他们的行为。}\BibleRef{玛 23:2} 祂又治愈癞病人时说：\BibleQuote{你去，把你自身给司祭看，并献上梅瑟所规定的礼物，作为对证据给他们。}\BibleRef{玛 8:4} 然而祢却说：\BibleQuote{你们使他一倍两倍地成为地狱之子，比你们更甚}\BibleRef{玛 23:15}。因此，祂说：\BibleQuote{不要做他们所做的事。}所以祂为受管辖者堵住了所有借口。这位宗徒在致弟茂德书中也说：\BibleQuote{那善于管理的长老，当受加倍的尊敬。}\BibleRef{弟前 5:17} 在致希伯来人书中又说：\BibleQuote{你们要服从那些管辖你们的人，并要顺服他们。}\BibleRef{希 13:17} 而在这里又：\BibleQuote{但是我们恳求你们，弟兄们，要认识那些在你们中间劳苦，并在主内管理你们、劝戒你们的人。}因他先前说过：\BibleQuote{彼此建造，}以免他们以为他把他们提升到教师的等级，他补充说：然而看，我也准许你们彼此建立，因为教师不能说出一切。\BibleQuote{那些在你们中间劳苦的人，}他说，\BibleQuote{并在主内管理你们、劝戒你们的人。}他说，这岂不荒谬吗？若一个人在你们面前为你们辩护，你们做任何事，都承认自己欠他很多；但他在天主面前为你们辩护，你们却不承认这恩惠。他如何为我辩护？你们说，因为他为你们祈祷，因为他藉圣洗圣事将属灵的恩赐分施给你们，他探望、劝告和劝戒你们，你们午夜召唤他，他便前来；他不过是你们口中的常客，并承受你们的辱骂之言。他有何必要？他做了好事还是坏事？你们确实有妻子，生活奢靡，选择商业生活。但\OriginalTerm{司铎}{Priest}因他的职务而禁止自己这样做；他的生命无非是服务于\OriginalTerm{教会}{Church}。\BibleQuote{因他们的工作，以爱格外敬重他们；与他们和平相处。}你们看他多么清楚地意识到不快之感会产生？他不只是说\BibleQuote{爱，}而是\BibleQuote{格外，}如同儿女爱父亲。因为藉着他们，你们由那永远的诞生而重生：藉着他们，你们获得了王国：藉着他们的手，一切完成；藉着他们，\OriginalTerm{天堂}{heaven}的门为你们打开。愿无人制造分裂，愿无人争辩。凡爱基督的人，无论司铎是怎样的，都会爱他，因为藉着他获得了可畏的奥迹。告诉我，若你们想看一座用许多黄金闪烁、以宝石光辉照亮的宫殿，你们能找到拿钥匙的人，他应一声立即开门，让你进去，你们岂不将他置于众人之上？岂不爱他如同自己的眼珠？岂不亲吻他？这人已为你们打开了天堂，你们却不亲吻他，也不向他致敬。若你们有妻子，岂不爱那为你们获取她的人胜过一切？因此，若你们爱基督，若你们爱天国，要承认藉着谁获得了这一切。为此他说：\BibleQuote{因他们的工作，与他们和平相处。}\par

\BibleQuote{并且我们劝戒你们，弟兄们：要劝戒不守规矩的人，抚慰胆怯的人，扶持软弱的人，对众人要忍耐。}\BibleRef{得前 5:14}\par

在此，他规劝那些有管理权的人。他说：\Quote{要告诫那些不守秩序的人}，不是以专横或任性去责备他们，而是用劝告。\Quote{鼓励胆怯的人，扶持软弱的人，对众人都要忍耐。}因为若严厉责备，那人会自暴自弃，变得更加藐视。因此，必须用劝告使药变得甘甜。但谁是不守秩序的人呢？所有做违背天主之愿之事的人。因为教会的秩序比军队的秩序更加和谐；所以毁谤者是不守秩序的，酒徒是不守秩序的，贪婪的人，以及一切犯罪的人；因为他们不按规矩行事，而是偏离行列，因此他们被推倒。但还有另一种恶，不是这种，但本身也是恶习，就是心胸狭窄。因为这与懒惰一样有害。不能忍受侮辱的人是懦弱。不能忍受考验的人是懦弱。这就是撒在石头上的那一种人。还有另一种，就是软弱。他说：\Quote{扶持软弱的人}；现在软弱发生在信德方面。但注意，他不允许他们被轻视。在别的书信中他也说：\BibleQuote{信德软弱的，你们要接纳}\BibleRef{罗 14:1}。因为连在我们的身体上，我们也不允许软弱的肢体死亡。他说：\Quote{对众人都要忍耐。}甚至对那不守秩序的人？是的，当然。因为没有药能比这更有效，尤其是对于教师，对于管辖者最适合。它能使比所有人更凶恶更无耻的人羞愧和难堪。\par

第15节。\BibleQuote{你们要小心，谁也不要以恶报恶。}\BibleRef{得前 5:15}\par

如果我们不应该以恶报恶，何况以恶报善？更何况当人未曾做恶时，我们报以恶？你说，有一个人是坏人，得罪了我，给了我很多伤害。你想报复他吗？不要报复。让他不受惩罚。那么，这就是终点吗？绝不是；\par

\Quote{但要常追求良善，彼此相待，也待众人。}\par

这是更高的智慧，不仅不以恶报恶，反而以善报恶。因为这才是真正的报复，对他有害而对你自己有益，甚至对他也有大益，如果他愿意。为使你不以为这是仅对信徒说的，因此他说：\Quote{彼此相待，也待众人。}\par

第16节。\BibleQuote{要常常喜乐。}\BibleRef{得前 5:16}\par

这是针对带来苦难的试探而说的。你们许多陷入贫穷或困境的人，请听。因为从这些中，喜乐被产生。因为当我们拥有这样的灵魂，我们对任何人都不报复，而是善待众人，那么，请告诉我，悲伤的刺痛如何能进入我们？因为那在以恶报善中欢喜的人，怎能以后悲伤？然而，你说，这可能吗？如果愿意，是可能的。然后他也指出了方法。\par

第17、18节。\BibleQuote{不住地祈祷；凡事感谢：因为这是天主的旨意。}\par

常常感恩，是智慧灵魂的标志。你遭受了什么恶吗？但如果你愿意，它并非恶。感谢天主，恶就变为善。你也要像约伯那样说：\BibleQuote{愿上主的名永远受赞美。}\BibleRef{约 1:21}因为告诉我，你遭受了什么大事？疾病临到你？这并不奇怪。因为我们的身体是可朽的，容易受苦。财产匮乏临到你？但这些也是可得可失，停留在此世。是敌人的阴谋和诬告吗？但受伤害的不是我们，而是制造它们的人。因为他说：\BibleQuote{犯罪的人，他自己也必死亡。}\BibleRef{则 18:4}\par

因此，对于已死的人，你不应当报复，而应为他祈祷，使你能将他从死亡中救出。你没有看见蜜蜂螫人后就死吗？天主藉着这动物教导我们不要使邻人忧伤。因为我们自己先接受死亡。因为也许我们螫了他们，使他们痛苦一时，但我们自己将不再活着，正如那动物一样。然而圣经称赞它，说它是勤劳的，君王和庶民都利用它的工作来保养身体。\BibleRef{德 11:3}但这不能使它免于死亡，它必须灭亡。如果它的其他优点在它伤害人时不能救它，更何况我们呢？\par

因为确实，最凶猛的野兽，当无人伤害你时，才会主动伤害，或者甚至连野兽也不这样做。因为它们，如果你允许它们在荒野中觅食，且不因逼迫而迫使它们陷入绝境，它们永不会伤害你，也不会靠近你、咬你，而是走自己的路。\par

但你作为一个理性的人，被授予如此多的管理、尊荣和荣耀，却没有在对待同类上效法野兽，反而伤害你的兄弟，并吞吃他。你怎能为自己辩护呢？你没有听见\Person{保禄}说：\BibleQuote{为什么不情愿受亏负？为什么不情愿被讹诈？你们反倒自己亏负人、讹诈人，且是对你们的弟兄。}\BibleRef{格前 6:7}你看见吗？受害是作恶，而无辜受害是受惠？因为告诉我，如果有人辱骂他的上司，如果有人侮辱掌权者，他伤害谁？自己，还是他们？显然是伤害自己。那么，辱骂上司的人，不是侮辱上司，而是侮辱自己——而侮辱基督徒的人，不是藉着他侮辱基督吗？你说，决不是。你说什么？向国王（皇帝）的肖像扔石头的人，他向谁扔石头？不是向他自己吗？那么，向地上君王的肖像扔石头的人，是向自己扔石头，而侮辱天主肖像的人（因为人是天主的肖像）不也是伤害自己吗？\par

我们要喜爱财富到何时呢？因为我不会停止斥责它们：因为它们是万事的起因。我们何时才能不再被这种贪得无厌的欲望所满足呢？黄金有什么好处？我对这件事感到惊讶！这件事中有些魔力，使金银在我们当中如此受重视。因为我们对自己的灵魂毫不在意，却对这些无生命的像投以极大的关注。这疾病从何侵入世界？谁能摧毁它？什么理性能够斩断这只恶兽，并将其彻底消灭？这种欲望深深植根于人的思想中，甚至包括那些看似虔诚的人。让我们因福音的命令而感到羞愧。圣经中只有言语，却没有付诸实行。\par

众人常用的托词是什么呢？有人说：我有孩子，我害怕自己陷入极端饥饿和匮乏的境地，害怕自己需要依赖他人。我羞于乞讨。因此你就使别人乞讨吗？你说：我不能忍受饥饿。因此你就使别人挨饿吗？你知道乞讨是多么可怕的事吗？挨饿是多么可怕？也怜悯你的弟兄吧！请告诉我，你以饥饿为耻，却不以抢劫为耻吗？你害怕饿死，却不害怕毁灭别人吗？饥饿既非耻辱，也非罪行；但使别人陷入这种境地，不仅带来耻辱，还有极重的惩罚。\par

这些都是借口、空话、琐事。因为那些没有孩子、也不会有孩子的人，他们同样劳苦、折磨自己、忙于积累财富，就好像有无数孩子可以继承一样，这证明了你们并非为了孩子才如此行。使人贪婪的不是对孩子的关心，而是灵魂的疾病。因此，许多没有孩子的人也迷恋财富，而另一些有众多子女的人却藐视财富。他们将在那一日控告你们。因为如果孩子的需要迫使人积累财富，那么他们也必然有同样的渴望、同样的贪欲。如果他们不这样，那么我们的疯狂就不是来自孩子的数量，而是来自对金钱的贪爱。你说：谁有孩子却藐视财富？这样的人很多，在各处都有。如果你允许，我也要提到古人中的例子。\par

雅各伯不是有十二个孩子吗？他不是过着雇佣工人的生活吗？他不是被亲戚亏待了吗？他不是常常使那亲戚失望吗？孩子的数量曾迫使他采用任何不诚实的计谋吗？亚巴郎的情况如何？依撒格不也有许多孩子吗？那又怎样？他不是将一切所有都用于款待外方人吗？你们看见，他不仅没有行不义，甚至舍弃了自己的财产，不仅行善，而且甘愿被侄子亏待。因为为了天主而忍受被抢夺，比行善伟大得多。为什么呢？因为行善是灵魂自由选择的果实，因此容易实践；但被抢夺却是伤害和暴力。一个人会自愿抛弃一万塔冷通，而不觉得自己受了任何损害，却很难在被抢走三个便士时不发怒。所以，这更是灵魂的修养。我们看到，这在亚巴郎身上实现了。经上说：\Quote{罗特举目看见整个平原，有水灌溉，如同天主的乐园，便选择了它。}\BibleRef{创 13:10}亚巴郎没有说什么反对的话。你们看见，他不仅没有亏待他，反而被他亏待了。人啊，你为什么责备自己的孩子？天主赐给我们孩子，并非为此，好使我们抢夺别人的财物。当心，不要因这样说而激怒天主。因为如果你说孩子是你贪婪和吝啬的原因，我恐怕你会失去他们，因为他们伤害了你、诱骗了你。天主赐给你孩子，是为了让他们在年老时支持你，从你那里学习德行。\par

天主愿意人类如此结合在一起，提供了两个最重要的目标：一方面是立父亲为教师，另一方面是植入伟大的爱情。因为如果人只是被生出来，没有人会对别人有任何关系。因为现在，即使有父亲、孩子、孙子的关系，许多人也不关心许多人，那么那时就更不关心了。为此，天主赐给你孩子。所以，不要控告孩子。\par

但如果有孩子的人没有借口，那么没有孩子却为积累财富而耗尽心力的人，他们又能为自己说什么呢？他们也有自己的说法，但毫无借口。那是什么呢？他们说：可以作为孩子的替代，让我们的财富成为纪念。这真是可笑。有人说：代替孩子，我的房子成为我不朽荣耀的纪念。人啊，那将不是你荣耀的纪念，而是你贪婪的纪念。难道你没有看见，如今许多人经过宏伟的房屋时互相说：某人为了建造这房子，犯了多少欺诈、抢夺；如今他已化为尘土，他的房子却成了别人的产业！所以你留下的不是荣耀的纪念，而是贪婪的纪念。你的身体确实被埋在地下，但你却不允许你贪婪的纪念被隐藏，尽管时间可能将其掩埋；你反而通过你的房子将它翻掘出来。因为这房子只要存在，带着你的名字，被称为某人的房子，所有人的口必然要开口攻击你。你看见了吗？一无所有总比遭受这样的控告要好。\par

这些事确实在此。但在那里我们将做什么呢？请告诉我，我们在此拥有如此多的财富，若我们不曾将我们的财产分施给任何人，或者至少很少施予，我们怎能摆脱不义之财呢？因为凡愿意摆脱贪财之利的，不是从大量中施予少许，而是偿还超过他所掠夺的许多倍，并且停止掠夺。请听\Person{匝凯}说什么：\BibleQuote{凡我讹诈来的，我偿还四倍。}\BibleRef{路 19:8} 但你讹诈了一万塔冷通，若只施予几块德拉克马，你以为就偿还了全部，并自以为是施予了很多。甚至这施予还是勉强的。为什么？因为你应当既偿还这些，又外加你自己私产中的其他。因为盗贼若只归还他所偷的，并不算免罪，反而往往加上他的性命；常常他通过归还许多倍来和解：贪财者也当如此。因为贪财者也是盗贼和强盗，远比那更坏，因为他更暴虐。那盗贼因隐蔽、夜间袭击，减少了尝试的胆量，仿佛羞愧、害怕犯罪。但另一个毫无羞耻，在光天化日之下，在市场中央偷窃所有人的财产，同时是盗贼和暴君。他不破墙，不灭灯，不开箱，不撕封。但怎样？他做更为无礼的事，在受害者眼前将东西从门口搬出，他自信地打开一切，强迫受害者自己暴露所有财产。这就是他暴力的过度。此人比那些更邪恶，因他更无耻、更残暴。因为受欺诈的人固然悲伤，但有一个不小的安慰：伤害他的人曾怕他。但那人既受伤害又被轻视，将无法忍受这暴力，因为嘲笑更大。请告诉我，若一人与妇人秘密通奸，另一人在她丈夫面前通奸，谁更使丈夫痛苦、更易伤害他？因为那人连同他所行的恶事，也藐视了丈夫。但前者，若没有做别的，至少显示他害怕他所伤害的人。同样在钱财的事上：秘密取走的人，在这方面尊敬了他，因为他秘密地做；但公开抢夺的人，连同损失还加上了羞耻。\par

因此，我们无论贫富，都当停止夺取他人的财产。因为我现在的讲论不仅是对富人，也是对穷人。因为他们也抢夺比他们更穷的人。那些景况较好、较有实力的工匠，排挤更穷困的工匠；商人排挤商人；所有在市场谋生的人无不如此。所以我愿从各方除去不义。因为伤害不在于被掠夺和偷窃的物的数量，而在于偷窃者的意图。并且那些不轻视小利的人，更是盗贼和欺诈者，我知道并记得我以前告诉过你们，如果你们也记得的话。但我们不要过于苛求。让他们与富人同样坏吧。我们当教导自己的心：不要贪图更大的事物，不要追求超过我们所有的。在属天的事物上，愿我们永不对更多有满足，而让人人永远贪图更多。但在世上，让人人只求需要和足够，不再寻求更多，这样他才能获得真正的美善，赖我们的主耶稣基督的恩宠和慈爱。愿光荣、能力、尊崇，藉着祂，归于圣父及圣神，自今至永远，世世无穷。阿们。\par

\SourceRecord{Translated by John A. Broadus.；Nicene and Post-Nicene Fathers, First Series；Vol. 13.；Edited by Philip Schaff.；Buffalo, NY: Christian Literature Publishing Co.,；1889.；<http://www.newadvent.org/fathers/230410.htm>.}

% structure: p0001 source=h1 target=section reason=page-title

\section{得撒洛尼前书第十一篇讲道}

\BibleRef{得前 5:19-22}\par

\BibleQuote{不要熄灭圣神。不要轻视先知预言。但要考验一切，好的要保持，各种邪恶的形态要远避。}\par

浓雾、黑暗和云彩笼罩着全地。为此，宗徒说：\BibleQuote{因为我们从前是黑暗。}\BibleRef{弗 5:8} 又说：\BibleQuote{弟兄们，你们不在黑暗里，以致那日子像贼一样袭击你们。} 既然可以说是一个没有月光的黑夜，我们行在夜里，天主就给了我们一盏明灯，点燃了我们灵魂中圣神的恩宠。但有些领受这光的人使它更加明亮耀眼，如\Person{保禄}、\Person{伯多禄}和所有诸位圣人；而另一些人甚至把它熄灭了，如五个童女，那些\BibleQuote{在信仰上遭遇了船难}的人，格林多的淫乱者，以及被败坏了的迦拉达人。\par

为此，\Person{保禄}说：\BibleQuote{不要熄灭圣神}，即恩宠的恩赐，因为他惯常这样称圣神的恩赐。但不洁的生活会熄灭它。如同一个人把水和灰尘洒在我们的灯上，就会熄灭火，或者不这样做，只是取走油——同样，恩宠的恩赐也是如此。如果你把世俗的事物和纷扰的忧虑倒在它上面，你就熄灭了圣神。即使你没有做这些，但从别处来的诱惑猛烈地攻击它，如同一阵风，而光并不强，油也不多，或者你没有堵住开口，或者没有关上门，一切就都毁了。什么是开口？如同在灯中，在我们内也是如此：就是眼睛和耳朵。不要让一股强烈的邪风袭击它们，否则会熄灭灯，而要用对天主的敬畏堵住它们。嘴是门。关上它，锁住它，这样既能发光，又能抵挡外来的攻击。例如，有人侮辱、谩骂你？你要闭口；因为如果你开口，你就加强了风力。你没有看到在房屋中，当两扇门正对着，有强风时，如果你关上一扇，没有对流，风就没有力量，其大部分力量就减弱了？现在也是如此，有两扇门：你的口，和那侮辱、冒犯你的人的口；如果你闭上你的口，不让另一边的气流进来，你就熄灭了整个风暴；但如果你开口，就无法抑制。因此，我们不要熄灭它。\par

并且，即使没有诱惑袭击，火焰也常常容易熄灭。当油不够、我们不行哀矜时，圣神就熄灭了。因为祂是从天主而来的哀矜施予你。然后祂看到这果实不在你内，祂不会停留在一个没有怜悯的灵魂里。但圣神一熄灭，你们中曾行走在无月之夜道路上的人，就知道接下来是什么。如果从陆地到陆地的夜间之路都难走，那么从地到天的路又怎能安全呢？难道你不知道中间空域有多少魔鬼，多少野兽，多少邪恶的灵体吗？如果我们确实有那光，它们就不能伤害我们；但如果我们熄灭了它，它们就会立刻俘获我们，立刻夺走我们的一切。因为甚至强盗也是先熄灯，然后抢劫我们。他们确实在黑暗中能看见，因为他们行黑暗的事；但我们不习惯那光。所以不要熄灭它。一切恶行都会熄灭那光，无论是辱骂、傲慢，还是你能提到的任何事。如同在火的情形中，一切与其本性相异的东西都会毁灭它，而与之相合的东西则点燃它：干燥的、温暖的、火性的东西，点燃圣神之火。因此，不要用任何冷淡或潮湿的东西覆盖它，因为这些东西会毁灭它。\par

还有另一种解释。他们中有许多确实说真预言，但有些人说假预言。他在致格林多人书里也说到这一点，因此祂赐给了\BibleQuote{辨别神恩}\BibleRef{格前 12:10}。因为魔鬼出于其邪恶的诡计，想借此恩赐颠覆教会的一切。既然恶魔和圣神都预言未来，一个说假话，一个说真话，并且无法从任何地方得到对两者的验证，而是各自不受质问地发言，如同\Person{耶肋米亚}和\Person{厄则克耳}所做的那样，但到时就会被定罪，祂也赐给了\BibleQuote{辨别神恩}。既然在得撒洛尼人中也有许多人预言，他暗指他们说：\BibleQuote{无论是言语，或是书信，好像出于我们，说主的日子已经来到}\BibleRef{得后 2:2}，他在这里这样讲。即，不要因为你们中间有假先知，就因此禁止这些，并远离它们；\BibleQuote{不要熄灭}它们，即\BibleQuote{不要轻视先知预言}。\par

你看到这正是他所说的\BibleQuote{考验一切}的意思吗？因为他曾说\BibleQuote{不要轻视先知预言}，恐怕他们以为他开放讲台给所有人，他就说：\BibleQuote{考验一切}，即那些确实是预言的事；\BibleQuote{好的要保持}。\BibleQuote{各种邪恶的形态要远避}；不是这个或那个，而是所有；使你们能通过考验辨别真假，远避后者，保持前者。因为这样，对一方的憎恶才会强烈，对另一方的爱才会产生，当我们做事不草率、不无审查，而是仔细查考时。\par

第23节。\BibleQuote{愿平安的天主亲自完全圣化你们，并愿你们整个的灵、魂和身体，在我们的主耶稣基督来临时，无可指摘地得以保全。}\BibleRef{得前 5:23}\par

请注意老师的情感。在劝诫之后，他加上了祈祷；不仅如此，他甚至将它写在信中。因为我们既需要劝告也需要祈祷。为此，我们也是先给你们劝告，然后为你们祈祷。这事\OriginalTerm{领洗者}{Initiated}知道。但保禄这样做，确有充分的理由，因为他对于天主有很大的信心，而我们却羞愧难当，没有发言的胆量。但由于我们被任命做这事，我们就做，即使我们不配站在祂面前，不配占有最卑微门徒的位置。但因为恩宠甚至借着不配的人工作，不是为我们自己，而是为那些将要受益的人，我们尽我们的一份。\par

\Quote{愿你们完全成圣，} 他说，并愿\Quote{你们的灵、魂和身体完全得蒙保守，在我们主耶稣基督来临时无可指摘。} 他在这里所说的灵是指什么？是指恩宠的恩赐。因为如果我们离开时点着灯，我们就能进入婚筵。但如果灯灭了，就不能。因此他说\Quote{你们的灵。} 因为如果灵保持纯洁，其他部分也保持纯洁。他说\Quote{灵魂和身体。} 因为那时二者都不容纳任何邪恶。\par

24. \BibleQuote{召唤你们的那位是信实的，他必会成就此事。}\BibleRef{得前 5:24}\par

请注意他的谦卑。因为他祈祷了，他说：不要以为这是因为我的祈祷而成就，乃是因为召唤你们的旨意。假如他召唤你们是为了救恩，而他是信实的，他必会拯救你们，因为他愿意这样做。\par

25. \BibleQuote{弟兄们，也请为我们祈祷。}\BibleRef{得前 5:25}\par

奇怪！这是何等的谦卑！但他说这话是为了谦卑，而我们，并非出于谦卑，而是为了极大的益处，并希望从你们那里获得一些大的益处，却说：\Quote{也请为我们祈祷。} 因为即使你们没有从我们这里得到任何大的或奇妙的益处，也要因着这份尊敬和这称号本身而这样做。有人有儿女，即使儿女从未从他那里得到益处，但因他是他们的父亲，他或许会将此放在他们面前，说：\Quote{有一天我没有被你们称为父亲。} 为此我们也说：\Quote{也请为我们祈祷。} 我这样说不仅仅是口头说说，而是真心渴望你们的祈祷。因为如果我必须为对你们众人的这份管理负责，并要交账，那么我更应当得到你们祈祷的益处。因你们的缘故，我的责任更大，因此从你们那里得来的帮助也应更大。\par

26. \BibleQuote{你们要用圣吻问候所有弟兄。}\BibleRef{得前 5:26}\par

哦！何等的热忱！哦！这是何等的炽爱！因为他不能当面向他们行吻礼，就通过别人来问候他们，如同我们所说：替我亲他。因此你们也要保持爱火。因为爱不因距离而止，它甚至能穿越长路而延伸，无处不在。\par

27. \BibleQuote{我指着主嘱咐你们，要把这封书信读给所有圣洁的弟兄听。}\BibleRef{得前 5:27}\par

这个命令更是出于爱，而并非出于教导；意思是，借此我也可以与他们交谈。\par

28. \BibleQuote{愿我们主耶稣基督的恩宠与你们同在。阿们。}\BibleRef{得前 5:28}\par

他不仅命令，而且用誓词嘱咐他们，这是出于热切的心，好让他们即使轻视他，也会因着这个誓词而遵行所吩咐的。因为那时人们非常惧怕这种恳求，但现在它也被践踏了。常常有奴隶被鞭打，恳求天主和祂的基督，说：\Quote{愿你也像一个基督徒那样死去，} 却无人理会，无人重视；但若他指着自己的儿子恳求，那么他立即会放下愤怒，尽管不情愿，咬牙切齿。又有另一个人被拖拽着穿过市场，在犹太人和外邦人面前，用最可怕的誓词恳求那拖走他的人，却无人理会。当一个信徒恳求另一个信徒和基督徒，却无人理会，我们甚至轻视他时，外邦人会怎么说呢？\par

请允许我告诉你们一个故事，是我亲耳听到的？因为我并不是凭自己杜撰，而是从一位可靠的人那里听来的。有一个女仆，嫁给了一个邪恶、卑劣的逃跑奴隶。她的丈夫犯了很多过错，女主人要把他卖掉；（因为过错太大，无法宽恕；那女人是个寡妇，无法惩罚家里的祸害，所以决定卖掉他；然后考虑到将丈夫与妻子分开是不圣洁的事，女主人尽管那个女仆有用，为了避免拆散他们，也决定将她一起卖掉；）女仆看到自己陷入这样的困境，就来到一位可敬的人面前，这人与女主人相熟，也把这事告诉我了。她抱住那人的膝盖，千般哀求，求她替她向女主人说情；说了许多话之后，最后她加上这句，以为特别能打动她，便加上一个最可怕的恳求，那恳求是：\Quote{愿你在审判之日看见基督，如同你不忽视我的请求一样。} 说完她就离开了。那位被恳求的人，因被一些世俗的事务打扰，如同家中常发生的那样，忘记了这事。然后到了下午晚些时候，那可怕的恳求突然进入她脑海，她深感悔恨，便去急切地请求，并获得了所求。就在那夜，她突然看见天开了，基督本人显现。但她看见了他，如同一个妇人可能看见的那样。因为她尊重了那恳求，因为她害怕了，便堪当了这异象。\par

我说这些话，为叫我们不要轻视恳求，尤其是当别人为善事恳求我们时，例如为施舍和慈善工作。但现在那些失去双脚的穷人，坐着看见你匆匆经过，当他们不能用脚跟随你时，他们期望像用钩子一样，借着恳求的恐惧留住你，伸出双手，恳求你只给他们一两文钱。但你却匆匆经过，尽管被你的主恳求。如果他以你外出的丈夫的眼睛，或你的儿子、你的女儿的眼睛恳求你，你立刻让步，你的心被感动，你变得温暖；但如果他以你的主恳求你，你却匆匆经过。我知道许多妇女，确实听到了基督的名，就匆匆经过；但若有人来夸赞她们的美貌，她们就被融化、变软，并伸出手来。\par

是的，他们就是这样使受苦且悲惨的乞丐沦落至此，甚至以制造笑料为业！因为当他们不能用激烈和苦涩的话触动他们的灵魂时，就诉诸这种方法，使他们大为快乐。我们的极大邪恶迫使那遭受灾难或为饥饿所困的人，去赞美那些怜悯他之人的美貌。我希望仅止于此。但还有另一种比这更坏的形式。它迫使穷人成为杂耍者、小丑和污秽的玩笑者。因为他将杯子、碗和罐子固定在手指上，像铙钹一样敲击它们，又有一支笛子，用它吹出那些低贱而淫荡的旋律，并用最大声音歌唱；然后许多人围成一圈，有人给他一块面包，有人给他一文钱，还有人给他别的东西，他们长久地留住他，男女都感到快乐；还有什么比这更可悲的呢？这些事岂不令人大为叹息吗？它们确实是琐碎的，也被视为琐碎，但它们在我们的品性中生出大罪。因为每当淫秽而甜美的旋律发出，它软化心灵，败坏灵魂本身。而那个呼求天主、为我们祈求千万祝福的穷人，你不肯赐他一言；但那个以这些事代替祝福而插入戏谑俏皮话的人，却被赞赏。\par

现在我心里想到要对你们说的话，我要说出来。是什么呢？当你陷入贫穷和疾病时，若没有别的途径，至少从那些在窄街中行乞的乞丐身上，学习感谢主。因为他们终其一生行乞，并不亵渎，不生气，也不急躁，反而将他们乞讨的全部叙述放在感恩中，尊崇天主，称祂为仁慈的。那个因饥饿而濒死的人，称祂为仁慈的；但你们生活在富足中，若不能获得所有人的财产，就称祂为残忍。他有多么好！他将如何定我们的罪！天主将穷人派遣到世上，作为我们在患难中的共同老师，以及患难中的安慰。你是否遭受过任何违背你意愿的事？然而没有一件像那个穷人所受的那样。你失去了一只眼睛，但他失去两只。你长期受疾病折磨，但他患的是不治之症。你失去了儿女，但他甚至失去了自身身体的健康。你遭受了巨大损失，但你还没有沦落到向他人乞求。感谢天主吧。你看见他们在贫穷的熔炉中，确实向所有人乞讨，但从少数人那里得到。当你厌倦祈祷，没有得到时，想想多少次你听到一个穷人呼求你，而你没有听他，他却没有生气也没有侮辱你。然而你那样做是出于残忍；但天主出于仁慈甚至拒绝垂听。因此，如果你自己出于残忍不听你的同仆，却期望不被责备，那么你责备那出于仁慈不听他仆人的主吗？你看见多么大的不平等，多么大的不义吗？\par

让我们不断思想这些事，那些比我们低下的人，那些遭受更大灾难的人，这样我们就能感谢天主。生命中有许多这样的实例。那清醒并愿意留意的人，从祈祷之所中得到不小的教训。因为正是为此，穷人坐在教会和殉道者小堂的门廊前，好使我们从这景象中得到极大的好处。因为想一想，当我们进入地上的宫殿时，看不到任何这类的事；但尊贵而著名、富有而聪明的人，到处匆匆来往。但进入真正的宫殿——我是指教会和殉道者的祈祷所——的，是附魔者、残废者、穷人、老人、盲人以及肢体扭曲的人。为什么呢？为叫你从这景象中得到教导：首先，如果你进入时带着任何从外面来的骄傲，看到这些后，放下你的傲慢，心中痛悔，然后你进去，听所讲的话；因为傲慢之心祈祷的人，不可能被垂听。当你看见一位老人时，你不要因自己的年轻而得意，因为这些老人曾经年轻。当你夸耀你的军事或王权时，你要想到，那些在君王宫廷中显赫的人也是从这些人中出来的。当你凭恃你的身体健康时，留心这些事，你就可以降低你的高傲精神。因为那经常进入这里的健康人，不会因他的身体健康而高傲；而病人会得到不小的安慰。\par

但他们坐在这里，不仅为这个缘故，也是为使你们有同情心，使你们倾向于怜悯；使你们钦佩天主的仁慈；因为如果天主不以他们为耻，反而将他们安置在自己的前院，你们更不该羞愧；使你们不因地上的宫殿而自高。当被穷人请求时，不要羞愧；如果他走近，如果他抓住你们的膝盖，不要推开他。因为这些是王庭中某些可佳的狗。我说他们是狗，并非侮辱他们——万万不可——反而是极称赞他们。他们守护着君王的宫殿。所以要喂养他们。因为这荣誉归于君王。在那里，一切都是骄傲——我说的是地上的宫殿——在这里，一切都是谦卑。你们尤其从这些前院本身学到，人什么也不是。从这些坐在他们面前的人，你们学到天主不以财富为乐。因为他们的坐与聚集几乎是一种劝诫，发出清晰的声音，关乎所有人的本性，说人的事什么也不是，它们是影子和烟雾。如果财富是好事，天主就不会将穷人安置在自己的前院前。如果祂也接纳富人，不必惊奇，因为祂接纳他们，并非为使他们继续富有，而是为使他们从重负中得解脱。因为请听基督对他们说的话：\BibleQuote{你们不能事奉天主而又事奉财帛} \BibleRef{玛 6:24} ；又说：\BibleQuote{富人进天国是难的} ；又说：\BibleQuote{骆驼穿过针眼，比富人进天国还容易。} \BibleRef{玛 19:23} 为此，祂接纳富人，为使他们听见这些话，为使他们渴求永久的财富，为使他们贪恋天上的事。你为什么惊奇祂不耻于将这样的人安置在祂的前院？因为祂不耻于召唤他们到祂的属灵筵席，使他们成为那宴席的分享者。而且，瘸腿的、跛足的、身穿破衣和污秽、患有鼻炎的老年人，前来与年轻貌美者，甚至与身穿紫袍、头戴王冠者同席——被认为配得上这属灵筵席，两者享受同样的恩惠，毫无区别。\par

那么，基督不耻于召唤他们与君王（皇帝）同席——因为二者同被召唤——而你们或许甚至耻于被看见周济穷人，或与他们交谈？可耻的傲慢和骄傲！要当心我们不会遭受从前那个富人的同样命运。他耻于看拉匝禄一眼，不允许他分享他的屋顶或遮蔽，反而让他被抛在外，在他门前遭弃，甚至不配从他那里得到一句话。但请看，当他陷入困境，需要他的帮助时，却得不到。因为如果我们以基督不耻于的人为耻，我们就是以基督为耻，以祂的朋友为耻。让你们的筵席坐满瘸腿的和跛足的。通过他们，基督来临，不是通过富人。也许你们听了这个发笑；因此，为使你们不以为这是我的话，请听基督亲自说话，使你们不笑，反而战栗：\BibleQuote{几时你设午宴或晚宴，不要请你的朋友、兄弟、亲戚及富有的邻人，免得他们也回请你，报答了你。但你设宴时，要请贫穷的、残废的、瘸腿的、瞎眼的，你就有福了，因为他们无力报答你；你将在义人复活时得到报答。} \BibleRef{路 14:12} 而且，即使在此世，你们的荣耀也更大，如果你们喜爱它。因为从前一类客人中产生嫉妒、恶意、诽谤、辱骂以及许多恐惧，唯恐发生不体面的事。你们像仆人站在主人面前，如果受邀请的是你们的上级，你们便会害怕他们的批评和他们的嘴唇。但在这些穷人身上，没有这类事；无论你们带给他们什么，他们都欣然接受；而且有丰盛的赞美，更光耀的荣耀，更高的钦佩。所有听见的人，并不像赞美前者那样赞美后者。但如果你们不信，你们这些富人，请做试验，你们这些邀请将军和长官的人。邀请穷人，用他们充满你们的餐桌，看是否不被所有人赞美，是否不被所有人爱戴，是否不将你们视为父亲。因为那些宴席没有益处，但为了这些宴席，有天堂和天堂的美好事物——愿我们所有人都能分享这些，藉着我们的主耶稣基督的恩宠和慈爱，愿光荣、权能、尊崇归于祂及圣父和圣神，从现在直到永远，万世无疆。阿们。\par

\SourceRecord{Translated by John A. Broadus.；Nicene and Post-Nicene Fathers, First Series；Vol. 13.；Edited by Philip Schaff.；Buffalo, NY: Christian Literature Publishing Co.,；1889.；<http://www.newadvent.org/fathers/230411.htm>.}
